% !TeX program = xelatex
% Standalone English volume. Use XeLaTeX until the references stabilize.
\documentclass[a4paper,11pt]{article}
\usepackage[T1]{fontenc}
\usepackage{lmodern}
\usepackage[english]{babel}
\usepackage[left=26mm,right=26mm,top=25mm,bottom=25mm,headheight=15pt,headsep=7mm,footskip=11mm]{geometry}
\usepackage{amsmath,amssymb,mathtools,amsthm}
\newtheorem{proposition}{Proposition}[section]
\newtheorem{lemma}[proposition]{Lemma}
\newtheorem{theorem}[proposition]{Theorem}
\newtheorem{corollary}[proposition]{Corollary}
\newtheorem{maintheorem}[proposition]{Main Theorem}
\theoremstyle{definition}
\newtheorem{definition}[proposition]{Definition}
\theoremstyle{remark}
\newtheorem{remark}[proposition]{Remark}
\makeatletter
\renewcommand*\l@section{\@dottedtocline{1}{0em}{2.4em}}
\renewcommand*\l@subsection{\@dottedtocline{2}{1.45em}{3.45em}}
\renewcommand*\l@subsubsection{\@dottedtocline{3}{3.8em}{4.55em}}
\makeatother
\usepackage{xcolor,fancyhdr,titlesec,needspace,xurl}
\usepackage{longtable,multicol}
\usepackage{hyperref}
\definecolor{ink}{HTML}{1A3542}
\definecolor{muted}{HTML}{4F606A}
\hypersetup{unicode=true,colorlinks=true,linkcolor=ink,filecolor=ink,urlcolor=ink,citecolor=ink,bookmarksnumbered=true,
 pdfdisplaydoctitle=true,pdflang={en},
 pdftitle={The generalized Riemann hypothesis for Dirichlet L-functions. Edition 930, Volume X},
 pdfsubject={A natural-language reconstruction of the proposed proof},
 pdfauthor={A mathematical reconstruction}}
\titleformat{\section}{\Large\bfseries\color{ink}}{\thesection}{.65em}{}
\titlespacing*{\section}{0pt}{.8em}{.65em}
\setlength{\parindent}{1.2em}
\linespread{1.045}
\setlength{\parskip}{.22em plus .05em minus .03em}
\setlength{\emergencystretch}{2em}
\setcounter{topnumber}{2}
\setcounter{tocdepth}{3}\setcounter{secnumdepth}{3}
\widowpenalty=4500
\clubpenalty=4500
\pagestyle{fancy}\fancyhf{}
\fancyhead[L]{\small\color{muted}Dirichlet GRH \textperiodcentered\ Volume I}
\fancyhead[R]{\small\color{muted}Dirichlet zeros and power intervals}
\fancyfoot[C]{\small\thepage}
\renewcommand{\headrulewidth}{.3pt}
\newcommand{\lead}[1]{\par\smallskip\noindent{\bfseries\color{ink}#1.}\enspace}
\newcommand{\sour}[2]{\unskip\penalty150\hspace{.35em}{\small\hyperlink{ref#1}{[#2]}}}
\titleformat{\subsection}{\large\bfseries}{\thesubsection}{.6em}{}
\titlespacing*{\subsection}{0pt}{1.15em plus .15em}{.4em}
\setlength{\abovedisplayskip}{9pt plus 2pt minus 3pt}
\setlength{\belowdisplayskip}{9pt plus 2pt minus 3pt}
\newcommand{\RR}{\mathbb R}\newcommand{\CC}{\mathbb C}\newcommand{\QQ}{\mathbb Q}
\newcommand{\ZZ}{\mathbb Z}\newcommand{\ord}{\operatorname{ord}}
\newcommand{\Frob}{\operatorname{Frob}}\newcommand{\Gal}{\operatorname{Gal}}
\newcommand{\len}{\operatorname{length}}\newcommand{\Res}{\operatorname*{Res}}
\newcommand{\wmean}{\operatorname{av}}\newcommand{\Ker}{\operatorname{ker}}
\newcommand{\Lh}{\mathcal L}\newcommand{\Gg}{\mathcal G}
\newcommand{\real}{\operatorname{Re}}\newcommand{\imag}{\operatorname{Im}}
\newcommand{\abs}[1]{\left|#1\right|}\newcommand{\eps}{\varepsilon}
\newcommand{\anchoritem}[2]{\par\Needspace{2.8\baselineskip}\hypertarget{ref#1}{}\noindent{\bfseries\color{ink}#2.}\enspace}
\newcommand{\srcloc}[2]{{\small\nolinkurl{#1.lean}}, line~#2}
\numberwithin{equation}{section}
\raggedbottom

\newcommand{\Sel}{\mathcal S}
\newcommand{\Cl}{\operatorname{Cl}}
\newcommand{\Aut}{\operatorname{Aut}}
\newcommand{\GL}{\operatorname{GL}}
\newcommand{\OO}{\mathcal O}
\hypersetup{bookmarksdepth=3}


\addto\captionsenglish{\renewcommand{\partname}{Volume}}


\makeatletter
\expandafter\def\csname r@cor:actualmarks714\endcsname{{12.10}{38}{}{proposition.12.10}{Dirichlet-GRH-v930-Volume-II.pdf}}
\expandafter\def\csname r@cor:actualperiodmultiples768\endcsname{{16.28}{59}{}{proposition.16.28}{Dirichlet-GRH-v930-Volume-III.pdf}}
\expandafter\def\csname r@cor:allscalartests770\endcsname{{16.30}{62}{}{proposition.16.30}{Dirichlet-GRH-v930-Volume-III.pdf}}
\expandafter\def\csname r@cor:baseseparation723\endcsname{{45.21}{58}{}{proposition.45.21}{Dirichlet-GRH-v930-Volume-IX.pdf}}
\expandafter\def\csname r@cor:congcost728\endcsname{{12.26}{50}{}{proposition.12.26}{Dirichlet-GRH-v930-Volume-II.pdf}}
\expandafter\def\csname r@cor:fieldrelations712\endcsname{{12.4}{32}{}{proposition.12.4}{Dirichlet-GRH-v930-Volume-II.pdf}}
\expandafter\def\csname r@cor:labelledtrace735\endcsname{{13.16}{75}{}{proposition.13.16}{Dirichlet-GRH-v930-Volume-II.pdf}}
\expandafter\def\csname r@cor:levelcost725\endcsname{{45.24}{61}{}{proposition.45.24}{Dirichlet-GRH-v930-Volume-IX.pdf}}
\expandafter\def\csname r@cor:maincost730\endcsname{{12.33}{59}{}{proposition.12.33}{Dirichlet-GRH-v930-Volume-II.pdf}}
\expandafter\def\csname r@cor:markeddet715\endcsname{{47.37}{102}{}{proposition.47.37}{Dirichlet-GRH-v930-Volume-IX.pdf}}
\expandafter\def\csname r@cor:markedlocal801\endcsname{{17.34}{120}{}{proposition.17.34}{Dirichlet-GRH-v930-Volume-III.pdf}}
\expandafter\def\csname r@cor:normalizedtransfer722\endcsname{{45.17}{55}{}{proposition.45.17}{Dirichlet-GRH-v930-Volume-IX.pdf}}
\expandafter\def\csname r@cor:oddwindow677\endcsname{{11.30}{27}{}{proposition.11.30}{Dirichlet-GRH-v930-Volume-II.pdf}}
\expandafter\def\csname r@cor:oldperiod752\endcsname{{15.29}{32}{}{proposition.15.29}{Dirichlet-GRH-v930-Volume-III.pdf}}
\expandafter\def\csname r@cor:periodlabels764\endcsname{{16.20}{52}{}{proposition.16.20}{Dirichlet-GRH-v930-Volume-III.pdf}}
\expandafter\def\csname r@cor:quotientchainheight734\endcsname{{13.14}{73}{}{proposition.13.14}{Dirichlet-GRH-v930-Volume-II.pdf}}
\expandafter\def\csname r@cor:scalarfl724\endcsname{{12.16}{42}{}{proposition.12.16}{Dirichlet-GRH-v930-Volume-II.pdf}}
\expandafter\def\csname r@cor:sphericalcupbound762\endcsname{{16.16}{48}{}{proposition.16.16}{Dirichlet-GRH-v930-Volume-III.pdf}}
\expandafter\def\csname r@cor:threeextensions761\endcsname{{16.12}{44}{}{proposition.16.12}{Dirichlet-GRH-v930-Volume-III.pdf}}
\expandafter\def\csname r@cor:torustests741\endcsname{{15.6}{9}{}{proposition.15.6}{Dirichlet-GRH-v930-Volume-III.pdf}}
\expandafter\def\csname r@lem:basestabilizer723\endcsname{{45.18}{57}{}{proposition.45.18}{Dirichlet-GRH-v930-Volume-IX.pdf}}
\expandafter\def\csname r@lem:graphrealization724\endcsname{{12.14}{41}{}{proposition.12.14}{Dirichlet-GRH-v930-Volume-II.pdf}}
\expandafter\def\csname r@lem:inertcrys789\endcsname{{17.4}{85}{}{proposition.17.4}{Dirichlet-GRH-v930-Volume-III.pdf}}
\expandafter\def\csname r@lem:localevaluation728\endcsname{{12.27}{51}{}{proposition.12.27}{Dirichlet-GRH-v930-Volume-II.pdf}}
\expandafter\def\csname r@lem:sphere48\endcsname{{6.25}{58}{}{proposition.6.25}{Dirichlet-GRH-v930-Volume-I.pdf}}
\expandafter\def\csname r@lem:sphericalchain716\endcsname{{45.1}{37}{}{proposition.45.1}{Dirichlet-GRH-v930-Volume-IX.pdf}}
\expandafter\def\csname r@lem:sphericalcore252\endcsname{{75.1}{112}{}{proposition.75.1}{Dirichlet-GRH-v930-Volume-XIV.pdf}}
\expandafter\def\csname r@lem:tameorbit715\endcsname{{47.35}{101}{}{proposition.47.35}{Dirichlet-GRH-v930-Volume-IX.pdf}}
\expandafter\def\csname r@lem:tametraceheight734\endcsname{{13.12}{72}{}{proposition.13.12}{Dirichlet-GRH-v930-Volume-II.pdf}}
\expandafter\def\csname r@lem:wittconstants757\endcsname{{16.1}{36}{}{proposition.16.1}{Dirichlet-GRH-v930-Volume-III.pdf}}
\expandafter\def\csname r@nav:VII-67\endcsname{{123.6}{82}{}{subsection.123.6}{Dirichlet-GRH-v930-Volume-XXIII.pdf}}
\expandafter\def\csname r@prop:basechange258\endcsname{{71.17}{34}{}{proposition.71.17}{Dirichlet-GRH-v930-Volume-XIV.pdf}}
\expandafter\def\csname r@prop:blockrows719\endcsname{{45.7}{44}{}{proposition.45.7}{Dirichlet-GRH-v930-Volume-IX.pdf}}
\expandafter\def\csname r@prop:commonlift762\endcsname{{16.14}{47}{}{proposition.16.14}{Dirichlet-GRH-v930-Volume-III.pdf}}
\expandafter\def\csname r@prop:freetests741\endcsname{{15.4}{7}{}{proposition.15.4}{Dirichlet-GRH-v930-Volume-III.pdf}}
\expandafter\def\csname r@prop:heckecore252\endcsname{{71.12}{28}{}{proposition.71.12}{Dirichlet-GRH-v930-Volume-XIV.pdf}}
\expandafter\def\csname r@prop:highwidth788\endcsname{{17.1}{82}{}{proposition.17.1}{Dirichlet-GRH-v930-Volume-III.pdf}}
\expandafter\def\csname r@prop:normfactor723\endcsname{{45.19}{57}{}{proposition.45.19}{Dirichlet-GRH-v930-Volume-IX.pdf}}
\expandafter\def\csname r@prop:ramifiedorders782\endcsname{{16.44}{80}{}{proposition.16.44}{Dirichlet-GRH-v930-Volume-III.pdf}}
\expandafter\def\csname r@prop:righttransfer721\endcsname{{45.15}{52}{}{proposition.45.15}{Dirichlet-GRH-v930-Volume-IX.pdf}}
\expandafter\def\csname r@prop:split175\endcsname{{72.17}{60}{}{proposition.72.17}{Dirichlet-GRH-v930-Volume-XIV.pdf}}
\expandafter\def\csname r@sec:heckechains716\endcsname{{45.1}{36}{}{subsection.45.1}{Dirichlet-GRH-v930-Volume-IX.pdf}}
\expandafter\def\csname r@sec:highsmall789\endcsname{{17.2}{85}{}{subsection.17.2}{Dirichlet-GRH-v930-Volume-III.pdf}}
\expandafter\def\csname r@sec:mainorbit776\endcsname{{16.14}{70}{}{subsection.16.14}{Dirichlet-GRH-v930-Volume-III.pdf}}
\expandafter\def\csname r@sec:markedlevel801\endcsname{{17.13}{118}{}{subsection.17.13}{Dirichlet-GRH-v930-Volume-III.pdf}}
\expandafter\def\csname r@sec:markedmatrix715\endcsname{{47.12}{100}{}{subsection.47.12}{Dirichlet-GRH-v930-Volume-IX.pdf}}
\expandafter\def\csname r@sec:missing\endcsname{{6.6}{32}{}{subsection.6.6}{Dirichlet-GRH-v930-Volume-I.pdf}}
\expandafter\def\csname r@sec:oddcocycles675\endcsname{{11.8}{22}{}{subsection.11.8}{Dirichlet-GRH-v930-Volume-II.pdf}}
\expandafter\def\csname r@sec:ramifiedorders782\endcsname{{16.17}{79}{}{subsection.16.17}{Dirichlet-GRH-v930-Volume-III.pdf}}
\expandafter\def\csname r@sec:relationmatrices719\endcsname{{45.3}{43}{}{subsection.45.3}{Dirichlet-GRH-v930-Volume-IX.pdf}}
\expandafter\def\csname r@sec:splitattachment722\endcsname{{45.6}{54}{}{subsection.45.6}{Dirichlet-GRH-v930-Volume-IX.pdf}}
\expandafter\def\csname r@sec:tameends743\endcsname{{15.3}{11}{}{subsection.15.3}{Dirichlet-GRH-v930-Volume-III.pdf}}
\expandafter\def\csname r@seriesguide\endcsname{{1.1}{4}{}{subsection.1.1}{Dirichlet-GRH-v930-Volume-I.pdf}}
\expandafter\def\csname r@thm:allramifiedlifts761\endcsname{{16.11}{44}{}{proposition.16.11}{Dirichlet-GRH-v930-Volume-III.pdf}}
\expandafter\def\csname r@thm:allspherical798\endcsname{{17.28}{113}{}{proposition.17.28}{Dirichlet-GRH-v930-Volume-III.pdf}}
\expandafter\def\csname r@thm:alltorsion741\endcsname{{15.5}{8}{}{proposition.15.5}{Dirichlet-GRH-v930-Volume-III.pdf}}
\expandafter\def\csname r@thm:balancedcount728\endcsname{{12.28}{51}{}{proposition.12.28}{Dirichlet-GRH-v930-Volume-II.pdf}}
\expandafter\def\csname r@thm:charrows719\endcsname{{45.8}{44}{}{proposition.45.8}{Dirichlet-GRH-v930-Volume-IX.pdf}}
\expandafter\def\csname r@thm:commoncritical731\endcsname{{13.4}{62}{}{proposition.13.4}{Dirichlet-GRH-v930-Volume-II.pdf}}
\expandafter\def\csname r@thm:commonhecke716\endcsname{{45.3}{38}{}{proposition.45.3}{Dirichlet-GRH-v930-Volume-IX.pdf}}
\expandafter\def\csname r@thm:coreeuler252\endcsname{{71.2}{14}{}{proposition.71.2}{Dirichlet-GRH-v930-Volume-XIV.pdf}}
\expandafter\def\csname r@thm:criticalheights734\endcsname{{13.13}{72}{}{proposition.13.13}{Dirichlet-GRH-v930-Volume-II.pdf}}
\expandafter\def\csname r@thm:cuppath757\endcsname{{16.2}{36}{}{proposition.16.2}{Dirichlet-GRH-v930-Volume-III.pdf}}
\expandafter\def\csname r@thm:integrallifts759\endcsname{{16.5}{39}{}{proposition.16.5}{Dirichlet-GRH-v930-Volume-III.pdf}}
\expandafter\def\csname r@thm:labelledlift735\endcsname{{13.15}{74}{}{proposition.13.15}{Dirichlet-GRH-v930-Volume-II.pdf}}
\expandafter\def\csname r@thm:mainorbit776\endcsname{{16.38}{70}{}{proposition.16.38}{Dirichlet-GRH-v930-Volume-III.pdf}}
\expandafter\def\csname r@thm:markedmatrix715\endcsname{{47.36}{101}{}{proposition.47.36}{Dirichlet-GRH-v930-Volume-IX.pdf}}
\expandafter\def\csname r@thm:markedtype801\endcsname{{17.33}{119}{}{proposition.17.33}{Dirichlet-GRH-v930-Volume-III.pdf}}
\expandafter\def\csname r@thm:mixedcup772\endcsname{{16.33}{64}{}{proposition.16.33}{Dirichlet-GRH-v930-Volume-III.pdf}}
\expandafter\def\csname r@thm:negativeinterval740\endcsname{{15.3}{4}{}{proposition.15.3}{Dirichlet-GRH-v930-Volume-III.pdf}}
\expandafter\def\csname r@thm:oddinteger677\endcsname{{11.29}{26}{}{proposition.11.29}{Dirichlet-GRH-v930-Volume-II.pdf}}
\expandafter\def\csname r@thm:oldtorustraces743\endcsname{{15.8}{12}{}{proposition.15.8}{Dirichlet-GRH-v930-Volume-III.pdf}}
\expandafter\def\csname r@thm:ramifiedcut782\endcsname{{16.43}{79}{}{proposition.16.43}{Dirichlet-GRH-v930-Volume-III.pdf}}
\expandafter\def\csname r@thm:ribettrace720\endcsname{{45.11}{48}{}{proposition.45.11}{Dirichlet-GRH-v930-Volume-IX.pdf}}
\expandafter\def\csname r@thm:selftwist728\endcsname{{12.25}{50}{}{proposition.12.25}{Dirichlet-GRH-v930-Volume-II.pdf}}
\expandafter\def\csname r@thm:sparsecups780\endcsname{{16.41}{76}{}{proposition.16.41}{Dirichlet-GRH-v930-Volume-III.pdf}}
\expandafter\def\csname r@thm:splitattachment722\endcsname{{45.16}{54}{}{proposition.45.16}{Dirichlet-GRH-v930-Volume-IX.pdf}}
\expandafter\def\csname r@thm:tametransfer721\endcsname{{45.14}{51}{}{proposition.45.14}{Dirichlet-GRH-v930-Volume-IX.pdf}}
\expandafter\def\csname r@thm:tiedlength729\endcsname{{12.30}{53}{}{proposition.12.30}{Dirichlet-GRH-v930-Volume-II.pdf}}
\expandafter\def\csname r@thm:topdescent723\endcsname{{45.20}{58}{}{proposition.45.20}{Dirichlet-GRH-v930-Volume-IX.pdf}}
\expandafter\def\csname r@thm:twistedselmer730\endcsname{{12.32}{57}{}{proposition.12.32}{Dirichlet-GRH-v930-Volume-II.pdf}}
\expandafter\def\csname r@thm:twolatticecup768\endcsname{{16.27}{58}{}{proposition.16.27}{Dirichlet-GRH-v930-Volume-III.pdf}}
\expandafter\def\csname r@thm:unrestrictedcups762\endcsname{{16.15}{47}{}{proposition.16.15}{Dirichlet-GRH-v930-Volume-III.pdf}}
\makeatother
\begin{document}
\setcounter{part}{9}
\fancyhead[L]{\small\color{muted}Dirichlet GRH \textperiodcentered\ Volume X}
\fancyhead[R]{\small\color{muted}Torus periods and congruences}
\thispagestyle{plain}\begin{center}
{\LARGE\bfseries The generalized Riemann hypothesis\\[2pt]for Dirichlet \(L\)-functions}\\[6pt]
{\small 27 November 2026 \quad\textperiodcentered\quad Edition 930}\end{center}
\part{Projected torus periods and congruence modules}\label{part:X}
\begin{abstract}\noindent Local inertia characters and crystalline flags enter the projected torus pairing. We compare cup-product bounds over the period fields and relate the integral period to the Hecke congruence module.\end{abstract}
\setcounter{section}{49}
\noindent The \hyperref[seriesguide]{guide in Volume~I} lists all twenty-seven volumes.
\begingroup\small\setlength{\parskip}{0pt}\tableofcontents\endgroup
\clearpage

\section{Local constituents of the projected torus pairing}\label{nav:torus-constituents}\label{nav:volume-X}

\subsection{The integral pairing of two tests projected to a torus lift}
\label{sec:rankonepair736}

In passing from a Hecke multiplier to a period, the tests for the
lift from \(U(1)\) are initially orthogonal to certain forms in the
kernel of the entire descent. The projection used in the period is
onto the lift of a single source character. Averaging one test over
the source classes relates these two spaces. We compute the integer
introduced by the averaging on the fields of the shorter family.

\begin{proposition}[The number of source classes at the two marked levels]
\leavevmode\label{prop:torusorder736}
Put \(F=E^+\), \(f=[F:\QQ]\), and \(T=U(1)\) on the fields of
Corollary~\ref{cor:fieldrelations712}. Write \(U\) for the norm-one
unit subgroup at every finite place, with the additional condition
that the component reduce to \(1\) at each marked split place over
\(p_0,p_1\). Let \(j:\mathrm{Cl}(F)\to\mathrm{Cl}(E)\) be extension
of ideals and let \(\mu_E\) be the roots of unity of \(E\). Then
\[
 C=T(F)\backslash T(\mathbb A_{F,f})/U,\qquad
 |C|=
 \frac{h_E}{|\operatorname{im}j|}
 \frac{H}{|\mu_E|},\qquad H=((p_0-1)(p_1-1))^f .
\]
All stabilizers in this class sum are trivial. For each coefficient
prime \(\ell\), representatives \(z_c\) of all classes can be chosen
integral and invertible at every place above \(\ell\).
\end{proposition}
\begin{proof}
First use the full norm-one unit subgroup \(U_0\). The map
\(a\mapsto a/\bar a\) from finite ideles of \(E\) onto
\(T(\mathbb A_{F,f})\) has kernel the finite ideles of \(F\).
Indeed Hilbert's Theorem 90 gives the assertion at a field place,
and at a split place one can lift \((z,z^{-1})\) to \((z,1)\).
At an inert place the extension is unramified, so multiplication
by a power of a uniformizer of \(F_v\) makes any lift of a unit
into a unit. These are all the
finite cases by Corollary~\ref{cor:fieldrelations712}. Thus the
map on products of unit groups is also onto \(U_0\). Global Hilbert
90 gives the same surjection on rational points. The quotient with
level \(U_0\) is therefore the ideal class group of \(E\) modulo
the images of finite ideles of \(F\), namely
\(\mathrm{Cl}(E)/\operatorname{im}j\).

Reduction gives \([U_0:U]=H\), since there are \(f\) marked
places of \(F\) of residue cardinality \(p_i\) over each \(p_i\).
The kernel of the map from \(U_0/U\) to \(C\) is the image of
\(T(F)\cap U_0\). This intersection equals \(\mu_E\) because a norm-one
integral unit has absolute value one at every complex embedding,
so is a root of unity by Kronecker's theorem, and every root of
unity has norm one. A root equal to \(1\) modulo both marked
characteristics has order a power of each characteristic, by
Lemma~\ref{lem:tametraceheight734}. It must equal \(1\).
We get the asserted order and \(T(F)\cap U=1\), which gives the
stabilizer statement because the group is abelian.
Finally lift a given representative to an idele \(a\) of \(E\).
By weak approximation choose \(b\in E^\times\) such that \(ba\)
is a local unit everywhere above \(\ell\). Replacing \(a/\bar a\)
by \(ba/\overline{ba}\) preserves its class and has the required
integral components.
\end{proof}

Write \(W\) for the finite space of definite target forms in the
rank-one lift, with the Hermitian class-sum pairing linear in the
second variable. Absorb the chosen source character into the theta
kernel, so that the selected vector of the source is the constant
function on \(C\). Write \(K_c\) for the target kernel at \(c\)
and let \(\phi=\sum_{c\in C}K_c\). Common nonzero changes of
normalization of the kernel give the same line
\(Z=\mathbb C\phi\) and the same kernel of descent. Write \(P\)
for orthogonal projection onto \(Z\) and \(D\) for descent,
so that \(D y(c)=\langle K_c,y\rangle\) up to a common
nonzero scalar. At the marked places \(e\) is the orthogonal
projection given by the mean of the finite residue-unit action
for the trivial source type. Its group has order \(H\).
For a split place \(v\) off the level let \(A_v\) be the reversed
spherical Hecke polynomial evaluated at the inverse of the top
root prescribed for this lift. A tested vector is orthogonal to
\(\operatorname{im} e\cap\ker A_v\cap\ker D\) for all such \(v\).

\begin{lemma}[Averaging one test]\leavevmode
\label{lem:centraltest736}
For each \(\ell\) there is a sum \(B\) of \(|C|\) normalized
central translations of the target satisfying
\[
 PB=|C|P,\qquad
 B\Xi\perp \operatorname{im}e\cap\ker A_v\cap Z^\perp
                  \quad\text{for every tested }\Xi\text{ and every }v .
\]
The sum preserves the lattice of integral forms at each embedding
above \(\ell\).
\end{lemma}
\begin{proof}
For a finite adelic norm-one scalar \(z\), take right translation
by the scalar on the target with the unitary splitting factor,
choosing the inverse convention so that its action \(U_z\) on
the kernel is \(U_zK_c=K_{cz}\). Indeed the tensor action of a
source and a target scalar of the same value on the underlying
homomorphism space is \(x\mapsto zxz^{-1}=x\). Their joint Weil
action is the scalar of the chosen splitting. Dividing it out
moves the source translation to the target, and the source
character absorbed into the kernel contributes its own unitary
factor. This proves the identity by applying the action to the
theta sum term by term. In particular \(U_z\phi=\phi\).
These translations and their inverses commute with the residue
unit action and with all \(A_v\), since their group elements are
central. The trivial source type gives \(e\phi=\phi\) because the
marked action on the lifted constant is the normalized translation
of the constant on the source, which fixes it. Equivalently the
finite change of variable in the source sum permutes its classes.

Take \(B=\sum_{c\in C}U_{z_c}\) using the representatives of the
proposition. Each summand is unitary and fixes \(\phi\), proving
\(PB=|C|P\). If \(y\perp Z\), then for every \(d\in C\),
\[
  D(B^*y)(d)=\langle B K_d,y\rangle
             =\Big\langle\sum_{c\in C}K_{d c},y\Big\rangle=0
\]
in the class-sum normalization. If \(y\) is also in
\(\operatorname{im}e\cap\ker A_v\), then \(B^*y\) stays in this
intersection. Testing \(\Xi\) against \(B^*y\) proves the claimed
orthogonality. Only the action on the kernels is indexed by the
finite quotient \(C\). The representatives define actual
translations on the whole target.

For integrality the scalar matrix of \(z_c\) and its inverse
preserve the coefficient lattice at every place above \(\ell\).
The splitting and source characters used for normalization have
\(\ell\)-adic unit values on these representatives. To check
this last assertion, raise an idele integral above \(\ell\) to
a positive power so that it is a rational torus point times
a compact unit, using finiteness of \(C\). Both factors remain
units above \(\ell\). A smooth character on the compact unit
subgroup has finite image. On the rational point the finite
character value is prescribed by the infinity type, a product
of powers of its conjugates (after squaring for a half-integral
normalization). Each conjugate is an \(\ell\)-adic unit. The
original value is therefore a unit as well. Applying these
two observations to the translated value of a form at each
adelic representative proves integrality of the sum.
\end{proof}

\begin{theorem}[From a trace multiplier to a projected inner product]
\leavevmode\label{thm:rankonepair736}
Fix the shorter period field, a twist and a coefficient embedding
\(\iota:\mathbb C\simeq\overline{\QQ}_\ell\). Use the rank-one
theta data just described at target rank \(n>1\). For a nonzero
\(\phi\) take its base Hecke system \(x_*\), with the split
normalized representation and auxiliary row hypotheses of
Corollary~\ref{cor:baseseparation723} for \(X_0=\{x_*\}\).
Let \(b_\ell\) be any of the bounds for \(v(\beta_{x_*})/v(\ell)\)
in Theorem~\ref{thm:ribettrace720}, including the replacement for
the two original level terms in Corollary~\ref{cor:levelcost725}
when applied. Let \(c_0\) be the stabilizer multiple of
Lemma~\ref{lem:basestabilizer723} for this \(n\). For any two
tested forms \(\Xi,\Xi'\) integral at all embeddings above \(\ell\),
\[
  \bigl|\iota\langle P\Xi,P\Xi'\rangle\bigr|_\ell
       \le \ell^{\,b_\ell+v_\ell(Hc_0^2|C|)} .
\]
For \(\phi=0\) the left side is zero. In particular the bound
counts the trace loss only once.
\end{theorem}
\begin{proof}
Put \(u=B\Xi'\), \(l=eu\) and \(y=Pu\). The vector \(Hl\)
is integral because \(He\) is the sum of the normalized residue-unit
actions, each of which preserves the integral function lattice.
The local matrices are integral with integral inverses, and the
finite character multipliers are roots of unity. The line of the
constant lift has a single good Hecke system, by the spherical
theta parameter formula of Section~\ref{sec:heckechains716}
and the unramified transfer used in Theorem~\ref{thm:coreeuler252}.
Hence \(y\) is supported on \(x_*\).

We check the agreement on every base system with a trivial
constituent, including possible multiplicities in the target.
At any good split \(v\nmid\ell\) take the geometric Frobenius
in that constituent. By Theorem~\ref{thm:splitattachment722}
the arithmetic-normalized spherical polynomial has the root
\(C_v=q_v^{1/2}\mu_v(\varpi_v)\), in the notation of
Corollary~\ref{cor:normalizedtransfer722} with source dimension
one. This follows because the normalizing character at that
Frobenius equals \(C_v^{-1}\) under the coefficient
identification. By Lemma~\ref{lem:sphericalchain716} the same
polynomial has the root \(C_v q_v^{\,n-2}\). Reversing it and
evaluating at the inverse of this root kills that whole system
by \(A_v\).

The commuting spherical Hecke algebra is closed under adjoints.
Its joint eigenspaces are mutually orthogonal, and the marked
residue action commutes with it. Each joint projector \(P_x\)
therefore commutes with \(e\) and with \(P\), since \(Z\) is an
eigenline. Also \(eP=P\). For a system with a trivial constituent,
\[
       d_x=P_x(l-y)=eP_x(u-Pu)
               \in\operatorname{im}e\cap\ker A_v\cap Z^\perp .
\]
The equality uses the commutation just proved. The lemma gives
\(\langle d_x,u\rangle=0\), and orthogonality to \(Z\) gives
\(\langle d_x,Pu\rangle=0\). Since \(eP_x\) is itself an
orthogonal projection, \(\|d_x\|^2=\langle d_x,u-Pu\rangle=0\).
This verifies the agreement required in base separation.

Apply Corollary~\ref{cor:baseseparation723} with \(h=H\).
It gives a scalar \(s_*\) on the chosen system, with
\(v_\ell(\iota s_*)\le b_\ell\), for which
\(Hc_0s_* y\) is integral at \(\iota\). To pair this form with
\(\Xi\), write the definite inner product as a sum over double
cosets of the coefficient pairing divided by the point
stabilizer order, using volume one for the defining level.
In the tensor basis the coefficient pairing has identity matrix.
Taking inner products of columns of the integral inclusion
of the Weyl lattice gives its integral Gram matrix. In the transported
coordinates of a representative, the two changes of basis at
\(\iota\) and \(\iota\circ\mathrm{conj}\) cancel in this
pairing by the unitary relation \(\bar g^{\,t}=g^{-1}\).
Thus every coefficient pairing is integral when the first
vector is integral at \(\iota\circ\mathrm{conj}\) and the
second at \(\iota\). Every stabilizer order divides \(c_0\)
by Lemma~\ref{lem:basestabilizer723}. Consequently
\[
   \iota\bigl(Hc_0^2 s_* \langle\Xi,y\rangle\bigr)
                   \in\overline{\ZZ}_\ell .
\]
Now \(y=|C|P\Xi'\), and self-adjointness and idempotence give
\(\langle\Xi,P\Xi'\rangle=\langle P\Xi,P\Xi'\rangle\).
Taking absolute values proves the estimate.
\end{proof}

\begin{corollary}[Total cost of the source classes and stabilizers]
\leavevmode\label{cor:rankonefixed736}
At every node to which the theorem applies, the additional
denominator has total logarithmic cost
\[
 \sum_{\ell}v_\ell(Hc_0^2|C|)\log\ell
 \le \log h_E+2f\log(p_0p_1)+2n^2\min_i\log p_i .
\]
This bound is independent of the auxiliary moduli of the twists
chosen separately at the coefficient embeddings.
\end{corollary}
\begin{proof}
The three integers \(H,c_0,|C|\) use the original levels.
Proposition~\ref{prop:torusorder736} gives \(|C|\mid h_E H\).
Use \(H\le(p_0p_1)^f\) and the bound for \(c_0\) in
Lemma~\ref{lem:basestabilizer723}, then factor the fixed integer
into primes. The upper bound for \(b_\ell\) in the theorem
continues to use the actual row and local cohomology estimates.
Corollary~\ref{cor:labelledtrace735} specifies which coefficient
completions can be omitted in them.
\end{proof}
\clearpage

\subsection{The rank-one period at the spherical coefficient primes}
\label{sec:rankonespherical737}

For a main pair with critical integer \(z\), the lift from \(U(1)\)
has target rank \(n=2z\). The string in its normalized parameter
has odd length \(t'=2z-1\). This parity determines the global
congruence term at every coefficient prime and eliminates it at
every odd prime unramified in the period field. We use this fact
to bound the inner products of the original tested forms of
Section~\ref{sec:rankonepair736}.

\begin{proposition}[The cyclotomic ideal for an odd string]\leavevmode
\label{prop:oddcyclo737}
Let \(E=\QQ(\mu_{q_1},\sqrt{-d})\), where \(4\mid q_1\),
\(d\equiv7\bmod 8\) is prime and \(d\nmid q_1\), as in the
shorter family. Set \(N=q_1d\). Let \(t'\) be odd and let
\(T\) be any continuous rank-one \(\ell\)-adic lattice character
of \(G_E\). The normalized valuation \(a_\ell(t')\) of the ideal
\(I_{t'}\) in Theorem~\ref{thm:selftwist728} is
\[
 a_\ell(t')=
 \begin{cases}
  v_\ell(q_1)+v_\ell(t'), &\ell\mid q_1,\\
  1+v_d(t'), &\ell=d\ \text{and }(d-1)/2\mid t',\\
  0,&\text{otherwise}.
 \end{cases}
\]
In particular, for an integer \(t_0\ge1\), the sum of the suprema
over odd \(t'\le t_0\), weighted by \(\log\ell\), is at most
\(\log\big(N\operatorname{lcm}(1,\ldots,t_0)\big)\).
\end{proposition}
\begin{proof}
A rank-one lattice character has unit trace on every group
element. The ideal is therefore generated by \(1-u^{t'}\) for
\(u\) in the cyclotomic image of \(G_E\). The cyclotomic Galois
group acts on roots of unity through the product of the groups
\(\ZZ_p^\times\). Fixing \(E\) requires that its component
at each \(p\mid q_1\) be \(1\) modulo \(p^{v_p(q_1)}\), and
that the component at \(d\) reduce to a square. The last
condition uses the quadratic subfield of \(\QQ(\mu_d)\), which
has discriminant \(-d\). These conditions are independent by the
Chinese remainder theorem on every finite cyclotomic level.

For \(\ell^b\Vert q_1\) the image is \(1+\ell^b\ZZ_\ell\).
Here \(b\ge2\) at two. The least valuation of \(u^{t'}-1\)
for \(u\) in this subgroup is \(b+v_\ell(t')\), attained at
\(1+\ell^b\). Indeed a power prime to \(\ell\) preserves the
depth of a principal unit, and an \(\ell\)-th power increases
it by one by the binomial expansion, for \(\ell\) odd or depth
at least two. At \(d\) the residual image has order \((d-1)/2\).
If that order divides \(t'\), the residual roots of unity lift
to roots of the same order and are killed by the power \(t'\).
The minimum on the principal units is then \(1+v_d(t')\).
Otherwise a lifted generator gives valuation zero. At every
other prime the image is the whole unit group and the prime is
odd. The element \(-1\) gives valuation zero because \(t'\) is
odd. Finally the product of the prime powers given by the
suprema divides \(N\operatorname{lcm}(1,\ldots,t_0)\), which
gives the sum bound.
\end{proof}

\begin{theorem}[Simultaneous rank-one Selmer bounds at every prime]
\leavevmode\label{thm:rankoneselmer737}
On these fields take every main character of the designated torus
lifts at a node, with tame part \(\omega_w^{\pm j}\) at each old
marked place, \(1\le j\le m\), in the notation of
Corollary~\ref{cor:quotientchainheight734}. For their untwisted
lattice characters \(T\), the constants \(B_i\) of
Theorem~\ref{thm:twistedselmer730} equal \(a_\ell(2z-1)\).
Thus, for the same twists and signed powers that control all the
critical ratios, the corresponding normalized divisible Selmer
lengths with the integral crystalline evaluation equations are
at most \(a_\ell(2z-1)+\epsilon_\ell\).
\end{theorem}
\begin{proof}
Write \(\nu\) for the source torus character in the main pair and
put \(u(a)=\nu(a/\bar a)^{-1}\) on the ideles of \(E\). The
untwisted rank-one main character is
\(T=u_\ell\kappa^{1-z}\), where \(u_\ell\) is the algebraic
avatar of \(u\). Indeed the source parameter of the constant is
one. In the normalization of Corollary~\ref{cor:normalizedtransfer722}
the \(2z-1\) added Satake entries at a good split place become
\(1,q,\ldots,q^{2z-2}\). The remaining entry is the inverse
source torus character on \(a/\bar a\) times \(q^{z-1}\).
Indeed the unitary entry is the splitting-character ratio, and the
arithmetic weight shift after dividing by the rank-one normalizer
is \(((2z-1)-1)/2=z-1\). At geometric Frobenius this is exactly
the value of \(T\). Replacing \(\nu\) by \(\nu\vartheta^s\)
multiplies the remaining entry by the inverse pullback of
\(\vartheta^s\). The unramified theta parameter identity used in
Theorem~\ref{thm:coreeuler252}, followed by Chebotarev and
semisimplicity as in Theorem~\ref{thm:splitattachment722}, gives
the direct sum of this character with the stated string whenever
the lift is nonzero.

The equality \(u^c=u^{-1}\) gives \(T^c=T^{-1}\kappa^{1-t'}\).
The character is de Rham above \(\ell\) since it is the avatar
of an algebraic Hecke character. Away from \(\ell\) it is
unramified except at the two marked levels. Take \(S\) to
consist of just these levels, their conjugates and the places
above \(\ell\). At every \(w\in S\) away from \(\ell\),
\(\kappa\) is unramified and the full tame generator has value
of order \(Q/\gcd(Q,j)\) under \(T\), also at the conjugate
place. In the proof of Corollary~\ref{cor:quotientchainheight734},
if \(p\) is either \(2\) or \(3\), the \(p\)-part of this order
is at least \(p^{h+6}/m>2\). Taking \(p\ne\ell\) shows that
the residual order is greater than two. Both \(T(g_*)-1\) and
\(T(g_*)^2-1\) are units for a tame generator at any chosen
old place away from \(\ell\). The same holds at its conjugate
inverse. This checks the marked scalar hypotheses. The full
matrix span in rank one follows already from the identity matrix.
On the Tate dual the inertia value is the inverse of that on
\(T\). The difference from one is still a unit, so
\(H^0(E_w,W_b^*)=0\) at every finite precision \(b\) for all the
old terms away from \(\ell\).

In the formula for \(B_i\), both distinguished determinants and
every local invariant term consequently have valuation or length
zero. Proposition~\ref{prop:oddcyclo737} gives the remaining ideal
term exactly. Apply Theorem~\ref{thm:criticalheights734} with
these finitely many lattices and all the critical pairs. This
also gives any prescribed positive bound \(\eta_\ell\) for the
total new tame inertia cost, with the same choice of twist.
\end{proof}

\begin{theorem}[The contribution of the spherical primes to projected periods]
\leavevmode\label{thm:rankonespherical737}
At a node on the shorter period interval use \(t_0=5r\),
\(Y=N+t_0+1\) and \(G_0\) from Corollary~\ref{cor:actualmarks714}.
Let \(\mathcal G\) consist of the odd coefficient primes prime to
\(Np_0p_1\) where all the rank-one lift data have hyperspecial
level above \(\ell\) and the labelled intervals of every relevant
trace family and scalar or main target satisfy
Theorem~\ref{thm:labelledlift735} at every completion and
coefficient embedding. This set contains every sufficiently large
prime at the node and has a common such threshold on each finite
interval. Given \(\delta>0\), the twists can be chosen so that
\[
\begin{split}
 &\sum_{\ell\in\mathcal G}
    \sup_{\iota,\,\nu,z,s,\,\Xi,\Xi'}
       \log\max\bigl(1,|\iota\langle P\Xi,P\Xi'\rangle|_\ell\bigr)\\
 &\qquad\le
     2\log h_E+\log\mathcal J_{N,t_0}
       +\bigl(6+(4t_0+3)Y^6\bigr)\log G_0+\delta .
\end{split}
\]
Here the supremum is over the original integral tested forms of
Section~\ref{sec:rankonepair736} and the nonzero signed twist powers
of magnitude at most \(m\) of the rank-one main data. Zero lifts give
zero contributions. The same twists can give the bound of
Corollary~\ref{cor:quotientchainheight734} for the chain of critical
values. All quantities on the right are independent of their
auxiliary moduli.
\end{theorem}
\begin{proof}
At a prime of \(\mathcal G\) one has \(a_\ell(t')=0\). Choose
\(\epsilon_\ell,\eta_\ell>0\) so that
\(\sum_\ell(\epsilon_\ell+2t_0\eta_\ell)\log\ell\le\delta\).
One may decrease \(\epsilon_\ell\) further for the product of the
critical ratios. By Corollary~\ref{cor:labelledtrace735} all
scalar colengths vanish, and the normalized main trace length
is at most \(\epsilon_\ell\) by Theorem~\ref{thm:rankoneselmer737}.
The large-prime assertion follows since the finitely many infinity
types give bounded jumps at every embedding, independently of all
auxiliary levels. The twists split above their respective
coefficient primes. The exclusions at the fixed levels and at the
ramified primes of the period fields are finite. Take maxima over
the node data on any one finite interval.

We verify the remaining row inputs for these rank-one targets.
At a twist place the character of the parahoric is the inverse
tame twist on the determinant of the one-dimensional quotient,
and it is trivial on the other block. In the torus datum a block
triangular transformation on the quotient line changes the
test by the inverse residue character of that scalar. The same
parahoric condition holds in every system at the later trace
levels, since those places are disjoint from the twist places.
All the selected powers ramify at every twist pair by
Theorem~\ref{thm:commoncritical731}. The nontrivial
line-character classification in the proof of
Theorem~\ref{thm:tametransfer721} gives a single ramified line
with this inertia character and an unramified spherical complement,
with zero monodromy. The conclusion for a quotient line follows
from the version for a subline by inverse-transpose, which inverts
the characters. The normalizer comes from the source splitting
character with unramified smooth part. By
Theorem~\ref{thm:splitattachment722} the actual normalized
inertia matrices are therefore killed by
\((X-1)(X-\xi_w)\) at a chosen tame generator, with
\(\xi_w\) the indicated twist value or its inverse.
At the conjugate place the same assertion follows by duality.
This supplies \(o_w(X)=X-\xi_w\) for
Theorem~\ref{thm:charrows719}. Its contribution after normalization
is \(t'v_\ell(1-\xi_w)\). Summing over both sides of all twist
pairs costs at most \(2t'\eta_\ell\).

At the original marked levels use the Frobenius rows of
Corollary~\ref{cor:levelcost725}. Together these cost at most
\(2v_\ell(G_0)\). Away from the marked, twist, trace and
coefficient places the level is hyperspecial. The spherical
base change and the zero-monodromy compatibility in the proof
of Theorem~\ref{thm:splitattachment722} give unramified
representations there. For the selected main summand the
prime choices of Corollary~\ref{cor:markeddet715} provide the
unit determinant required at the subsidiary trace primes.
The unramified parameters of the selected lift have the direct
sum found above. Thus Corollary~\ref{cor:normalizedtransfer722}
applies to the full trace family, including conjugate primes.
The construction in Section~\ref{sec:relationmatrices719}
adjoins any remaining generators of the matrix span independently
of the row estimate. These choices verify
the row inputs for base separation on the actual selected line.

Theorem~\ref{thm:oddinteger677} bounds the odd-cocycle term by
\(v_\ell(h_E\mathcal J_{N,t_0})\), since the scalar weights
used here form a subset of \(0\le a<t_0\). The final interpolation loss in
Theorem~\ref{thm:ribettrace720} is at most
\((4t_0+3)Y^6v_\ell(G_0)\), by Corollary~\ref{cor:actualmarks714}.
Adding all the terms gives the bound in
Theorem~\ref{thm:rankonepair736} with
\[
 b_\ell=v_\ell(h_E\mathcal J_{N,t_0})
  +\big(2+(4t_0+3)Y^6\big)v_\ell(G_0)
  +\epsilon_\ell+2t_0\eta_\ell .
\]
Finally \(n=2z\le 2r\le r t_0\). The standard block inclusions
of finite general linear groups give \(c_0\mid G_0\), and the
split residue group orders give \(H\mid G_0\).
Proposition~\ref{prop:torusorder736} gives
\(Hc_0^2|C|\mid h_E G_0^4\). Substitute this and \(b_\ell\)
into the pairing estimate, take the supremum, and factor these
fixed integers over the primes. This proves the stated sum.
\end{proof}

The theorem bounds the rank-one projected inner products at every
included coefficient prime, on the actual fields of the period
calculation. Bounding the omitted coefficient places as the interval
grows and comparing the full chain of higher-rank projections with
its archimedean periods are further steps in Section~\ref{sec:missing}.
\clearpage

\subsection{The two inertia characters on the twist parahoric}
\label{sec:twistblocks738}

At a new split place the period test singles out a subspace of the
residue vector space. A change of basis of this subspace acts on
the test through the tame character of its determinant. We first
determine the local Langlands parameter from this transformation
law, including when both parts of the flag have dimension greater
than one. We then measure the cost of splitting an integral Galois
lattice along the two inertia characters.

\begin{theorem}[A nontrivial determinant character on one parahoric block]
\leavevmode\label{thm:twistblocks738}
Let \(K\) be a nonarchimedean local field of characteristic zero,
with ring of integers \(O\), residue field \(k\) and uniformizer
\(\varpi\). Write \(n=a+b\) with \(a,b>0\). Let \(P\subset\GL_n\)
be the standard upper parabolic of type \((a,b)\) and let \(J\)
be the inverse image of \(P(k)\) in \(\GL_n(O)\). For a
nontrivial character \(\xi:k^\times\to\CC^\times\) let
\(\lambda(j)=\xi(\det\bar j_{22})\) on \(J\), and extend \(\xi\)
to a character \(\widetilde\xi\) of \(K^\times\) with value one on
\(\varpi\). If an irreducible smooth complex representation \(\pi\)
of \(\GL_n(K)\) contains a vector of type \((J,\lambda)\), then
\[
 \pi\simeq i_P^{\GL_n(K)}
       \bigl(\sigma_a\boxtimes
                 ((\widetilde\xi\circ\det)\otimes\sigma_b)\bigr)
\]
for irreducible spherical representations \(\sigma_a,\sigma_b\).
Here \(i_P\) denotes normalized induction. The space of vectors of
this type in \(\pi\) is one-dimensional. The Weil--Deligne
parameter of \(\pi\) has monodromy zero and inertia representation
\(\mathbf1^{\oplus a}\oplus\xi_I^{\oplus b}\), where
\(\xi_I=\widetilde\xi\circ\operatorname{Art}_K^{-1}\) on inertia
with the reciprocity convention of Section~\ref{sec:splitattachment722}.
\end{theorem}
\begin{proof}
Restrict the given vector to the standard Iwahori, the inverse
image of the upper triangular group. It has diagonal characters
\(1,\ldots,1,\xi,\ldots,\xi\), in the indicated order.
The tame Iwahori type theorem, in the form proved using Roche's
theorem in the proof of Clozel--Harris--Taylor, Lemma~3.1.6,
puts \(\pi\) in a normalized principal series as a composition
factor. Its inducing characters \(\gamma_1,\ldots,\gamma_n\)
restrict to \(O^\times\) in precisely the indicated order.
The theorem applies to \(\GL_n\) in every residue
characteristic.~\cite[Lemma 3.1.6 and its proof]{X-VII-VII-chttransfer721}

We calculate the dimension of the \((J,\lambda)\)-space in this
entire principal series. A function in it is determined by its
values on \(B(K)\backslash\GL_n(K)/J\), with \(B\) the upper
Borel. By Iwasawa decomposition followed by reduction of the
maximal compact subgroup, this set is \(B(k)\backslash\GL_n(k)/P(k)\).
Its permutation representatives are indexed by the subsets of size
\(b\) giving the coordinates of the second block. The diagonal
unit group belongs to the stabilizer of each representative.
An arbitrary unit in a single coordinate shows that a value can
be nonzero only if exactly the last \(b\) inducing characters are
assigned to the second block. Since \(\xi\ne1\), precisely one
representative survives. At that representative, which can be
taken to be the identity, the characters agree on \(B(K)\cap J\).
Indeed its diagonal entries are units, the modulus is one and
both characters are trivial on the upper unipotent. A function
supported on \(B(K)J\) with value one at the identity therefore
exists. This proves the dimension assertion for the principal series.

For the first \(a\) inducing characters take the spherical
constituent \(\sigma_a\) of their principal series. For the
remaining \(b\), divide them first by \(\widetilde\xi\) and take
the spherical constituent \(\sigma_b\). In each case Iwasawa
decomposition gives a one-dimensional space of maximal compact
fixed vectors. Compact averaging makes passage to invariants exact,
so the spherical constituent exists and occurs once. Its parameter is
unramified with monodromy zero, by the spherical correspondence.
Normalized induction is exact and transitive, so the product in the
statement is a subquotient of the original principal series. This
product is irreducible by Zelevinsky's criterion for two collections
of segments. Every character in the first collection is trivial
on units, and every character in the second restricts to \(\xi\).
Multiplication by any power of the norm preserves this restriction.
Consequently no segment in one collection can be linked to a segment
in the other.~\cite[Prop.~8.5]{X-VII-zelevinsky738}

In this irreducible product, take a section on \(P(K)J\) with
value at one the tensor of the two spherical vectors, with the
second factor twisted as above. The inducing action of \(P(K)\cap J\)
on this value is \(\lambda\), which gives the required
equivariance of the section. Since \(J\) is compact, averaging
against \(\lambda^{-1}\) makes its character space an exact
functor in characteristic zero. The dimension one just computed
for the whole principal series can thus be accounted for by only
one composition factor, with multiplicity one. This identifies
\(\pi\) with the product. Compatibility of local Langlands with
Langlands quotients gives the direct sum of the two spherical
parameters with the second twisted by \(\widetilde\xi\). More
explicitly, interchanging factors from opposite collections preserves
the induction since their segments are unlinked. This puts the
combined Langlands data in standard order, and the quotient
then has the union of the two parameter multisets. This proves
the claimed restriction and monodromy.
\end{proof}

\begin{proposition}[Integral splitting of the two inertia spaces]
\leavevmode\label{prop:inertialattice738}
Let \(v(\ell)=1\) on a finite coefficient field \(L/\QQ_\ell\)
with ring of integers \(\mathcal O\). Suppose inertia acts on
\(V=V_0\oplus V_1\) by the two characters \(1,\xi_I\) in the
theorem, of dimensions \(a,b\), with values in \(L\).
For any stable lattice \(\Lambda\subset V\) put
\(\Lambda_i=\Lambda\cap V_i\). If \(\tau\) maps to a generator of
the tame character image and \(c=\xi_I(\tau)-1\), then
\[
       c\Lambda\subseteq\Lambda_0\oplus\Lambda_1,\qquad
       \operatorname{len}_{\rm norm}
          \bigl(\Lambda/(\Lambda_0\oplus\Lambda_1)\bigr)
                     \le\min(a,b)\,v(c).
\]
Here the normalized length of a cyclic module \(\mathcal O/u\)
is \(v(u)\). The bound is attained among stable local lattices
with the given ranks and tame character.
\end{proposition}
\begin{proof}
Let \(A\) denote the action of \(\tau\) and \(P_i\) the projection
onto \(V_i\). Then \(cP_1=A-1\) and \(cP_0=c-(A-1)\) preserve
\(\Lambda\). Their images lie in \(\Lambda_1,\Lambda_0\),
respectively. Projection along either summand identifies the stated
quotient with \(P_i\Lambda/\Lambda_i\). In fact it is onto, and if
\(P_i x\in\Lambda_i\) then \(x-P_i x\) is in the other intersection.
This quotient is killed by \(c\) and generated by at most
\(\dim V_i\) elements because \(P_i\Lambda\) is a lattice in
\(V_i\). Taking the smaller of the two bounds proves the inequality.

For sharpness, put \(z=\xi_I(\tau)\). Take
\(\min(a,b)\) copies of the matrix
\(\left(\begin{smallmatrix}1&1\\0&z\end{smallmatrix}\right)\)
on \(\mathcal O^2\), and scalar blocks for the remaining dimensions.
The two eigenlattices in a displayed block have respective bases
\((1,0)\) and \((1,z-1)\), and their sum has determinant \(z-1\).
These matrices define an action of the cyclic tame image, since
the sum of all powers of \(z\) up to its order minus one is zero.
They extend to a local representation by the quotient character
\(\widetilde\xi\circ\operatorname{Art}_K^{-1}\).
\end{proof}

\begin{corollary}[The total twist-place cost in all block ranks]
\leavevmode\label{cor:alltwistcost738}
At a node of the shorter family, take the split twist places of
Corollary~\ref{cor:quotientchainheight734}, with every prescribed
nonzero signed power. For each of the finitely many descent types
\(\alpha\), denote by \(b_\alpha\) the dimension of its marked
block, by \(a_\alpha\) the complementary dimension and by
\(t_\alpha\) the number of scalar characters used in
Theorem~\ref{thm:charrows719}. Use the integral space of forms of
the prescribed parahoric character at every twist place and all
the auxiliary trace levels away from those places. Let \(u_x\)
be the normalized Galois representation of any constituent of
one of these spaces. Write \(\zeta_w\) for the value at a full
tame generator \(\tau_w\) of the prescribed nontrivial local
character, transported to the coefficient field. At the conjugate
place use the conjugate dual character. Then
\[
          (u_x(\tau_w)-1)(u_x(\tau_w)-\zeta_w)=0 .
\]
In the row construction one may use \(o_w(X)=X-\zeta_w\)
in all these ranks. If \(d_{x,w}\) is the normalized length in
Proposition~\ref{prop:inertialattice738} for any stable lattice
in \(u_x\), then, given \(\delta>0\), the same twists used for
the critical ratios can be chosen to satisfy
\[
 \sum_\ell(\log\ell)\,
 \sup_{\iota,s,\alpha,x,\Lambda}
 \sum_{w\in Q\cup Q^c}
       \bigl(t_\alpha v_\ell(1-\zeta_w)+d_{x,w}\bigr)
       \ \le\ \delta .
\]
The set \(Q\) contains one place from each new pair for the twist
belonging to \((\ell,\iota)\). Empty sets contribute zero. The
same bound is available at every node in a finite common period
interval.
\end{corollary}
\begin{proof}
At a twist place the finite test in the integral frames asks that
the reduction of the matrix map the marked space isomorphically
onto its marked image. Its coefficient there is the inverse tame
character of the determinant of this map. A block triangular
change of frame preserves this support and multiplies the
coefficient by the character of the corresponding diagonal
block. The inverse frame convention replaces the character by its
inverse. This is the transformation law in the split test of
Section~\ref{nav:VII-67}. An opposite flag gives the same result by
changing the ordered blocks or taking the contragredient. In either
convention \(\zeta_w\) is the value of the actual character on the
space of forms. Those spaces are defined integrally by the
equations \(j f=\lambda(j)f\). Hecke operators and changes of level
at disjoint places commute with these equations. The other local
tests and the integration in the source variable preserve this
target transformation law.

Every chosen power ramifies at each specified twist place by the
construction of Theorem~\ref{thm:commoncritical731}. Apply
Theorem~\ref{thm:twistblocks738}. The place is split and is away
from \(\ell\). Theorem~\ref{thm:splitattachment722} therefore
transfers the parameter with zero monodromy to the attached
Galois representation. Frobenius semisimplification preserves
its action on inertia, and zero monodromy removes the unipotent
tame factor. The normalizer of
Corollary~\ref{cor:normalizedtransfer722} is unramified at the
new places because its defining source character is unramified there
and the norm power is trivial on local units. Polarization gives
the dual action at the conjugate place up to a cyclotomic power
which is also unramified there. This proves the polynomial
identity. It holds on every stable lattice, before reduction to
any finite precision, because that lattice injects into its
fraction vector space.

The row multiplier from this place is \(o_w(1)^{t_\alpha}\) by
Theorem~\ref{thm:charrows719}. Let
\(D=\max_\alpha(t_\alpha+\min(a_\alpha,b_\alpha))\), taking the
conclusion as empty when there are no types. The proposition
bounds the sum inside the supremum by
\(D\sum_{w\in Q\cup Q^c}v_\ell(1-\zeta_w)\). A character and
its inverse give the same valuation. For each \(\ell\), request
the per-side tame tolerance \(\eta_\ell>0\) in
Theorem~\ref{thm:criticalheights734}, with
\(2D\sum_\ell\eta_\ell\log\ell\le\delta\).
For any embedding and each signed power the last sum over
places is at most \(2\eta_\ell\), counting both members of every
pair. The tolerance can be imposed together with all the earlier
critical-value and Selmer tests, since it is chosen before the
auxiliary primes and roots of unity. This proves the bound
uniformly in all later trace levels. On a finite interval one can
use the largest of the finitely many values of \(D\) and the same
tolerances at each node.
\end{proof}

These local terms concern the primes introduced by the twists.
The local equations above the coefficient prime and the
comparison of the full projected period in Section~\ref{sec:missing}
require separate estimates.
\clearpage

\subsection{The projected torus pairing at every coefficient prime}
\label{sec:projectedcups742}

We compare the integral pairing of two projected tests with the
free cup products of Corollary~\ref{cor:torustests741}. Both
quantities will be evaluated on the given period field, for the
same twist, at the precision used to separate its Hecke system.
This gives an estimate on the entire prime sum of the difference.

Use the shorter intervals and the main torus pairs of
Theorem~\ref{thm:negativeinterval740}. At a node \(x\) put
\(N=q_1d\), \(t_0=5r\), \(Y=N+t_0+1\), and use the integer
\(G_0\) of Corollary~\ref{cor:actualmarks714}. Let
\[
        Z_x=h_E^2\mathcal J_{N,t_0}\,
                       G_0^{\,6+(4t_0+3)Y^6}.
\]
Here \(h_E\) is the class number of \(E=E_x\). The integer
\(m_x<z\le r\) of a designated main pair gives target rank
\(2z\) and string length \(t'=2z-1\le t_0\).
At a coefficient prime \(\ell\), use one of the simultaneous
twists for the ratios and Selmer groups from
Theorems~\ref{thm:criticalheights734} and~\ref{thm:rankoneselmer737}.
For a pair of integral tested forms \(\Xi,\Xi'\) of
Section~\ref{sec:rankonepair736}, let \(P\) project onto the
line of the chosen torus lift. Set \(P=0\) for a zero lift.
Measure its denominator by
\[
 p(D)=\frac{\log\max\{1,|\iota\langle P\Xi,P\Xi'\rangle|_\ell\}}
                  {\log\ell}.
\]
The datum \(D\) includes the node, prime, coefficient embedding,
main pair, twist power, and two tests. For a nonzero lift it also
includes any choice of a separating trace family and precision
\(\ell^\nu\) as in Corollary~\ref{cor:baseseparation723}.

For this family let \(T\) denote the twisted main line and put
\(T_e=\mathcal O(-e)\) for \(0\le e<t'\). Let \(C_T\) be the
image in \(\bigoplus_{w\mid\ell}H^1(E_w,T/\ell^\nu T)\)
of the normalized global main trace cocycles killed by
\(\ell^\nu\) in Section~\ref{sec:relationmatrices719}. Identify
\(T/\ell^\nu T\) with the \(\ell^\nu\)-torsion of
\(T[1/\ell]/T\) by division by \(\ell^\nu\).
Let \(C_e\) be the local image for the scalar cocycles in the
same section, with the same trace relations and precision.
In each case apply the joint map of Proposition~\ref{prop:freetests741}
at \(a=\nu e_L\). Denote it by \(P_a^M\) for lattice \(M\), and
define the finite number
\[
 r(D)=\frac{1}{e_L}\left(
    \operatorname{length}_{\mathcal O}P_a^T(C_T)
       +\sum_{e=0}^{t'-1}
         \operatorname{length}_{\mathcal O}P_a^{T_e}(C_e)\right).
\]
The identification of \(\ell^\nu\mathcal O\) with
\(\pi^a\mathcal O\) changes coordinates by a unit at most. The
choices of free generators may be made separately for each
lattice. Put \(r(D)=0\) for a zero lift.

\begin{theorem}[The prime sum for the projected integral pairing]
\leavevmode\label{thm:projectedcups742}
Fix \(q_0,d\) and make one choice of the shorter period data at
each integer node in an interval starting at \(V_0\). The twists
can be chosen simultaneously with the bounds for the critical
ratios so that
\[
 \sum_\ell(\log\ell)\sup_{D\ {\rm at}\ \ell}
          \frac{\max\{0,p(D)-r(D)\}}{\sqrt{b_x}}
                    \ll_{q_0,d} V_0^{-10}.
 \tag{10}\label{eq:projectedcups742}
\]
The supremum includes every node of the interval, all the
indicated integral tests, coefficient embeddings, signed twist
powers, and separating families. At a nonzero lift each \(r(D)\)
uses only \([E:\QQ]/2\) free main coordinates and
\([E:\QQ]\) free coordinates for each of the \(t'\) scalar
summands. In the common labelled lifting range of
Theorem~\ref{thm:rankonespherical737}, \(r(D)=0\).
\end{theorem}
\begin{proof}
Fix all the data at a nonzero lift. Use positive tolerances
\(\epsilon_{x,\ell},\eta_{x,\ell}\) for the normalized Selmer
length and the sum of the new inertia costs respectively. The
simultaneous construction allows both to be decreased, for the
same list of signed critical ratios. Write
\[
   B_{x,\ell}=a_\ell(t')+\epsilon_{x,\ell},
\]
using the explicit odd cyclotomic ideal of
Proposition~\ref{prop:oddcyclo737}. The global main group
satisfying the integral local equations has normalized length
at most \(B_{x,\ell}\) by Theorem~\ref{thm:rankoneselmer737}.

We verify that the single finite precision in \(r(D)\) suffices.
The row lattice in the proof of Proposition~\ref{prop:blockrows719}
contains \(\ell^\nu\mathcal O^{kr'}\), where \(k\) is its number of
columns before reading the blocks and \(r'\) is the main rank
(here \(r'=1\)). Its annihilator therefore maps into global
normalized raw trace cocycles killed by \(\ell^\nu\). Reading
their values at the row generators recovers the annihilator
vectors. The determinant index in that proof is consequently
bounded by the length of this finite module of cocycles.
The kernel of restriction to the local quotients by
\(F_a(T)\) satisfies the integral evaluation equations, by
Lemma~\ref{lem:localevaluation728} and division by \(\ell^\nu\).
These cocycles vanish at the chosen trace inertia generators,
so are unramified at the later trace places as in the proof of
Corollary~\ref{cor:maincost730}. At the marked normalization
element \(\eta\), \(T(\eta)-1\) is a unit. A coboundary with
values in the divisible module that vanishes at \(\eta\) must
therefore be zero. Thus the kernel injects into the global
group just bounded. Its normalized length is at most
\(B_{x,\ell}\). This argument does not require crystallinity
of an auxiliary matrix at the coefficient prime.

On the finite image \(C_T\), the index of the local conditions
is at most the normalized length of its free cup-product image
plus \(\sum_{w\mid\ell}h_w(T^\vee(1))\), by
Proposition~\ref{prop:freetests741}. For \(T_e\) the colength
\(l_e\) in Theorem~\ref{thm:ribettrace720} is the index in
the scalar trace module of those same coefficient-place
evaluation equations for \(T_e\). Applying the proposition
to \(C_e\) bounds \(l_e/e_L\) by its free image length divided
by \(e_L\) and the torsion error for \(\mathcal O(e+1)\).
Put
\[
 h(D)=\sum_{w\mid\ell}
       \left(h_w(T^\vee(1))+\sum_{e=0}^{t'-1}h_w(\mathcal O(e+1))\right).
\]
In total, the normalized contribution of the two trace terms
to the coefficient of the separating multiplier is at most
\(B_{x,\ell}+r(D)+h(D)\).

The other rows of this multiplier use only places away from
the coefficient prime. At each old level different from
\(\ell\) the Frobenius rows of Corollary~\ref{cor:levelcost725}
cost in total at most \(2v_\ell(G_0)\). An old level equal to
\(\ell\) is already included in \(C_T,C_e\). At the new twist
pairs the parahoric line character is nontrivial for every
selected power. Its inertia polynomial, with the other block
spherical, is the one in Theorem~\ref{thm:tametransfer721}.
The two sides of all pairs cost at most \(2t_0\eta_{x,\ell}\)
as in Theorem~\ref{thm:rankonespherical737}. At the subsidiary
split trace places the determinant unit and normalized Hecke
rows follow from Corollaries~\ref{cor:markeddet715}
and~\ref{cor:normalizedtransfer722}. They are chosen away from
\(\ell\). At other places away from \(\ell\) the spherical
representation is unramified. Thus these row calculations are
valid for every coefficient prime. The remaining terms of
Theorem~\ref{thm:ribettrace720} are bounded after normalization
by \(v_\ell(h_E\mathcal J_{N,t_0})+(4t_0+3)Y^6v_\ell(G_0)\),
using \(t'\le t_0\) and Corollary~\ref{cor:actualmarks714}.

Apply Theorem~\ref{thm:rankonepair736} to the original two
tests. Its averaging and stabilizer factor \(Hc_0^2|C|\)
divides \(h_E G_0^4\). Indeed \(|C|\mid h_E H\) by
Proposition~\ref{prop:torusorder736}, the residue unit orders
give \(H\mid G_0\), and the stabilizer multiple divides \(G_0\)
by the block inclusions \(2z\le r t_0\). Combining the bounds
gives the pointwise estimate
\[
   p(D)\le r(D)+h(D)+a_\ell(t')+v_\ell(Z_x)
                          +\epsilon_{x,\ell}+2t_0\eta_{x,\ell}.
 \tag{11}\label{eq:periodpoint742}
\]
At a zero lift the estimate holds with the torsion terms omitted.
The precision required in base separation was chosen for this
base family, before its subsidiary levels were introduced.
There is no factor of \(\nu\) in any term on the right other
than the actual image length \(r(D)\). The count of coordinates
is Corollary~\ref{cor:torustests741}. In the labelled lifting
range all restricted cocycles are reductions of local finite
cocycles by Corollary~\ref{cor:labelledtrace735}. They annihilate
the free tests, yielding \(r(D)=0\).

It remains to sum terms that vary with the node.
Corollary~\ref{cor:torustests741} bounds the weighted sum of
the suprema of \(h(D)/\sqrt{b_x}\) by \(O_{q_0,d}(V_0^{-10})\).
At a fixed node the prime sum of \(\sup_{t'}a_\ell(t')\)
with weights \(\log\ell\) is at most
\(\log(N\operatorname{lcm}(1,\ldots,t_0))\), by
Proposition~\ref{prop:oddcyclo737}. Prime factorization gives
exactly \(\log Z_x\) for the corresponding \(Z_x\) term.
We have \(N,[E:\QQ]\ll_{q_0,d}x^{12}\),
\(r,t_0\ll x^5\), and \(\log(p_0p_1)\le22\log b_x\).
The formula for the orders of the finite general linear groups
gives \(\log G_0\le [E^+:\QQ](r t_0)^2\log(p_0p_1)\).
The class number estimate and the bound
\(\log\mathcal J_{N,t_0}\ll N^4t_0^2\log(2Nt_0)\) in
Section~\ref{sec:oddcocycles675} and
Theorem~\ref{thm:oddinteger677} consequently give
\[
 \log Z_x+\log\!\bigl(N\operatorname{lcm}(1,\ldots,t_0)\bigr)
                 \ll_{q_0,d} x^{110}(1+\log b_x).
\]
For the tolerance terms choose
\(\epsilon_{x,\ell}+2t_0\eta_{x,\ell}
 \le 2^{-\ell}/(1+\log\ell)\), decreasing them further for
any prescribed critical-ratio tolerances. Their weighted sum
at one node is at most one. Use \(1+\log b_x\ll b_x^{1/8}\)
and \(b_x\ge x^{500}\). After division by \(\sqrt{b_x}\)
all the fixed and tolerance terms at a node have prime sum
\(O_{q_0,d}(x^{-77})\). Since all terms are nonnegative, their
supremum across nodes is at most their sum across those nodes
before taking the prime sum. The tail over integers \(x\ge V_0\)
is \(O_{q_0,d}(V_0^{-76})\). This proves
\eqref{eq:projectedcups742}.
\end{proof}

Inequality~\eqref{eq:projectedcups742} controls the difference
from a specified image length for the projected torus tests.
A bound for \(r(D)\) on the remaining coefficient places and
the comparison for the higher-rank chain in
Section~\ref{sec:missing} are still required for the conclusion
about zeros.
\clearpage

\subsection{The contribution of abelian local systems to the torus pairing}
\label{sec:abelianperiod744}

Fix the shorter intervals, the rank-one marked separating systems,
and the integral tests of Theorems~\ref{thm:projectedcups742}
and~\ref{thm:oldtorustraces743}. Write \(K=E_{x,w}\) for a
completion at a coefficient prime \(\ell\). We call a member
\(u_y\) of the trace family abelian at \(w\) if the matrices of
\(u_y(G_K)\) commute. This includes the normalized representation
of the selected lift itself. We shall evaluate the remaining
local conditions on explicit matrices whose components at every
such member are zero. The prime sum of the resulting error is
uniform in the congruence precision and the number of systems.

Use \(a=\nu e_L\), \(e_L=v_\pi(\ell)\), and the notation for
\(C_M,P_a^M\) of Section~\ref{sec:projectedcups742}. For the
matrix tuple \(D(g)=(u_y(g)-1)_y\) set
\(\mathcal N=\sum_g\mathcal O D(g)\). Each raw trace cocycle
\(z\) has an \(\mathcal O\)-linear map \(f_z\) on \(\mathcal N\)
with \(f_z(D(g))=z(g)\), by its defining relations. For a
finite group-ring expression \(X=\sum_g b_g[g]\) of augmentation
zero set \(u(X)=(\sum_g b_g u_y(g))_y\). Then \(u(X)\in\mathcal N\)
and \(f_z(u(X))=\sum_g b_gz(g)\). This uses only integral
relations among the actual matrices. Coefficients can be extended
to contain all local embeddings while retaining the ideal
\(\ell^\nu\) and extending the image modules by flatness.
Normalized lengths are invariant under this operation.

\begin{lemma}[Unit characters visible in an abelian trace system]
\leavevmode\label{lem:abelianunits744}
At a spherical coefficient place, the inertia action of an
abelian member is a sum of characters which, on
\(\mathcal O_K^\times\) by local reciprocity, have the form
\[
        v\longmapsto\prod_{\tau:K\hookrightarrow L}
                             \tau(v)^{\epsilon b_\tau},
   \qquad
   b_\tau\in\{0,\ldots,t'-1,\ z-1+\epsilon_\tau k_\tau\}.
 \tag{12}\label{eq:unitweights744}
\]
Here \(\epsilon=\pm1\) is the reciprocity convention and
\(\epsilon_\tau\) records the two CM orientations. The coefficients
may be enlarged. At an old coefficient place the formula has one
additional smooth unit character of order dividing \(Q\).

There is an element \(H_w\in\mathcal O[G_K]\), of nonzero
augmentation \(\delta_w\), such that
\[
        (u_y(g)-1)\,u_y(H_w)=0
             \qquad(g\in I_K,\ u_y\ {\rm abelian\ at}\ w).
 \tag{13}\label{eq:abelianann744}
\]
At both old primes \(\delta_w\) is a unit. On each node,
putting \(n_x=[E_x:\QQ]\), one can choose the elements so that
\[
 \sum_\ell(\log\ell)
       \sup_D\sum_{w\mid\ell}[E_{x,w}:\QQ_\ell]\,v_\ell(\delta_w)
            \ll_{q_0,d} x^{60} b_x^{1/4}.
 \tag{14}\label{eq:unitbound744}
\]
The supremum is over the choices at that prime and node. The
bound includes places where \(E_{x,w}/\QQ_\ell\) is ramified.
\end{lemma}
\begin{proof}
All normalized trace systems of this weight have the same
labelled Hodge multiset as the selected nonzero base lift:
the weight formula depends only on the infinity type, and the
base representation is the sum of the main character and the
scalar string (Theorem~\ref{thm:rankoneselmer737}). They are
crystalline at the spherical coefficient places. Indeed the
spherical base change has an unramified parameter with zero
monodromy, and the compatibility argument in
Corollary~\ref{cor:scalarfl724} gives crystallinity before any
condition on the length of the interval of weights. This also
applies when \(K\) itself is ramified over \(\QQ_\ell\).
The smooth normalizer is unramified here, the twist splits above
\(\ell\), and the subsidiary levels avoid \(\ell\). At the
original levels Theorem~\ref{thm:oldtorustraces743} gives one
crystallizing extension of tame degree \(Q\).

For commuting matrices extend coefficients and use the generalized
eigenspaces of the finite-dimensional commutative algebra they
generate. On each such space divide by the eigencharacter
\(\psi\). The quotient representation has a filtration with
trivial successive quotients and is crystalline over the same
crystallizing field as the original representation and \(\psi\).
To see that it is unramified there, argue by induction along the
filtration. For an unramified weight-zero representation \(V\),
unramified extension classes lie in \(H^1_f(K',V)\), since
unramified representations are crystalline. Their dimension is
\(\dim V/(\mathrm{Fr}-1)V=\dim H^0(K',V)\), the dimension
of \(H^1_f(K',V)\) by the Bloch--Kato formula. Thus an extension
of the trivial line by such a \(V\) which is crystalline is
unramified. This proves the induction. Over \(K\) the inertia
in this unipotent part is now finite, hence trivial in
characteristic zero. Thus the inertia action is semisimple even
when the abelian representation has a nontrivial unramified
extension. For a crystalline character, local reciprocity on
units gives precisely the algebraic monomial with its labelled
weights. For a potentially crystalline character divide by
that crystalline algebraic character; the smooth unit factor
is then the finite inertia type. This proves
\eqref{eq:unitweights744} with the asserted modification.
~\cite[Appendix~B, Proposition~B.4]{X-VII-conradunits744}

At an old prime put \(p=\ell\). Choose \(s\in I_K\) lifting
the reciprocity image of a Teichm\"uller generator of
\(\mathbb F_p^\times\), and take \(H_w=\sum_{j=0}^{p-2}[s]^j\).
Every nontrivial unit character in question takes a nonidentity
\((p-1)\)-st root of unity at \(s\). Indeed if its labelled
exponent \(b\) is nonzero then \(0<Q|b|<p-1\), so raising the
value to \(Q\) is already nontrivial. If \(b=0\), a nontrivial
smooth character of the cyclic residue group is nontrivial on
its generator. Therefore \(H_w\) kills every nontrivial inertia
character. On the remaining spaces inertia acts trivially.
This gives \eqref{eq:abelianann744}, with \(\delta_w=p-1\).

At the other primes we first describe how to detect all nontrivial
characters in \eqref{eq:unitweights744}. Call the choice
\(z-1+\epsilon_\tau k_\tau\) at a label the main jump. Every
\(k_\tau\) exceeds \(3n_xt_0\) for sufficiently large \(x\):
the formula in Theorem~\ref{thm:negativeinterval740} gives
\(k_\tau\ge3r_x^8-4r_x\), while \(n_x\ll_{q_0,d}x^{12}\),
\(x^5\le r_x\le3x^5\), and \(t_0=5r_x\). If only scalar
jumps occur and some are nonzero, rational units detect their
positive sum. If the quadratic CM field splits at \(\ell\),
all main jumps at \(w\) have the same orientation. One main
jump suffices to make the sum of all selected jumps nonzero,
even when the scalar jumps have the opposite sign. Use rational
units in this case also. At a nonsplit prime, as soon as a main
jump occurs, use norm-one units in the quadratic completion
inside \(K\). The exponent on such a unit, up to one overall
sign and a choice of its embedding, is \(\sum_\tau\epsilon_\tau
b_\tau\). Every selected main jump contributes \(+k_\tau\)
to this sum apart from its term \(\epsilon_\tau(z-1)\).
All the latter terms and the scalar terms together have absolute
value at most \(n_xt_0\). This proves detection, including when
the rational sum is zero.

Here is a common finite list of absolute exponents for those
tests. For a fixed \(1\le j_0\le m_x\) let \(w_1,w_2\) be
the two base weights after adding \((2j_0+1)r_x^8\), as in
Theorem~\ref{thm:alltorsion741}. Include \(1,\ldots,n_xt_0\)
and the positive integers of the form
\[
   |n_1 w_1+n_2 w_2+D_0|,\qquad
   n_1,n_2\ge0,\quad 1\le n_1+n_2\le n_x,\quad
   |D_0|\le 3n_xt_0 .
\]
Take the union over \(j_0\) and call it \(\mathcal B_x\).
Indeed the selected main magnitudes are \(w_i-h_\tau\) with
\(0\le h_\tau\le4r_x\); all the discarded scalar and central
shifts fit in the indicated bound. This list has size
\(O((m_x+1)(n_x+1)^2(1+n_xt_0))\) and maximum
\(O(n_xb_x)\).

For each \(B\in\mathcal B_x\) and each required direction
choose one unit \(v_B\). For the rational direction use a
Teichm\"uller generator if \(\ell-1\nmid B\), and
\(1+\ell^{c_\ell}\) otherwise, where \(c_2=2\) and
\(c_\ell=1\) for odd \(\ell\). In the inert norm-one direction
use a Teichm\"uller generator of order \(\ell+1\) when
\(\ell+1\nmid B\), and the norm-one unit of depth \(c_\ell\)
in the proof of Theorem~\ref{thm:alltorsion741} otherwise.
At the ramified quadratic prime \(d\) use the norm-one unit
of depth \(3/2\) of that proof. Lift each reciprocity value
to some \(s_B\in I_K\). For each direction and each sign take
the factor \([s_B]-v_B^{\pm B}[1]\), reading \(v_B\) in the
chosen embedding of its rational or quadratic field. Define
\(H_w\) to be the product of these factors in any fixed order.
Every nontrivial inertia character in the abelian member vanishes
on at least one factor, in the sense that its linear extension to
the group ring takes that factor to zero. All factors preserve
the inertia character spaces of this member and act on them by
scalars. This proves \eqref{eq:abelianann744} over the fraction
field and hence on the original integral lattices.

Each factor has nonzero augmentation. For a fixed \(B\), its
valuation is zero on the residue tests. On the principal-unit
tests it is \(c_\ell+v_\ell(B)\), or \(3/2+v_d(B)\) for the
ramified norm-one unit, by the binomial evaluations in
Theorem~\ref{thm:alltorsion741}. Summed over primes, both
directions and both signs with weights \(\log\ell\), this costs
\(O_d(\tau(B)\log(B+1)+\log(2B))=O_d((1+B)^{1/4})\).
Indeed the rational and inert tests with positive cost require
\(\ell-1\mid B\) and \(\ell+1\mid B\), respectively, and the
ramified prime is fixed. Sum over \(\mathcal B_x\), then use
\(\sum_{w\mid\ell}[E_{x,w}:\QQ_\ell]=n_x\). The explicit bound
is \(O_d(n_x|\mathcal B_x|(1+n_xb_x)^{1/4})\). Inserting
\(m_x\le2\sqrt{x}\) and the bounds for \(n_x,t_0\) proves
\eqref{eq:unitbound744}.
\end{proof}

\begin{theorem}[Projected pairings and the nonabelian local matrices]
\leavevmode\label{thm:abelianperiod744}
For the indicated marked separating families at each prime and
node there is a number \(r_*(D)\) computable from their integral
trace functionals on tuples with zero components at every member
abelian at the place of evaluation. It satisfies
\[
  0\le r_*(D)\le r(D),\qquad
  \sum_\ell(\log\ell)\sup_{D\ {\rm at}\ \ell}
          \frac{(p(D)-r_*(D))_+}{\sqrt{b_x}}
                      \ll_{q_0,d} V_0^{-10}.
 \tag{15}\label{eq:abelianperiod744}
\]
In particular the whole normalized prime sum of the projected
pairing on the subfamily in which every \(u_y\) is abelian at
every \(w\mid\ell\) is \(O_{q_0,d}(V_0^{-10})\). This conclusion
holds for the same twists and integral tests used to bound the
critical ratios.
\end{theorem}
\begin{proof}
First take \(M=T\) or \(M=\mathcal O(-h)\) with \(1\le h<t'\),
and write \(\alpha\) for its character. At each \(w\) choose
\(s_w\in I_K\) minimizing the valuation of
\(\lambda_{M,w}=\alpha(s_w)-1\). This minimum is finite by the
unit tests of Theorem~\ref{thm:alltorsion741} for \(T\), and by
the rational-unit test with exponent \(h[K:\QQ_\ell]\) for the
scalar. For \([s,g]=sgs^{-1}g^{-1}\), the cocycle rule gives
\[
 f_z\bigl(D([s_w,g])\bigr)=
     (\alpha(s_w)-1)z(g)-(\alpha(g)-1)z(s_w)
                                  \qquad(g\in G_K).
\]
As a function of \(g\) this represents \(\lambda_{M,w}[z]\).
The second term is a coboundary and all abelian matrix components
of \(D([s_w,g])\) vanish.

For \(M=\mathcal O\) take \(\lambda_{M,w}=\delta_w\) of
Lemma~\ref{lem:abelianunits744}. A cocycle is now an additive
homomorphism. For every \(g\in I_K\),
\[
       f_z\bigl(u(([g]-[1])H_w)\bigr)=\delta_w z(g).
\]
Indeed write \(H_w=\sum b_v[v]\) and use
\(z(gv)-z(v)=z(g)\). The tuple inside \(f_z\) is in
\(\mathcal N\) and has zero abelian components by
\eqref{eq:abelianann744}. The indicated inertia values determine
all the free cup products of \(\delta_w[z]\). Two extensions of
these values to \(G_K\) differ by an unramified homomorphism,
and every unramified homomorphism to \(\mathcal O/\pi^a\)
lifts integrally and belongs to the reduced finite condition.
Proposition~\ref{prop:freetests741} shows its free cups vanish.

In the definition of \(r(D)\), multiply all free coordinates at
\(w\) for \(M\) by \(\lambda_{M,w}\) and take the normalized
length of each resulting joint image. Their sum is \(r_*(D)\);
take zero at a zero lift. The descriptions above specify exactly
which evaluations on \(\mathcal N\) compute these images. In
particular they do not require a functional on \(\mathcal N\)
to extend to the full matrix lattice. On \(d_{M,w}\) copies of
\(\mathcal O/\pi^a\), the kernel of scalar multiplication has
length \(d_{M,w}\min(a,v_\pi(\lambda_{M,w}))\).
Restricting to the joint image can only decrease the kernel
length, including when classes at different places are dependent.
As \(d_{M,w}\le[K:\QQ_\ell]\), it follows that
\[
   0\le r(D)-r_*(D)
       \le\sum_{M,w}[E_{x,w}:\QQ_\ell]\,v_\ell(\lambda_{M,w}).
\]
For the main line, the proof of Theorem~\ref{thm:alltorsion741}
bounds valuations by actual inertia evaluations at all its places.
Its fixed-node prime sum before division by \(\sqrt{b_x}\) is
\(O_{q_0,d}(x^{62}b_x^{1/4})\), and insertion of the local degree
costs at most \(n_x\). For positive scalar weights use the proof
of Corollary~\ref{cor:torustests741} with \(h[K:\QQ_\ell]\)
in place of \((h+1)[K:\QQ_\ell]\). Both exponents lie between
one and \(n_xt_0\), so the same factorial estimate gives
\(O(n_x(n_xt_0+1)^2\log(2+n_xt_0))\) after inserting the
local degree. The weight-zero cost is \eqref{eq:unitbound744}.
Together these give \(O_{q_0,d}(x^{80}b_x^{1/4})\) at the node.
After division by \(\sqrt{b_x}\) this is \(O_{q_0,d}(x^{-45})\)
because \(b_x\ge x^{500}\). Bound the supremum across integer
nodes by their sum and apply Theorem~\ref{thm:projectedcups742}.
This proves \eqref{eq:abelianperiod744}. If all the indicated
systems have abelian image locally, every displayed matrix
evaluation vanishes and hence \(r_*(D)=0\).
\end{proof}

The selected torus representation has abelian image, but other
systems in a separating family can have noncommuting local
matrices. Their integral evaluations in
\eqref{eq:abelianperiod744} and the higher-rank projected period
in Section~\ref{sec:missing} still require bounds.
\clearpage

\subsection{Two inertia elements at the original level primes}
\label{sec:oldwild745}
At either original level prime, Theorem~\ref{thm:oldtorustraces743}
leaves several scalar weights at the same completions. We give
one simultaneous choice of inertia elements for all these weights,
at every precision. This identifies the remaining joint image inside
the actual integral trace module.

Write \(G=G_{\QQ_p}\), let \(I\) be inertia, let \(P\) be wild
inertia and let \(\kappa\) denote the \(p\)-adic cyclotomic
character. For the integers \(\mathcal O\) of a finite extension
of \(\QQ_p\), put \(R_a=\mathcal O/\pi^a\), \(a\ge1\). All
cohomology below is continuous. A choice of Tate basis identifies
the underlying module of \(R_a(-h)\) with \(R_a\).

\begin{lemma}[Simultaneous evaluation of the scalar ramification classes]
\leavevmode\label{lem:wildgenerators745}
For each \(p\ge5\) there are \(s\in I\) of pro-prime-to-\(p\)
order and \(g\in P\) with \(\kappa(s)\) a primitive
\((p-1)\)-st root and \(\kappa(g)=1+p\). Put \(b=[s,g]\).
For every \(\mathcal O,a\) as above one has isomorphisms
\[
 \begin{aligned}
 H^1(G,R_a(-h))&\longrightarrow R_a,& [c]&\longmapsto c(b)
                          &&(1\le h\le p-3),\\
 H^1(G,R_a)/H^1_{\rm unr}(G,R_a)&\longrightarrow R_a,
                & [c]&\longmapsto c(g).&&
 \end{aligned}
\]
The same \(s,g\) work for all these maps.
\end{lemma}
\begin{proof}
Lift a topological generator of \(I/P\) to inertia. Its procyclic
closure is the product of its pro-\(p\) part and its
prime-to-\(p\) part. The prime-to-\(p\) component \(s\) has
the same image in \(I/P\). Write \(u=\overline{\kappa(s)}\).
It generates \(\mathbb F_p^\times\) by total ramification of
\(\QQ_p(\mu_p)\), and \(\kappa(s)\) is its Teichm\"uller lift.

For \(1\le h\le p-3\), both \(\bar\kappa^{-h}\) and
\(\bar\kappa^{h+1}\) are nontrivial on inertia. Local Tate
duality and the local Euler characteristic formula give
\(\dim_{\mathbb F_p}H^1(G,\mathbb F_p(-h))=1\).
Choose a cocycle \(c_h\) representing a nonzero class.
Put \(J_0=I\cap\ker\kappa\) and \(J=P\cap\ker\kappa\). On
\(J_0\) the cocycle is a homomorphism. If it vanished there,
it would factor through \(G/J_0\), which embeds in
\((G/I)\times\mathbb Z_p^\times\) and is abelian. Commuting
the images of \(s\) and an arbitrary \(v\) would give
\[
 (u^{-h}-1)c_h(v)=(\bar\kappa(v)^{-h}-1)c_h(s).
\]
The first scalar is a unit, so \(c_h\) would be a coboundary.
Its restriction to \(J\) is also nonzero, since \(J_0/J\)
embeds in tame inertia and has no quotient of order \(p\).

The homomorphisms \(\psi_h=c_h|_J\) are linearly independent
over \(\mathbb F_p\). Indeed \(j\mapsto sjs^{-1}\) acts on their
arguments with distinct eigenvalues \(u^{-h}\), by the cocycle
rule. Thus \(j\mapsto(\psi_h(j))_{1\le h\le p-3}\) is surjective
because its image is a vector subspace and any proper subspace has a
nonzero annihilating linear functional. The image of \(P\) under
\(\kappa\) is \(1+p\mathbb Z_p\). This follows from total
ramification of the cyclotomic tower and the prime-to-\(p\)
order of \(I/P\). Lift \(1+p\) to \(g_0\in P\), and choose
\(j\in J\) with \(\psi_h(j)=1-c_h(g_0)\) for all \(h\).
Then \(g=g_0j\) has \(c_h(g)=1\). Since \(b=[s,g]\) belongs
to \(\ker\kappa\), evaluation there kills coboundaries. Direct
substitution in the cocycle identity gives \(c_h(b)=u^{-h}-1\).

For passage to all precisions, \(H^0(G,(\QQ_p/\mathbb Z_p)(-h))=0\)
by the tame unit \(u^{-h}-1\), so \(H^1(G,\mathbb Z_p(-h))\)
is torsion-free. Its rank is one by the rational local Euler
characteristic. The group \(H^2(G,\mathbb Z_p(-h))\) vanishes
by integral local duality since
\(H^0(G,(\QQ_p/\mathbb Z_p)(h+1))=0\). Thus reduction
of integral classes gives all classes modulo \(p\).
Evaluation at \(b\) is an isomorphism on the rank-one integral
module because it is nonzero modulo \(p\). Tensoring with the
finite free \(\mathbb Z_p\)-module \(\mathcal O\) commutes with
continuous cohomology here, as is seen on a \(\mathbb Z_p\)-basis.
The exact sequence for multiplication by \(\pi^a\) and the
vanishing of \(H^2\) give the claimed evaluation for \(R_a\).
For \(h=0\), local reciprocity identifies ramified additive
characters modulo unramified ones with
\(\operatorname{Hom}_{\rm cont}(1+p\mathbb Z_p,R_a)\).
On units the cyclotomic character agrees with the reciprocity
coordinate up to inversion. The image of \(g\) generates that
procyclic group, giving the last isomorphism.
~\cite[Chapter~I, \S2]{X-VII-localduality745}
\end{proof}

\begin{proposition}[A column of integral tame blocks]
\leavevmode\label{prop:tameblocks745}
Let \((u_y)_y\) be a finite family of continuous integral
representations of \(G\). Write \(u= (u_y)_y\) also for the
linear extension to the group algebra, and put
\(\mathcal N=\sum_{v\in G}\mathcal O\,(u(v)-1)\).
Suppose a cocycle \(c:G\to R_a(-h)\) admits an
\(\mathcal O\)-linear functional \(f_c\) on \(\mathcal N\)
with \(f_c(u(v)-1)=c(v)\). Set \(\eta_h=\kappa(s)^{-h}\),
and let \(P_{y,\eta}\) be the spectral projector of \(u_y(s)\)
at \(\eta\), taking zero for an absent eigenvalue. For the
pair of the lemma define tuples
\[
 Z_h=\big(P_{y,\eta_h}u_y(g)P_{y,1}\big)_y\quad(1\le h\le p-3),
 \qquad
 Z_0=\big(P_{y,1}(u_y(g)-1)P_{y,1}\big)_y .
\]
These are integral elements of \(\mathcal N\). For the functional
of the indicated weight they satisfy
\[
       (\eta_h-1)f_c(Z_h)=c([s,g])\quad(h>0),\qquad
                  f_c(Z_0)=c(g).
\]
\end{proposition}
\begin{proof}
The closure generated by \(u_y(s)\) is pro-prime-to-\(p\).
Reduction modulo \(\pi\) has pro-\(p\) kernel, so this closure
injects into a finite residue matrix group. Take \(m'\) divisible
by the orders of all \(u_y(s)\) and by \(p-1\), and prime to \(p\).
For each required \(\eta\), the element
\[
             e_\eta=(m')^{-1}\sum_{i=0}^{m'-1}\eta^{-i}[s]^i
                         \ \in\mathcal O[G]
\]
evaluates in \(u_y\) to \(P_{y,\eta}\), by summation of roots
of unity. This formula works over \(\mathcal O\) even when
other eigenvalues require a coefficient extension. Its augmentation
is one for \(\eta=1\) and zero otherwise. Thus
\(e_{\eta_h}[g]e_1\) for \(h>0\), and
\(e_1([g]-[1])e_1\) for \(h=0\), have zero augmentation.
Every element of zero augmentation is an integral sum of
\([v]-[1]\). This proves membership in the specified lattice.

Write \(\epsilon\) for augmentation, \(\alpha\) for the linear
extension of \(\kappa^{-h}\), and \(d_c([v])=c(v)\), extended
linearly. On products the cocycle identity gives
\[
       d_c(XY)=\alpha(X)d_c(Y)+d_c(X)\epsilon(Y).
\]
On elements of zero augmentation, \(d_c(X)=f_c(u(X))\).
For \(B=[s][g]-[g][s]\) both \(\alpha(B)\) and
\(\epsilon(B)\) vanish, and \(d_c(B)=c([s,g])\).
Since \(\alpha(e_{\eta_h})=\epsilon(e_1)=1\), we have
\(d_c(e_{\eta_h}Be_1)=d_c(B)\). Evaluation in the matrices
gives \(u(e_{\eta_h}Be_1)=(\eta_h-1)Z_h\).
For \(h=0\), \(\alpha=\epsilon\) and the same product rule
gives \(d_c(e_1([g]-[1])e_1)=c(g)\). This uses \(f_c\)
only on its original integral domain.
\end{proof}

\begin{theorem}[The remaining joint image at the level primes]
\leavevmode\label{thm:oldcolumn745}
Use the marked separating families, twists, nodes and shorter
intervals of Theorems~\ref{thm:projectedcups742}
and~\ref{thm:oldtorustraces743}. At a node \(x\) let \(p\) be
either original level prime. For a nonzero lift let \(W^-\) be
the completions \(w\mid p\) of \(E_x\) with \(J_w<0\).
Transport the pair \(s,g\) for \(\QQ_p\) to each of the chosen
decomposition groups at \(w\). For each \(0\le h<t'-1\) let
\(\Gamma_h(D)\) be the joint image
\[
  c\longmapsto\big(f_c(Z_{h,w})\big)_{w\in W^-}
                \quad\text{in }\bigoplus_{w\in W^-} R_a
\]
of the original normalized global scalar trace cocycles of weight
\(h\). Here \(Z_{h,w}\) is formed from the restricted family,
and its local trace lattice maps into the global one.
At the original trace precision \(a=e_L\nu\) there is an equality
\[
       r(D)=B_{\rm old}(D),\qquad
       B_{\rm old}(D)=\frac1{e_L}\sum_{0\le h<t'-1}
                             \operatorname{length}_{\mathcal O}\Gamma_h(D).
\]
Take \(B_{\rm old}=0\) at a zero lift. In particular
\[
  \sum_p(\log p)\!
  \sup_{\substack{D\ {\rm at}\ p\\p\ {\rm an\ original\ level\ at}\ x}}
       \frac{(p(D)-B_{\rm old}(D))_+}{\sqrt{b_x}}
          \ll_{q_0,d} V_0^{-10},
\]
with the empty supremum set to zero.
\end{theorem}
\begin{proof}
All the completions under consideration equal \(\QQ_p\),
and \(t'<p-2\) by \(t'\le t_0\le b_x\), \(p\ge b_x^{10}\).
At \(h>0\) the rational finite local condition on
\(\QQ_p(-h)\) is zero by the crystalline dimension formula
(its de Rham tangent space and its invariants vanish).
Its integral reduced finite condition is zero as well, since
the integral first cohomology of this character is torsion-free.
For \(h=0\) the rational finite condition is the unramified line
because unramified extensions are crystalline and the crystalline
dimension formula gives dimension one. The integral condition and
its reduction consist of all unramified classes by local reciprocity.
For every \(0\le h<t'\) the dual \(H^1_f(\QQ_p,\mathcal O(h+1))\)
has rank one by the same dimension formula. It has no torsion
since \(\bar\kappa^{h+1}\) is nontrivial on inertia. The free
cup test of Proposition~\ref{prop:freetests741} therefore
has exactly the indicated finite condition as its kernel.
By Lemma~\ref{lem:wildgenerators745} its quotient has length
\(a\), so its injection into one copy of \(R_a\) is an
isomorphism. The cup coordinate and the evaluation of the lemma
differ by a unit. For \(h>0\), \(\eta_h-1\) is itself a unit.
Proposition~\ref{prop:tameblocks745} now identifies each
surviving cup coordinate up to a unit with \(f_c(Z_{h,w})\).

Theorem~\ref{thm:oldtorustraces743} kills the main coordinates
everywhere, every scalar coordinate at the positive completions,
and the top scalar coordinate at the negative ones. At each weight
the comparison on the other coordinates is an automorphism of
their direct sum, hence preserves the length of the joint image
even for dependent restrictions. This proves the equality. The
prime sum follows by restricting the supremum in
\eqref{eq:projectedcups742} to the displayed data.
\end{proof}

There is a further simplification for a member whose residual
characteristic polynomial at \(s\) agrees with that of the base
member. In this case both projectors in \(Z_{h,w}\) have rank one.
Indeed the base eigenvalues are \(\kappa(s)^{-e}\) for
\(0\le e<t'\) and \(\xi\kappa(s)^{-J_w}\) with \(\xi^Q=1\),
by the local character formula of Theorem~\ref{thm:oldtorustraces743}.
The first roots are distinct. Equality of the last with any of
them after reduction would imply \(p-1\mid Q(J_w-e)\).
The jumps and \(Q\) in that theorem give
\(0<Q|J_w-e|<p-1\). Prime-to-\(p\) roots remain distinct after
reduction, so the spectral summands of the specified eigenvalues
have multiplicity one also in the given member. Their integral
images are direct summands and admit bases over \(\mathcal O\).
In these bases each displayed block is a single entry. The
equality for \(B_{\rm old}\) still involves functionals on the
joint trace lattice. A bound for the lengths of their images,
uniform in the precision and the family, is needed to remove
\(B_{\rm old}\) from the period estimate.
\clearpage

\subsection{The projected pairing for potentially abelian local systems}
\label{sec:potentialperiod746}
A local representation is potentially abelian if its image on
some open subgroup is abelian. The subgroup and its index may
vary with the representation. We use the common crystallizing
field from Section~\ref{sec:tameends743} to compare all these
systems on a single inertia group.

\begin{lemma}[Inertia of a potentially abelian crystalline representation]
\leavevmode\label{lem:potinertia746}
Let \(F\) be a finite extension of \(\QQ_\ell\). Let \(V\) be a
potentially abelian crystalline representation of \(G_F\) of
dimension \(R\) over a finite extension of \(\QQ_\ell\).
After a finite extension of coefficients its restriction to \(I_F\)
is a sum of characters. Every such character is the restriction
of a crystalline character on the unramified extension \(F_f/F\)
of some degree \(1\le f\le R\). At each embedding of \(F_f\) its
jump belongs to the Hodge multiset of \(V\) at the underlying
embedding of \(F\).
\end{lemma}
\begin{proof}
Choose a finite Galois extension \(L/F\) on whose absolute
Galois group the matrices commute. Write \(A\) for the
coefficient algebra spanned by those matrices, and \(A'\)
for its centralizer in \(\operatorname{End}(V)\).
Conjugation makes \(A'\) a subrepresentation of the crystalline
representation \(\operatorname{End}(V)\). Its action factors through
\(\operatorname{Gal}(L/F)\), and hence is unramified. For the last
assertion, a finite-image crystalline representation \(U\) trivialized
by \(L\) has crystalline comparison over \(L\) with period space
\(L_0\otimes_{\QQ_\ell} U\), where \(L_0\) is the maximal absolutely
unramified subfield. Its \(F\)-crystalline periods span this space
over \(L_0\) by the crystalline comparison and base change. Relative
inertia fixes \(L_0\) and these periods, so acts trivially on \(U\).
One can forget the coefficient field for this argument.
It follows that the inertia matrices for \(V\) belong to the
double centralizer \(A''\). Since \(A\) is commutative,
\(A\subset A'\) and \(A''\subset A'\). Any two elements of
\(A''\) commute.

Over \(L\), each generalized eigenspace of the commuting matrices,
twisted by the inverse of its eigencharacter, is crystalline with
a filtration by trivial lines. Such a representation is unramified.
Indeed the unramified extension classes of an unramified weight-zero
representation \(W\) have dimension
\(\dim W/(\mathrm{Fr}-1)W=\dim H^0(L,W)\), and they are
crystalline. The crystalline dimension formula gives precisely
this dimension for \(H^1_f(L,W)\). Induction on the filtration
proves the assertion. Each inertia matrix over \(F\) has a positive
power in \(I_L\) which is therefore semisimple. An invertible
matrix in characteristic zero with a semisimple positive power
is semisimple, as one sees on Jordan blocks. These commuting
semisimple matrices are simultaneously diagonalizable after a
finite coefficient extension.

Frobenius permutes their distinct characters. For one character
choose \(f\) to be its orbit length. Its isotypic space is stable
under \(G_{F_f}\), since inertia preserves it and a lift of
\(\mathrm{Fr}^f\) preserves it. An eigenline of that lift in
the space is stable under all of \(G_{F_f}\), by continuity and
the density of the subgroup generated by inertia and the lift.
It is crystalline as a subrepresentation of the restriction of
\(V\). Exactness and strictness for the filtered de Rham functor,
and base change from \(F\) to \(F_f\), give the assertion about
its labelled jumps. The orbit contains at most \(\dim V\)
characters.
\end{proof}

At the spherical coefficient places put \(F=K=E_{x,w}\).
At an old coefficient place put \(K=\QQ_p\) and \(F=K'\),
the totally ramified cyclic extension of degree \(Q\) in
Theorem~\ref{thm:oldtorustraces743}. In both cases \(I^0=I_F\)
is normal in \(G_K\). At a nonzero lift write \(R=2z=t'+1\),
and put \(R_x=t_0+1\), so \(R\le R_x\). Every trace system
under discussion is crystalline over \(F\).

\begin{proposition}[A common inertia test with controlled augmentation]
\leavevmode\label{prop:potunits746}
There is an integral group algebra element \(H_w\) supported on
\(I^0\), with nonzero augmentation \(\delta_w\), such that
\[
       (u_y(v)-1)u_y(H_w)=0\qquad (v\in I^0)
\]
for every member \(u_y\) potentially abelian at \(w\).
Coefficient extension is allowed with the ideal \(\ell^\nu\)
retained. At an old coefficient place \(\delta_w=1\) is possible.
At a fixed node one has
\[
 \sum_\ell(\log\ell)\sup_D
       \sum_{w\mid\ell}[E_{x,w}:\QQ_\ell]\,v_\ell(\delta_w)
              \ll_{q_0,d} x^{80} b_x^{1/4}.
\]
Here \(v_\ell(\ell)=1\) and the choices in \(D\) are those of
Theorem~\ref{thm:projectedcups742}.
\end{proposition}
\begin{proof}
First suppose the place is spherical. On each character of \(I^0\)
in Lemma~\ref{lem:potinertia746}, the units of \(F_f\) act by
the algebraic monomial of its labelled jumps via reciprocity,
with one overall sign convention. This is the classification of
crystalline characters.~\cite[Appendix~B, Proposition~B.4]{X-VII-conradunits744}
There are at most \(N_x=n_x R_x\) labels, where \(n_x=[E_x:\QQ]\).
Each selected jump is either in \([0,t'-1]\) or is
\(z-1+\epsilon_\tau k_\tau\) at the underlying label of \(K\).
If all selected jumps vanish, inertia acts trivially. If only
scalar jumps occur and one is positive, their sum is positive,
and rational units suffice. If a main jump occurs and \(\ell\)
splits in the imaginary quadratic field, all the main signs
agree. The rational sum is nonzero since
\[
 \min_\tau k_\tau\ge3r_x^8-4r_x>3N_xt_0
\]
for all sufficiently large initial nodes. This uses
\(n_x\ll x^{12}\), \(x^5\le r_x\le3x^5\) and \(t_0=5r_x\).
If a main jump occurs at a nonsplit prime, test on the norm-one
units of the quadratic completion. The corresponding signed sum
of jumps has a contribution \(+k_\tau\) from each main choice,
up to total central and scalar terms of absolute value at most
\(N_xt_0\). Hence this exponent is also nonzero.

For clarity, a single list valid at the node consists of the
integers \(1,\ldots,N_xt_0\), together with the nonzero absolute
values of
\[
    i_1w_1+i_2w_2+D_0,\qquad i_1,i_2\ge0,\quad
    1\le i_1+i_2\le N_x,\quad |D_0|\le3N_xt_0 .
\]
Use the two shifted base weights \(w_1,w_2\) for each of the
\(m_x\) allowed indices. Indeed \(k_\tau=w_i-h_\tau\) with
\(0\le h_\tau\le4r_x\), so all the discarded shifts satisfy
the displayed bound. This list \(\mathcal B\) has cardinality
\(O((m_x+1)(1+N_x)^2(1+N_xt_0))\) and maximum \(O(N_x b_x)\).
For each \(1\le f\le R\), each \(B\in\mathcal B\) and each
required direction, choose the unit \(v_B\) from the proof of
Lemma~\ref{lem:abelianunits744}, now viewed in \(F_f\).
Lift its reciprocity value for \(F_f\) to \(s_{f,B}\in I_{F_f}=I^0\).
Put in \(H_w\) the two factors \([s_{f,B}]-v_B^{\pm B}[1]\),
in any fixed order. On a nontrivial inertia character of any
of the members, the factor for its own \(f\), a detecting direction
and its signed exponent vanishes. All factors preserve that
character space. On the trivial character space \(v\in I^0\)
already acts as one.

The unit estimates in Lemma~\ref{lem:abelianunits744} depend
only on \(B\), the direction and \(\ell\). Lifting by reciprocity
for \(F_f\) changes none of their augmentations. For one \(B\)
the sum of the valuations of these augmentations over primes,
with weights \(\log\ell\), is \(O_d((1+B)^{1/4})\) for each \(f\).
The fixed-node bound is therefore
\[
    O_{q_0,d}\big(n_x R_x|\mathcal B|(1+N_x b_x)^{1/4}\big).
\]
Since \(m_x\le2\sqrt x\), \(N_x\ll x^{17}\), and \(t_0\ll x^5\),
the power of \(x\) here is at most
\(17+\tfrac12+34+22+17/4<80\).

At an old prime there are \(Qf\) embeddings of \(F_f\) over
\(\QQ_p\), with \(Q\) above each of the \(f\) residue embeddings.
Let \(\chi\) be any character of \(I^0\) occurring in a potentially
abelian member. Its stabilizer under
\(\Delta=I_K/I^0=\operatorname{Gal}(F/K)\) has index \(j\le R\).
The action here is well-defined since conjugation by \(I^0\)
fixes every such character. Use its crystalline line on \(G_{F_f}\)
from Lemma~\ref{lem:potinertia746}. Each stabilizer element can
be lifted to \(I_K\), acting trivially on the unramified degree
\(f\) subfield. The field \(F_f\), the compositum of that subfield
and \(F\), is preserved. Conjugation of the line character by this
lift has the same restriction to \(I^0\). Thus their algebraic
monomials agree on units by naturality of local reciprocity.
Their labelled exponents agree. Indeed take logarithms on sufficiently
deep principal units to obtain equality of the corresponding
linear combinations of embeddings on a full lattice in \(F_f\),
hence on \(F_f\). Distinct field embeddings are linearly
independent. The jumps are therefore constant along each orbit
of the stabilizer on embeddings.

At each residue embedding \(i\), \(0\le i<f\), let \(A_i\)
be the sum of the \(Q\) exponents before the common reciprocity
sign. There are \(j\) stabilizer orbits there, each of size
\(Q/j\), with jumps in \(\{0,\ldots,t'-1,J_w\}\). If \(J_w<0\)
then \(|J_w|>R_xt_0\), so even one negative choice among the \(j\)
terms outweighs every possible sum of the nonnegative choices.
If \(J_w>0\), all terms are nonnegative. Thus \(A_i=0\) precisely
when all its jumps vanish. Also \(|A_i|\le5Qb_x<p-1\)
by the size of the old primes. If \(\chi\) is nontrivial,
some \(A_i\ne0\). Its exponent on a Teichm\"uller generator
of the residue field of \(F_f\) is, up to sign,
\(\sum_{i=0}^{f-1}p^i A_i\), with the residue embeddings in
Frobenius order. This integer has absolute value less than
\(p^f-1\). If it were zero, reduction modulo \(p\) and
\(|A_i|<p-1\) would imply \(A_0=0\), and division and induction
would force all the others to vanish. The character is consequently
nontrivial on such a generator.

Take the element \(s\) of Lemma~\ref{lem:wildgenerators745}.
Its power \(s^Q\) generates the tame quotient of \(I^0\)
topologically, since that quotient is the index \(Q\) subgroup
of the procyclic tame quotient for \(K\). For every \(f\)
its image in the abelian inertia of \(F_f\) is a Teichm\"uller
generator. Indeed wild inertia maps onto principal units and
the prime-to-\(p\) unit quotient is the cyclic residue group,
while \(s^Q\) has pro-prime-to-\(p\) order. Thus a character
\(\chi\) as above with \(\chi(s^Q)=1\) is trivial on all of \(I^0\).
Take \(H_w\) to be the average of the powers of \([s^Q]\) from
zero to \(m'-1\), with \(m'\) a common prime-to-\(p\) multiple
of the finite orders of \(u_y(s^Q)\) as in
Proposition~\ref{prop:tameblocks745}. Its augmentation is one,
and its evaluation projects onto the \(s^Q\)-fixed space. This
proves the assertion for these places, including integral
lattices by embedding them in their rational spaces.
\end{proof}

\begin{theorem}[Period contribution outside the potentially abelian members]
\leavevmode\label{thm:potentialperiod746}
For the marked separating families of Theorem~\ref{thm:projectedcups742}
there are numbers \(0\le r_{\rm pot}(D)\le r(D)\) determined by
trace evaluations on integral matrix tuples with zero component
at every member potentially abelian at the place of evaluation.
They satisfy
\[
    \sum_\ell(\log\ell)\sup_{D\ {\rm at}\ \ell}
       \frac{(p(D)-r_{\rm pot}(D))_+}{\sqrt{b_x}}
                 \ll_{q_0,d} V_0^{-10}.
\]
In particular the full prime sum of \(p(D)\) satisfies this bound
on the subfamily for which every member is potentially abelian
at every coefficient place. At an original level prime all
the blocks in Theorem~\ref{thm:oldcolumn745} already have zero
component at each locally potentially abelian member.
\end{theorem}
\begin{proof}
At a spherical place take \(M=T\) or \(M=\mathcal O(-h)\),
\(1\le h<t'\), and write \(\alpha\) for its character. Choose
\(s_w\in I^0=I_K\) minimizing \(v_\ell(\lambda_w)\) for
\(\lambda_w=\alpha(s_w)-1\), as in
Theorem~\ref{thm:abelianperiod744}. For every \(v\in I^0\),
\[
 f_c(u([s_w,v])-1)=
      \lambda_w c(v)-(\alpha(v)-1)c(s_w),
\]
where \([s_w,v]\) denotes the group commutator. By
Lemma~\ref{lem:potinertia746} its matrix difference is zero at
all potentially abelian members. If the indicated evaluations
of two trace cocycles agree, the restriction of the difference
of their classes multiplied by \(\lambda_w\) is zero on \(I^0\).
By inflation and restriction this multiplied class is inflated
from \(H^1(G_K/I^0,(M/\pi^a)^{I^0})\). Its coefficients are
killed by \(\lambda_w\), as \(s_w\in I^0\). Thus the evaluations
determine all free cup coordinates of \(\lambda_w^2[c]\)
on \(G_K\).

For \(M=\mathcal O\) at a spherical place, evaluate \(f_c\)
on \(u(([v]-[1])H_w)\) for \(v\in I^0\), now using group
algebra symbols. The tuple has zero augmentation before
evaluation and belongs to \(\mathcal N\). It has zero components
as in the proposition, and the additive cocycle evaluates to
\(\delta_w c(v)\). These values determine all free cup
coordinates of \(\delta_w[c]\), since unramified additive
classes lift and have zero free cups.

At an original level prime keep the exact calculation of
Theorem~\ref{thm:oldcolumn745}. In a potentially abelian member,
\(g\in P\subset I^0\) commutes in its matrix action with
\(s^Q\) by Lemma~\ref{lem:potinertia746}. For \(h>0\) the
\(s\)-spaces at \(1\) and \(\eta_h\) have distinct \(s^Q\)
eigenvalues since \(0<Qh<p-1\). Hence the block with these
source and target spaces is zero. The \(s\)-fixed space is
contained in the \(s^Q\)-fixed space, on which \(I^0\) acts
trivially by the old-prime proof of Proposition~\ref{prop:potunits746}.
This gives the assertion for \(Z_0\) also.

Define \(r_{\rm pot}(D)\) from the same joint cup images as \(r(D)\),
multiplying each spherical coordinate by \(\lambda_w^2\) or
\(\delta_w\) according to its weight. At an old coefficient prime
use \(B_{\rm old}(D)\) unchanged. Set \(r_{\rm pot}=0\) at a zero lift.
The kernel of multiplication by \(v\in\mathcal O\setminus\{0\}\)
on \(d\) copies of \(\mathcal O/\pi^a\) has normalized length
at most \(d\,v_\ell(v)\). The same upper bound holds on any
joint submodule by inclusion, adding the bounds for its coordinates.
Each local test has at most \([E_{x,w}:\QQ_\ell]\) coordinates
per weight. For the main and positive scalar weights, the sum
of the costs from \(\lambda_w^2\) is at most twice the unit
cost already bounded at the end of Theorem~\ref{thm:abelianperiod744}.
For weight zero use Proposition~\ref{prop:potunits746}.
The resulting prime sum at each node for \(r(D)-r_{\rm pot}(D)\)
is \(O_{q_0,d}(x^{80}b_x^{1/4})\), uniformly in the same data.
Divide by \(\sqrt{b_x}\) and sum over nodes \(x\ge V_0\).
Since \(b_x\ge x^{500}\), this costs \(O_{q_0,d}(V_0^{-44})\).
Combining with \eqref{eq:projectedcups742} proves the displayed
estimate. When all members satisfy the stated local hypothesis
all determining tuples are zero, so \(r_{\rm pot}=0\) there.
\end{proof}

The original separating families allow representations having
nonabelian image on every open subgroup at a coefficient place.
The theorem leaves their integral trace evaluations to be bounded.
Its prime sum applies on the same fields, intervals and signed
twists as the critical-ratio estimates.
\clearpage

\subsection{Determinants of two-dimensional local constituents}
\label{sec:twodeterminants754}
We return to the coefficient primes away from the two original
levels. There the local trace matrices are crystalline, but
the width of their weights can exceed the residue characteristic.
A determinant identity removes two-dimensional direct summands
from the trace equations whenever their determinants and the
target character have different exponents on suitable units.
Its cost in the projected period is summable on the original
interval, including at ramified completions of the period field.

\begin{lemma}[Removing two-dimensional blocks from trace equations]
\leavevmode\label{lem:twoblocks754}
Let \(\alpha:G\to\mathcal O^\times\) be a character and \(c\) an
\(\alpha\)-cocycle modulo \(\pi^a\) satisfying the integral
matrix relations of a finite family \(u_y\), in the sense of
Lemma~\ref{lem:graphrealization724}. Fix rational
\(G\)-stable direct sum decompositions of the \(u_y\), after
extending coefficients if necessary. Choose \(s\in G\),
\(\alpha(s)\ne1\), and a finite set \(\mathcal D\) of units of
\(\mathcal O\) different from \(\alpha(s)\). Put
\[
                k_s=(\alpha(s)-1)
                        \prod_{d\in\mathcal D}(\alpha(s)-d).
\]
The class \(k_s[c]\) has a representative satisfying the matrix
relations on the remaining summands after deleting every line
and every two-dimensional block \(W\) with
\(\det W(s)\in\mathcal D\). There is no dependence on the index
of an integral direct sum inside the supplied lattice.
\end{lemma}
\begin{proof}
On the group ring write \(\epsilon\) for augmentation, extend
\(\alpha\) linearly, and put \(d_c(\sum a_g[g])=\sum a_g c(g)\).
The cocycle equation gives
\[
      d_c(XY)=\alpha(X)d_c(Y)+d_c(X)\epsilon(Y).
\]
Let \(S=[s]\), let \(B_s(X)=SXS\), and put
\[
   F_s(X)=
      \left(\prod_{d\in\mathcal D}(B_s-d\,\mathrm{id})\right)
                     (SX-XS).
\]
For a two by two matrix \(A\) and any matrix \(X_0\) of the
same size, write \(C=AX_0-X_0A\). Cayley--Hamilton gives
\(AC+CA=A^2X_0-X_0A^2=(\operatorname{tr}A)C\).
Multiplying by \(A\) on the left and substituting its
characteristic equation shows \(ACA=(\det A)C\). Thus
\(F_s\) is zero on the specified two-dimensional blocks, as
well as on lines. On each other block it is a linear operator
on the evaluated matrix.

The commutator has zero augmentation and zero \(\alpha\)-value.
These equalities are preserved by \(B_s\) and on this subspace
\(d_c(B_sX)=\alpha(s)d_c(X)\). For \(X=[g]-1\) we obtain
\[
 d_c(F_s([g]-1))
 = k_s c(g)-
    (\alpha(g)-1)c(s)\prod_{d\in\mathcal D}(\alpha(s)-d).
\]
The last term is a coboundary as a function of \(g\). A finite
integral relation on the remaining matrices maps under \(F_s\)
to zero on every original block. It is therefore zero on the
original matrices in their original lattices. Its augmentation
is zero, so the defining relation for \(c\) says its \(d_c\)
value is zero. This proves the lemma. Repeating the construction
on the remaining family multiplies the scalars \(k_s\).
\end{proof}

\begin{proposition}[The two weight patterns left by the unit tests]
\leavevmode\label{prop:twoweights754}
Use the marked torus families of Theorem~\ref{thm:projectedcups742}
at a spherical coefficient place \(w\) away from the original
levels. Write \(K=E_{x,w}\), \(n_x=[E_x:\QQ]\), \(B=t'-1\),
and \(J_\tau=z-1+\epsilon_\tau k_\tau\) for the main jump at
an embedding \(\tau\) of \(K\).
Take \(\alpha\) to be a positive scalar target \(\kappa^{-h}\),
\(1\le h\le B\), or the character of the twisted main line.
For any chosen rational direct sum decompositions of the
local matrices there is a scalar \(k_{\alpha,w}\ne0\), independent
of precision, such that every scaled local trace class has a
representative satisfying the relations on the following family:
\begin{enumerate}
\item keep the blocks of dimension at least three;
\item for a scalar target keep only those two-dimensional
blocks with both jumps in \([0,B]\) at every embedding;
\item for the main target keep only those two-dimensional
blocks whose jumps at every embedding are \(\{0,J_\tau\}\).
\end{enumerate}
One may retain additional summands in all three statements.
With \(v_\ell(\ell)=1\), the cost satisfies at each node
\[
 \sum_\ell(\log\ell)
    \sup_{D\ {\rm at}\ (x,\ell)}
     \sum_{\alpha,w}[E_{x,w}:\QQ_\ell]\,v_\ell(k_{\alpha,w})
                 \ll_{q_0,d} x^{63}b_x^{1/4}.
\]
The sum inside the supremum runs over the specified nontrivial
targets and spherical places. The bound is independent of the
number of trace systems and the choice of decompositions.
\end{proposition}
\begin{proof}
Crystallinity at the spherical places and the common labelled
weights were verified in Lemma~\ref{lem:abelianunits744}.
Any rational direct summand \(W\) is crystalline. If its
dimension is two, filtered exactness gives its determinant
weight as the sum of its two jumps from the ambient multiset,
with multiplicities. On local units \(\det W/\alpha\) has
the algebraic character with exponents
\[
 d_\tau=(\text{sum of the two jumps of }W\text{ at }\tau)
                          -(\text{jump of }\alpha\text{ at }\tau)
\]
up to the common reciprocity sign. This uses the rank-one
description on units in Lemma~\ref{lem:abelianunits744} applied
to the determinant. The same description holds for \(\alpha\).

The rational-unit direction measures \(\sum_\tau d_\tau\).
At a place where the quadratic CM field is nonsplit, the
norm-one units of that quadratic completion measure
\(\sum_\tau\epsilon_\tau d_\tau\), up to sign. For a positive
scalar target, if \(W\) contains the main jump at any label,
each such occurrence contributes one \(k_\tau\) with the
same sign to the oriented sum. At a split prime one can instead
use the rational sum, since the orientations there agree.
The other parts together have absolute value at most
\(4n_xt_0\). The inequality
\(\min k_\tau\ge3r_x^8-4r_x>4n_xt_0\) holds throughout the
interval after increasing its initial node. Thus a block
containing any main jump is detected.

For the main target, wherever \(W\) contains \(J_\tau\) that
jump cancels. At every other label exactly one far term occurs
with the opposite sign. The same oriented or rational sum
detects the block if there is any such label. If \(W\) contains
\(J_\tau\) everywhere, the difference consists of its other
scalar jumps, all nonnegative. The rational sum detects it
unless all those jumps are zero. The target \(\alpha\) itself
is nontrivial on at least one available direction: use \(h>0\)
for scalar targets and the same far-jump argument for the main
line. This proves the required distinctions between characters.

We give the group elements and the cost explicitly. Choose
one unit in each available direction. In the rational direction
take a Teichm\"uller generator times \(1+\ell^{c_\ell}\),
where \(c_2=2\) and \(c_\ell=1\) at odd primes. For an inert
quadratic field take a norm-one Teichm\"uller generator times
the norm-one principal unit of depth \(c_\ell\) from the proof
of Theorem~\ref{thm:alltorsion741}. At the fixed ramified
quadratic prime use its norm-one unit of depth \(3/2\) there.
For any nonzero integer \(j\) the cost of a direction unit \(v\)
is \(v_\ell(v^j-1)\). It is zero for a nontrivial residue value.
Otherwise it is the indicated depth plus \(v_\ell(j)\) by
the binomial formula. The rational and inert positive costs
require \(\ell-1\mid j\) and \(\ell+1\mid j\), respectively.
The ramified contribution is at one fixed prime.
Thus, for fixed \(j\), the sum of all these costs with weights
\(\log\ell\) is \(O_d((1+|j|)^{1/4})\). Indeed the residue
conditions involve at most \(2\tau(|j|)\) primes, each at
most \(|j|+1\), while
\(\sum_\ell v_\ell(j)\log\ell=\log|j|\).
Use the divisor bound as in Lemma~\ref{lem:abelianunits744}.

At the node form a list \(\mathcal J_x\) of signed nonzero
integers. For each of the at most \(m_x\) possible chain
indices take the two base weights \(w_1,w_2\) used in
Lemma~\ref{lem:abelianunits744}. Include all nonzero
\[
         i_1w_1+i_2w_2+D_0,\qquad
     |i_1|,|i_2|\le n_x,\qquad |D_0|\le 6n_x t_0 .
\]
This includes all signed exponents just used, for both
\(\det W/\alpha\) and \(\alpha\). At a label the far part
is present at most once after the subtraction, with either
sign. Its magnitude \(k_\tau\) is one of \(w_1,w_2\) minus
an integer in \([0,4r_x]\); the centering and the scalar
jumps fit in the stated bound. Consequently
\(\#\mathcal J_x\ll(m_x+1)(n_x+1)^2(1+n_xt_0)\), and
\(\max|j|\ll n_xb_x\).

Choose as \(v_1\) a direction on which \(\alpha\) has nonzero
exponent, and lift its reciprocity value to \(s\in I_K\).
Use Lemma~\ref{lem:twoblocks754} with
\(\mathcal D=\{\alpha(s)v_1^j:j\in\mathcal J_x\}\).
All blocks detected in this direction are removed. If there
is a second direction \(v_2\), the determinant ratios left
to detect are trivial on \(v_1\). Use either
\(v_1v_2\) or \(v_1^2v_2\), lifting its reciprocity value to
\(s'\). One can choose it so that
\[
  v_\ell(\alpha(s')-1)\le v_\ell(\alpha(s)-1).
\]
In fact the quotient of the two candidate \(\alpha\)-values is
\(\alpha(s)\), so the ultrametric inequality excludes a larger
valuation for both. The remaining nontrivial ratios on this
unit are unchanged from \(v_2\). Apply the lemma with
\(\{\alpha(s')v_2^j:j\in\mathcal J_x\}\). These sets of units
avoid the excluded value in both uses because \(j\ne0\) and
the direction units have nonzero principal logarithm.
Any further summands kept in the statement impose fewer
zero relations on a cocycle.

The valuations of the factors \(\alpha(s)-d\) are precisely
the direction costs for the listed integers, and the additional
two commutator factors cost at most twice the first direction
cost for the exponent of \(\alpha\). The full list works before
taking any suprema. Summing the local degrees gives \(n_x\)
for each of at most \(t_0+1\) targets. The resulting prime sum
at the node is
\[
 \ll_d n_x(t_0+1)(m_x+1)(n_x+1)^2
                      (1+n_xt_0)(1+n_x b_x)^{1/4}.
\]
Use \(n_x\ll_{q_0,d}x^{12}\), \(m_x\le2\sqrt x\) and
\(t_0\ll x^5\) to obtain the stated bound.
\end{proof}

\begin{corollary}[The residual joint image on the period interval]
\leavevmode\label{cor:twoperiod754}
In Theorem~\ref{thm:projectedcups742} there are numbers
\(0\le r_2(D)\le r(D)\) such that
\[
 \sum_\ell(\log\ell)\sup_{D\ {\rm at}\ \ell}
        \frac{(p(D)-r_2(D))_+}{\sqrt{b_x}}
                  \ll_{q_0,d} V_0^{-10}.
\]
At both original coefficient primes \(r_2(D)=0\). At every
spherical place the nontrivial targets in \(r_2(D)\) can be
evaluated on cocycles satisfying the reduced matrix relations
of Proposition~\ref{prop:twoweights754}. At an unramified
place with \(\ell>B+1\), \(\ell>2\), the positive scalar
coordinates of the reduced image vanish whenever all local
blocks have dimension at most two. At any spherical place the
main coordinate vanishes if all local blocks have dimension at
most two and no two-dimensional block has jumps \(\{0,J_\tau\}\)
at every label.
\end{corollary}
\begin{proof}
In the joint free cup map defining \(r(D)\), multiply each
nontrivial target coordinate at a spherical place by the
corresponding \(k_{\alpha,w}\). Keep the weight-zero map there
and any other untreated coordinates as given. Take \(r_2(D)\)
to be the sum of the normalized lengths of the images of the
original global trace classes under these joint maps.
Corollary~\ref{cor:oldperiod752} covers the two original primes.
Adding a coboundary preserves each local cup evaluation, so the
new values can indeed be computed with the reduced relations.
The kernel of multiplication by \(k_{\alpha,w}\) in each
coordinate of \(\mathcal O/\pi^a\) has normalized length at
most \(v_\ell(k_{\alpha,w})\). The local free ranks are at
most \([E_{x,w}:\QQ_\ell]\) for each target by
Corollary~\ref{cor:torustests741}. Intersecting this direct
sum of kernels with the joint global image cannot increase its
length. The proposition bounds the weighted loss at a node by
\(O_{q_0,d}(x^{63} b_x^{1/4})\); dividing by \(\sqrt{b_x}\)
and summing over \(x\ge V_0\) gives \(O_{q_0,d}(V_0^{-61})\)
because \(b_x\ge x^{500}\). Adding the inequality of
Theorem~\ref{thm:projectedcups742} proves the claim.

For the stated scalar coordinates the remaining family, its
target and the quotient all have jumps in \([0,B]\). At the
specified unramified places \(B\le\ell-2\), so
Theorem~\ref{thm:labelledlift735} puts the cocycles in the
finite condition. For the stated main coordinate no summand
remains; the reduced relations force the representative to
vanish on every group element. Both assertions apply to the
actual trace precision. The weight-zero target and the
unrestricted higher-dimensional blocks at the spherical primes
are still present in the general interval estimate.
\end{proof}
\clearpage

\subsection{Residual constituents and local trace relations}
\label{sec:residualtrace755}
The local trace systems can be nonsplit in characteristic zero.
We identify which part of each rational filtration can affect
the class of a scalar cocycle. This identification is exact
over the original coefficient ring. For the trivial target it
also applies to the cocycle itself, including its restriction
to inertia.

\begin{theorem}[An interval in a filtration that preserves trace classes]
\leavevmode\label{thm:residualinterval755}
Let \(\mathcal O\) be the ring of integers of a finite extension
of \(\QQ_\ell\), with uniformizer \(\pi\) and residue field \(k\).
Let \(u_y\) be a finite family of continuous representations of
a profinite group \(G\) on free \(\mathcal O\)-lattices, and let
\(\alpha:G\to\mathcal O^\times\) be continuous. In each rational
representation choose a \(G\)-stable filtration
\[
                  0=V_0\subset V_1\subset\cdots\subset V_m=V .
\]
Give each factor \(V_j/V_{j-1}\) the lattice induced from the
original one. Let \(i\) be the first index whose factor contains
\(\bar\alpha\) as a composition factor after reduction modulo
\(\pi\); let \(j\) be the last index whose reduced factor
contains the trivial \(k\)-character. Replace \(V\) by
\(V_j/V_{i-1}\). Use the zero space if either index is absent
or if \(j<i\).

At every precision \(\pi^a\), the image in
\(H^1(G,(\mathcal O/\pi^a)(\alpha))\) of the cocycles satisfying
the matrix relations is identical for the original and the
replacement families. If \(\bar\alpha=1\), the modules of
raw cocycles satisfying the two families of relations are
themselves identical.
\end{theorem}
\begin{proof}
Use the \(\mathcal O\)-algebra \(A\) which is the image of
\(\mathcal O[G]\) under the product of the \(u_y\), the
character \(\alpha\), and augmentation \(\epsilon\). It is a
finite free \(\mathcal O\)-module, being a submodule of the
finite free product of matrix and scalar rings. In particular
it is complete. The two scalar factors give characters of \(A\).

For each distinct character \(\lambda\) in
\(\{\bar\alpha,\bar\epsilon\}\), choose the idempotent of the
semisimple quotient of \(A/\pi A\) that acts as the identity
on that simple module and as zero on the others. Such an
idempotent exists by the structure theorem for semisimple
algebras. The simple component for a \(k\)-valued character is
\(k\) itself since the scalar subalgebra already maps onto
\(k\). Lift the idempotents to orthogonal idempotents \(P_\lambda\)
of \(A\). For details, the radical of the finite-dimensional
residue algebra is nilpotent. Starting with any preimage \(x\),
the substitution \(x\mapsto 3x^2-2x^3\) preserves its residue
in the semisimple quotient and squares the ideal of its
idempotency error, since the new error is a multiple of
\((x^2-x)^2\). This first lifts through the nilpotent radical,
then through all powers of \(\pi\) by completeness; successive
changes are multiples of the old error. Lift a second idempotent
in the corner complementary to the first. When the two
characters agree use the same lift for both. We denote them by
\(P_\alpha,P_1\). Their scalar values
\(\alpha(P_\alpha)=\epsilon(P_1)=1\) follow from their residues,
since the only idempotents in \(\mathcal O\) are zero and one.

For a lattice factor \(L\), projection by \(P_\lambda\) on
\(L/\pi L\) has rank the number of its \(\lambda\)-composition
factors. Indeed projection by this fixed idempotent is exact
on underlying vector spaces in any exact sequence of modules:
one projects a lift of a vector to prove surjectivity, and
intersects with the kernel to prove exactness in the middle.
On simple modules its ranks are one or zero as prescribed.
The image \(P_\lambda L\) is a direct summand as an
\(\mathcal O\)-module, so its rank before reduction is the same.
This argument also works over \(k\) when other simple modules
are not absolutely irreducible.

Consequently \(P_1V\subset V_j\), and \(P_\alpha\) kills
\(V_{i-1}\). For instance the projection on \(V/V_j\) has
rank zero by applying the rank calculation to its filtration.
When either character is missing or \(j<i\), the operator
\(P_\alpha X P_1\) is zero on \(V\) for every \(X\in A\).
Otherwise it vanishes on \(V\) whenever \(X\) vanishes on
\(V_j/V_{i-1}\), since then \(X(P_1V)\subset V_{i-1}\).

For a raw trace cocycle \(c\) use the linear map \(d_c\) of
Lemma~\ref{lem:twoblocks754}. It factors through \(A\): an
element of the group-ring kernel has augmentation zero and
acts as zero on all the given matrices, so its value under
\(d_c\) vanishes by the original relations. In \(A\), for
\(X=[g]-1\), the product rule gives
\[
 d_c(P_\alpha X P_1)
        = c(g)+(\alpha(g)-1)d_c(P_1).
\]
This differs from \(c(g)\) by a coboundary. If a linear
combination of the \(X\)'s vanishes on all replacement
modules, its projection between the two idempotents vanishes
on every original matrix. The projected element has
augmentation zero; lifting it to the group ring gives a
relation evaluated as zero by \(d_c\). Hence the displayed
representative satisfies the replacement relations. In the
reverse direction a relation on all original matrices always
holds on all their rational subquotients, so every replacement
trace cocycle already satisfies the original relations. All
identities concern zero operators, and no lattice projection
has been divided by an index.

If \(\bar\alpha=1\), then \(\alpha(P_1)=\epsilon(P_1)=1\).
Evaluating \(P_1^2=P_1\) gives \(d_c(P_1)=2d_c(P_1)\), so
the displayed representative is \(c\) itself. No division
by two occurs, including when \(\ell=2\).
\end{proof}

\begin{corollary}[The weight-zero test in an arbitrary rational filtration]
\leavevmode\label{cor:zerointerval755}
At a spherical coefficient place in the marked torus family,
suppose every local trace representation has a rational
filtration for which the part from the first to the last
occurrence of the trivial residual character is potentially
abelian; an empty part is permitted. Then multiplication by
the augmentation \(\delta_w\) of
Proposition~\ref{prop:potunits746} makes every trivial-target
trace cocycle unramified at \(w\), at every precision. In
particular this applies if the trivial residual character
occurs only in a single rank-one rational filtration factor
of each representation, or is absent there.
\end{corollary}
\begin{proof}
The theorem for \(\alpha=1\) gives the original cocycle on the
reduced family at this place. Those rational subquotients are
crystalline and each of their labelled weights occurs in the
original trace system, by filtered exactness. Their ranks are
at most \(2z\). The same \(H_w\) as in
Proposition~\ref{prop:potunits746} therefore kills their
nontrivial inertia parts. To see the uniformity, its product
includes each unramified orbit length up to \(2z\) and every
nonzero exponent formed from allowed jumps at the embeddings
of that extension. Passing to a rational subquotient only
restricts those possibilities. Thus for \(v\in I_{E_{x,w}}\)
the element \(([v]-1)H_w\) acts as zero on the replacement
matrices and has augmentation zero. Evaluating it by the
additive cocycle gives \(\delta_w c(v)=0\).
\end{proof}

\begin{theorem}[Omitting places detected by the residual constituents]
\leavevmode\label{thm:filteredperiod755}
Use the same fields, tests and interval as
Theorem~\ref{thm:projectedcups742}. For each target \(M\) and
each spherical coefficient place \(w\), apply
Theorem~\ref{thm:residualinterval755} to the restricted family
with \(\alpha\) the character of \(M\). Omit the local free cup
coordinates for that target at \(w\) whenever one of the
following conditions holds for the resulting family:
\begin{enumerate}
\item all its representations are potentially abelian;
\item \(M=\mathcal O(-h)\), \(1\le h\le B=t'-1\), the place
is unramified over \(\QQ_\ell\) with \(\ell>\max(2,B+1)\),
and the representations are rational direct sums of blocks
of dimension at most two;
\item \(M\) is the main line, the representations are rational
direct sums of blocks of dimension at most two, and none of
their two-dimensional blocks has jumps \(\{0,J_\tau\}\) at
all the labels of \(w\).
\end{enumerate}
The filtrations and direct sums may depend on the target.
Omit both original coefficient primes by
Corollary~\ref{cor:oldperiod752}. Define \(r_{\rm fil}(D)\)
by projecting the original joint images onto the coordinates
retained and using the same normalized lengths. Then
\[
 \sum_\ell(\log\ell)\sup_{D\ {\rm at}\ \ell}
      \frac{(p(D)-r_{\rm fil}(D))_+}{\sqrt{b_x}}
                         \ll_{q_0,d}V_0^{-10}.
\]
In particular the estimate holds for the full projected
pairing on the data where these omissions cover every
coefficient place and target.
\end{theorem}
\begin{proof}
Local finite conditions and cup maps depend only on the
cohomology class. Hence for any global trace class the
local representatives supplied by
Theorem~\ref{thm:residualinterval755} can be used at the
individual places, without changing its local values.
The second and third cases of the statement now follow
pointwise from Proposition~\ref{prop:twoweights754} and the
last paragraph of Corollary~\ref{cor:twoperiod754}. Indeed
every labelled jump of a replacement block occurs in the
ambient Hodge multiset with the same multiplicity bound; this
is all that the determinant test requires. Its local
annihilating scalar has the same valuation bound, using the
same exponent list and the same degrees.

For the first case at weight zero use
Corollary~\ref{cor:zerointerval755}. At a nontrivial target
choose \(s\in I_K\), \(K=E_{x,w}\), with
\(\lambda=\alpha(s)-1\ne0\) as in
Theorem~\ref{thm:abelianperiod744}. The inertia matrices of
every potentially abelian replacement commute, by
Lemma~\ref{lem:potinertia746}. Evaluating the group commutator
of \(s\) and \(v\in I_K\) therefore gives
\(\lambda c(v)=(\alpha(v)-1)c(s)\) for the replacement
representative. Subtract this coboundary after multiplying by
\(\lambda\). The class then restricts trivially to inertia.
Inflation and restriction place it in the image of
\(H^1(G_K/I_K,(M/\pi^a)^{I_K})\). Its coefficients are killed
by \(\lambda\) since \(s\in I_K\). Thus \(\lambda^2\) kills
the original local class. The exponents and unit bounds for
these choices concern the target alone.

Each omitted coordinate is consequently killed by the
corresponding scalar. Its normalized length is at most the
valuation of that scalar. Projection of a joint image has
kernel contained in the direct sum of the omitted coordinates;
this accounts for possible dependencies between places.
The nodewise prime sum for the first case is at most
\(O_{q_0,d}(x^{80}b_x^{1/4})\) by the unit estimates in
Theorem~\ref{thm:potentialperiod746}; the other two cases
cost \(O_{q_0,d}(x^{63}b_x^{1/4})\) by
Proposition~\ref{prop:twoweights754}. Both bounds are uniform
in the filtrations even before taking the supremum. Now
\(\sum_{x\ge V_0} x^{80}b_x^{-1/4}\ll V_0^{-44}\).
Apply the prime-sum inequality of
Theorem~\ref{thm:projectedcups742} on this same interval.
\end{proof}
The theorem uses a filtration of the actual local matrices.
When the two residual characters occur together in a large
irreducible rational constituent, that constituent survives
the reduction. The full period inequality still requires
control of the surviving joint image for unrestricted data.
\clearpage

\subsection{Ordered crystalline flags and the integral projected period}
\label{sec:orderedflags756}
At some coefficient places the crystalline matrices admit an invariant
flag whose characters have increasing Hodge jumps. An extension in
such a flag need not split, even over the fraction field. We bound
its contribution to the integral trace equations using the positions
of the trivial character and of the target in the flag.

\begin{lemma}[Two operators on a filtered trace representation]
\leavevmode\label{lem:flagoperators756}
Let the notation for integral matrix relations and for a character
\(\alpha\) be as in Lemma~\ref{lem:twoblocks754}. Suppose the rational
representations have \(G\)-stable full flags. Write \(\beta_i\) for the
character on \(V_i/V_{i-1}\), with indices increasing from the bottom
of a flag, and let \(R\) bound the ranks. Take \(A,B\in\mathcal O[G]\)
with \(\alpha(A)\epsilon(B)\ne0\), where \(\epsilon\) denotes
augmentation. In each flag let \(i\) be the first index with
\(\beta_i(A)\ne0\) and \(j\) the last index with \(\beta_j(B)\ne0\).
Use the quotient \(V_j/V_{i-1}\) if both indices exist and \(i\le j\).
In all other cases use zero. If \(c\) satisfies the original matrix
relations modulo \(\pi^a\), then
\[
                     k[c],\qquad k=\alpha(A)^R\epsilon(B)^R,
\]
has a representative satisfying the relations on the resulting
quotients. For \(\alpha=1\) the representative is \(kc\).
\end{lemma}
\begin{proof}
On a string of successive factors on which \(B\) is zero, \(B\)
lowers the filtration by one step. Applied to \(V/V_j\), this gives
\(B^RV\subset V_j\). In the same way \(A^RV_{i-1}=0\). Absence
of an index makes the corresponding power zero on the entire
representation. For \(X\) in the augmentation ideal, multiplication
by \(A^R\) on the left and \(B^R\) on the right therefore gives
the zero matrix on every original representation whenever \(X\)
acts as zero on the indicated quotients. This also covers \(j<i\)
because every group ring element preserves \(V_j\).

Let \(d_c(\sum a_g[g])=\sum a_gc(g)\). Applying the cocycle product
rule of Lemma~\ref{lem:twoblocks754} twice gives
\[
 d_c\big(A^R([g]-1)B^R\big)
  =k c(g)+(\alpha(g)-1)\alpha(A)^R d_c(B^R).
\]
The last summand is a coboundary. For a linear combination of the
\([g]-1\) giving a relation on the quotients, the product in
parentheses has augmentation zero and acts as zero on the original
lattices, as we just proved on their rational spaces. Its value
under \(d_c\) hence vanishes. This proves the assertion at the
original precision.
\end{proof}

\begin{proposition}[The two ends of a crystalline flag]
\leavevmode\label{prop:ordinaryunits756}
Use the fields and weights of Proposition~\ref{prop:twoweights754}
at a spherical coefficient place \(w\). Write \(K=E_{x,w}\).
We allow each local representation to be a rational subquotient
of one of the trace systems. For a full \(G_K\)-stable flag over a
finite coefficient extension call it ordered if, at every embedding
\(\tau\) of \(K\), its line characters have increasing jumps in flag
order. For the trivial target assume only that full stable flags
exist. Then there is a nonzero scalar \(k_{1,w}\) for which every
trivial-target trace cocycle satisfies \(k_{1,w}c(I_K)=0\).

For a nontrivial target assume the flags are ordered. For each
positive scalar target and for the main target when at least one
\(J_\tau>0\), some nonzero \(k_{\alpha,w}\) kills every local
trace class in \(H^1(K,(\mathcal O/\pi^a)(\alpha))\). These assertions
hold at every precision. The scalars can be chosen with
\[
 \sum_\ell(\log\ell)\sup_{D\ {\rm at}\ (x,\ell)}
       \sum_{\alpha,w}[E_{x,w}:\QQ_\ell]v_\ell(k_{\alpha,w})
                   \ll_{q_0,d} x^{68}b_x^{1/4}.
\]
For this estimate a place or target may be omitted. Ordered flags
may be assumed to exist after a finite extension of \(K\).
\end{proposition}
\begin{proof}
We first justify the last assertion. Pass to a finite normal
extension over which a proposed ordered flag is stable. On
restriction the labelled jumps above an embedding of \(K\) are
repetitions of its original multiset. At each individual embedding
they are distinct. This regularity follows from the list
\(0,\ldots,B,J_\tau\), where \(B=t'-1\) and \(J_\tau\notin[0,B]\).
An ordered flag on this restriction is unique. Indeed two distinct
initial lines would give a direct sum subrepresentation with two
copies of the least jump at any fixed embedding. Exactness of the
de Rham filtration on crystalline representations excludes this.
Apply the same argument after quotienting by the common first line.
Conjugation by \(G_K\) permutes the embeddings above each base
embedding and takes an ordered flag to an ordered flag. Uniqueness
therefore makes each subspace \(G_K\)-stable. The characters of
this flag have the claimed order over \(K\) by base change.

Use now the signed nonzero exponent list \(\mathcal J_x\) and the
units of Proposition~\ref{prop:twoweights754}. Choose a rational
unit \(v\) there and, at a nonsplit prime of the quadratic CM field,
also its specified norm-one unit. For each direction choose
\(s_v\in I_K\) lifting its local reciprocity image. For
\(\gamma=1,\alpha\) form in the integral group algebra
\[
 H_\gamma=\prod_{v}\prod_{j\in\mathcal J_x}
               \big([s_v]-\gamma(s_v)v^j[1]\big).
\]
Enlarge the coefficient field to contain the displayed units and
all line characters and use any fixed order of the factors.
Every \(\gamma(H_\gamma)\) is nonzero. Indeed every \(v\) has a
principal part of infinite order, and the exponents are nonzero.
On a line of character \(\beta\), a factor vanishes when
\(\beta(s_v)/\gamma(s_v)=v^j\).

For completeness, the exponents that this construction detects
can be read directly from the jumps. Write \(b_\tau\) and \(a_\tau\)
for the jumps of \(\beta\) and \(\gamma\). The characters are
crystalline and their unit characters are the monomials of
Lemma~\ref{lem:abelianunits744}. Up to the reciprocity sign,
the rational-unit exponent is \(\sum_\tau(b_\tau-a_\tau)\).
The quadratic norm-one exponent is
\(\sum_\tau\epsilon_\tau(b_\tau-a_\tau)\), with the sign chosen
by either embedding of the quadratic completion. Each nonzero
exponent is in \(\mathcal J_x\). At each label a far term, when
present after subtraction, is one signed \(k_\tau\); it is one
of the two base weights minus a shift of at most \(4r_x\).
The remaining scalar and central terms are covered by the
\(6n_xt_0\) allowance in the definition of the list.

For \(\gamma=1\), any far choice is detected. In the nonsplit case
every such choice contributes \(+k_\tau\) to the oriented sum and
the short terms total at most \(2n_xt_0\) in absolute value. In
the split case use the rational sum since the signs at \(w\) agree.
Here \(\min k_\tau>4n_xt_0\) on the specified interval. With
only short jumps a nonzero jump gives positive rational sum.
Thus \(\beta(H_1)\ne0\) precisely on the lines whose jumps are
all zero. Such a line is unramified by its unit monomial. In any
one flag there is at most one such line, since the jump zero
occurs at most once at a fixed embedding. With \(A=B'=H_1\)
in Lemma~\ref{lem:flagoperators756} (where \(B'\) denotes its
right operator), the retained quotient is either zero or that
unramified line. The scaled additive cocycle thus vanishes on
inertia.

For \(\alpha=\kappa^{-h}\), \(1\le h\le B\), the same far-term
test works after subtraction of the short jump \(h\). Any line
with \(\beta(H_\alpha)\ne0\) consequently has only short jumps,
with sum \(h[K:\QQ_\ell]>0\). At least one of its jumps is
positive. In an ordered flag it must follow the line with all
jumps zero if that line exists. For the main target a label at
which \(b_\tau\ne J_\tau\) contributes the opposite far term in
the oriented sum at a nonsplit prime, or in the common-sign
rational sum at a split prime. The same size inequality detects
any such label. A surviving line has \(b_\tau=J_\tau\) at every
label, and again follows the zero line if some \(J_\tau>0\).
Take \(A=H_\alpha\), \(B'=H_1\). In both cases every quotient
of the lemma is zero. A representative satisfying the relations
on the zero family vanishes on each \(g\), giving the asserted
annihilation of the class.

Take \(R=t_0+1\), which bounds every rank. The valuation of
\(\gamma(H_\gamma)\) is the sum of \(v_\ell(v^j-1)\) for the
displayed factors. This is independent of the matrices and of
\(\gamma\). The cost of \(k_{\alpha,w}\) is at most \(2R\) times
that one-product cost. Proposition~\ref{prop:twoweights754}
bounds the weighted sum of the latter over primes, local degrees
and targets by \(O_{q_0,d}(x^{63}b_x^{1/4})\), using the same
list before all suprema. Since \(R\ll x^5\), this proves the
bound. Extending coefficients preserves normalized lengths:
each original cyclic quotient retains the valuation of its
annihilating ideal when \(v_\ell(\ell)=1\).
\end{proof}

\begin{theorem}[Projected periods for families admitting ordered flags]
\leavevmode\label{thm:ordinaryperiod756}
Use the same marked torus data, twists and interval \(x\ge V_0\)
as in Theorem~\ref{thm:projectedcups742}. If each trace system
at every spherical coefficient place admits a potentially ordered
flag as above, on this subfamily one has
\[
            \sum_\ell(\log\ell)\sup_{D\ {\rm at}\ \ell}
                   \frac{p(D)}{\sqrt{b_x}}\ll_{q_0,d}V_0^{-10}.
\]
More generally start with \(r_{\rm fil}(D)\) of
Theorem~\ref{thm:filteredperiod755}. In each target the rational
subquotients supplied by that theorem may be used. Omit also the
weight-zero coordinates where all these representations admit
full \(G_K\)-stable flags, and all nontrivial-target coordinates
where all admit potentially ordered flags. Write \(r_{\rm ord}(D)\)
for the normalized lengths of the resulting joint images. On all
the data the displayed estimate holds with \(p(D)\) replaced
by \((p(D)-r_{\rm ord}(D))_+\).
\end{theorem}
\begin{proof}
At an omitted weight-zero place a multiple of the trace is
unramified by Proposition~\ref{prop:ordinaryunits756}. This
multiple lies in the finite local condition. Indeed unramified
classes for the trivial torsion line are values at Frobenius
and lift to integral unramified classes. At each of the other
places in the proposition its multiple is zero in cohomology.
If every \(J_\tau\) at a place is negative, then \(J_\tau<-1\)
at all its embeddings. The dual main line twisted by one has
strictly positive jumps there and no invariants. The crystalline
dimension formula makes its finite rational cohomology zero,
so the main free cup coordinate has rank zero by
Corollary~\ref{cor:torustests741}. Both original coefficient
primes have already been treated by Corollary~\ref{cor:oldperiod752}.

The exact equality of local classes in
Theorem~\ref{thm:residualinterval755} permits these tests after
passing to its subquotients at each place separately. Every new
coordinate omitted is killed by the corresponding scalar. Its
normalized length is at most the scalar valuation, and at each
place there are at most \([K:\QQ_\ell]\) coordinates per target.
The kernel of projection of the joint global image is contained
in the direct sum of the omitted coordinate images. Hence the
additional loss has weighted prime sum
\(O_{q_0,d}(x^{68}b_x^{1/4})\) per node, uniformly in the flags.
After division by \(\sqrt{b_x}\) the sum over \(x\ge V_0\) is
\(O_{q_0,d}(V_0^{-56})\) since \(b_x\ge x^{500}\).
Adding the bound of Theorem~\ref{thm:filteredperiod755} gives
the general assertion. For the subfamily in the first assertion,
one may use the original representations directly, and every
free coordinate is accounted for.
\end{proof}
The period estimate for unrestricted data still requires the
coordinates for which these local hypotheses have not been proved.
\clearpage

\subsection{The unramified constituent in the local period tests}
\label{sec:unramifiedfactor758}
We examine the absolute irreducible factors of the spherical
trace representations. Only factors whose reduction contains
the trivial character enter the following estimate. At any
one place a given representation has at most one unramified
factor. Its position controls the free local cup coordinates,
even when the extensions joining the other factors are nonsplit.

\begin{lemma}[An integral operator supported on one filtration factor]
\leavevmode\label{lem:unramifiedfactor758}
Let \(K=E_{x,w}\) be a spherical coefficient completion in
Theorem~\ref{thm:projectedcups742}. Allow rational subquotients
of the local trace systems. Enlarge coefficients to a finite
extension over which the rational composition factors are
absolutely irreducible. Assume that each of these factors
whose reduction contains the trivial residual character is
potentially abelian. Then for the given finite family there
is \(Q\in\mathcal O[G_K]\) of augmentation \(\delta_w\) from
Proposition~\ref{prop:potunits746} with the following property.
In a rational \(G_K\)-composition series for each representation,
\(Q\) acts as zero on every factor except possibly a single
unramified character \(\gamma\) with trivial reduction. The
same \(Q\) works for every choice of composition series.
\end{lemma}
\begin{proof}
In the product algebra of the matrices and the augmentation
choose the idempotent for the trivial residual character as
in Theorem~\ref{thm:residualinterval755}. Lift it to an element
\(P_1\) of \(\mathcal O[G_K]\). Its augmentation is one. It
acts as zero on every rational factor \(W\) whose reduction
has no trivial composition factor. Indeed its rank on an
induced lattice in \(W\) equals the dimension of its image
on the reduction. The latter is the multiplicity of the
trivial character, by exactness of the idempotent on vector
spaces. This computation is independent of an induced lattice:
the idempotent rank equals its rank on \(W\) itself.

Let \(H_w\) be the element constructed in
Proposition~\ref{prop:potunits746} and set \(Q=H_wP_1\).
For each remaining \(W\) that proposition puts the image of
\(H_w\) inside \(W^{I_K}\). Its construction applies to \(W\):
the crystalline factors have rank at most \(R=t_0+1\) and
every jump belongs to the original weight multiset. These
are precisely the restrictions used to choose the units for
each inertia orbit in that proposition. Since \(I_K\) is
normal, \(W^{I_K}\) is a \(G_K\)-submodule. Absolute
irreducibility makes it either zero or all of \(W\). In the
second case \(W\) is unramified and is one-dimensional.
Indeed an eigenline of Frobenius is then preserved by its
powers, and by continuity by the entire procyclic quotient.

An unramified line contributes jump zero at every labelled
embedding of \(K\). The original weights are all distinct
at each label. Exactness of the de Rham filtration in the
crystalline category therefore allows at most one unramified
factor in any one representation, counted with multiplicity.
The preceding projection removes this factor too when its
reduction is nontrivial. Finally
\(\epsilon(Q)=\epsilon(H_w)\epsilon(P_1)=\delta_w\).
\end{proof}

\begin{proposition}[Inertia at weight zero and a criterion for every target]
\leavevmode\label{prop:unramifiedsplit758}
Under the assumptions of Lemma~\ref{lem:unramifiedfactor758},
at every precision and at every spherical \(w\) the weight-zero
trace cocycles satisfy
\[
                        \delta_w^{2R}c(I_K)=0 .
\]
Suppose in addition that for each representation the possible
factor \(\gamma\) in that lemma is absent or occurs as a
\(G_K\)-stable subline. For every target \(\alpha\), the
restriction of \(\delta_w^R[c]\) to \(I_K\) is then zero.
For a nontrivial target this implies
\(\lambda_{\alpha,w}\delta_w^R[c]=0\) on \(G_K\), where
\(\lambda_{\alpha,w}=\alpha(s)-1\ne0\) for the inertia unit
test \(s\) of Proposition~\ref{prop:twoweights754}.

The additional subline condition holds whenever, in some
composition series of each member having a factor \(\gamma\),
all jumps of the part below \(\gamma\) are nonnegative at all
embeddings. In particular it holds if all jumps of the ambient
representation are nonnegative at the place.
\end{proposition}
\begin{proof}
Write \(V_j/V_{j-1}=\gamma\) when the possible factor occurs.
Because \(Q\) is zero on the factors of \(V/V_j\) and of
\(V_{j-1}\), its \(R\)-th power sends \(V\) into \(V_j\) and
kills \(V_{j-1}\). With the factor absent it kills all of \(V\).
For \(v\in I_K\), \([v]-1\) sends \(V_j\) into \(V_{j-1}\),
so \(Q^R([v]-1)Q^R\) acts as zero on every original lattice.
Its augmentation is zero. The additive trace cocycle evaluates
this relation as \(\delta_w^{2R}c(v)=0\), by the group ring
product rule of Lemma~\ref{lem:twoblocks754}.

If the factor is a subline, use a composition series beginning
with this subline. It remains the only possible exceptional
factor. Now \(([v]-1)Q^R\) is zero on every representation,
including those with no exceptional factor. Applying the
product rule with target \(\alpha\) gives
\[
          \delta_w^R c(v)=-(\alpha(v)-1)d_c(Q^R).
\]
This is a single coboundary as \(v\) varies through inertia.
Inflation and restriction then put \(\delta_w^R[c]\) in the
image of \(H^1(G_K/I_K,((\mathcal O/\pi^a)(\alpha))^{I_K})\).
For a nontrivial target \(\alpha\), the indicated element \(s\)
acts on the coefficients as \(\alpha(s)\). Its difference from
one kills the invariants, hence also the inflated class.

To prove the subline assertion, let \(U=V_{j-1}\). The short
exact sequence \(0\to U\to V_j\to\gamma\to0\) defines a
finite class in \(H^1_f(K,U\otimes\gamma^{-1})\), since its
middle representation is crystalline. Twisting by the
unramified character preserves all labelled jumps. Under the
stated condition there are no negative jumps, so the de Rham
tangent space of \(U\otimes\gamma^{-1}\) is zero in our jump
convention. Its invariants are zero as well: an invariant would
give a copy of \(\gamma\) in \(U\), repeating the jump zero
from \(V_j/U\). By the Bloch--Kato dimension formula these
two dimensions add to \(\dim H_f^1(K,U\otimes\gamma^{-1})\).
The extension thus splits over the coefficient field, and the
lift of \(\gamma\) in \(V_j\) is the required subline of \(V\).
All group ring identities above were first proved as identities
of rational operators and hence hold exactly in the supplied
lattices.
\end{proof}

\begin{theorem}[The period estimate from the unramified constituent]
\leavevmode\label{thm:unramifiedperiod758}
Use the fields, twists and full interval of
Theorem~\ref{thm:projectedcups742}. For a given target \(M\)
at a spherical place \(w\), use the original local representations
or any of the residual filtration replacements of
Theorem~\ref{thm:residualinterval755}. Omit the weight-zero free
cup coordinates when Lemma~\ref{lem:unramifiedfactor758} applies.
Omit those of any other target when the subline assertion of
Proposition~\ref{prop:unramifiedsplit758} applies. These omissions
may be combined with those of Theorem~\ref{thm:ordinaryperiod756}.
Let \(r_{\rm unr}(D)\) be the normalized length on the coordinates
retained, always using the images of the original joint maps.
Then
\[
 \sum_\ell(\log\ell)\sup_{D\ {\rm at}\ \ell}
       \frac{(p(D)-r_{\rm unr}(D))_+}{\sqrt{b_x}}
                                            \ll_{q_0,d}V_0^{-10}.
\]
In particular, at places where the matrices become triangular
over a finite field extension, the weight-zero tests are covered
in arbitrary flag order. If \(J_\tau>0\) at every embedding
of such a place, the estimate covers every target there.
\end{theorem}
\begin{proof}
First consider a potentially triangular representation. Choose
a finite normal extension \(K'/K\) making it triangular after
enlarging coefficients. Any absolutely irreducible rational
\(G_K\)-factor \(W\) then has a \(G_{K'}\)-stable line, by the
filtration induced on the factor. Its translates under \(G_K\)
span \(W\) and are still \(G_{K'}\)-stable by normality. A
basis can be selected among these lines, in which every
\(G_{K'}\)-matrix on \(W\) is diagonal. Thus all such factors
are potentially abelian. When every \(J_\tau\) is positive,
the remaining jumps are the nonnegative short jumps, giving
the subline condition by the proposition.

For weight zero the scaled class is unramified and is finite
at the torsion precision: its Frobenius value lifts to an
integral unramified class. For other targets the scaled class
vanishes locally. The residual filtration replacements preserve
the class itself in local cohomology, so the same multipliers
annihilate the original free cup outputs. At each place there
are at most \([E_{x,w}:\QQ_\ell]\) coordinates per target.
The kernel of the projection on joint images lies inside the
sum of the omitted coordinate images, each of length bounded
after normalization by the valuation of its annihilator.

Proposition~\ref{prop:potunits746} bounds the weighted sum,
at a fixed node, of the local degrees times \(v_\ell(\delta_w)\)
by \(O_{q_0,d}(x^{80}b_x^{1/4})\), uniformly before the data
are selected at each prime. There are at most \(R=t_0+1\)
targets and each new annihilator uses at most \(2R\) powers
of \(\delta_w\). The \(\lambda_{\alpha,w}\) add at most
\(O_{q_0,d}(x^{63}b_x^{1/4})\) by the direction-unit and
exponent-list bound in Proposition~\ref{prop:twoweights754}.
The total new cost is therefore \(O_{q_0,d}(x^{90}b_x^{1/4})\).
Dividing by \(\sqrt{b_x}\) and summing over the interval gives
\(O_{q_0,d}(V_0^{-34})\), since \(b_x\ge x^{500}\).
Combining with Theorem~\ref{thm:ordinaryperiod756} proves the
claim. If an irreducible factor has nonabelian image over
every finite extension and its reduction contains the trivial
character, the action of \(Q\) on that factor still needs
to be controlled for the unrestricted family.
\end{proof}
\clearpage

\section{Cup bounds and torus periods at the coefficient primes}\label{nav:torus-primes}

\subsection{Summing local cup bounds on the period fields}
\label{sec:sumcupheights763}
Consider the shorter period families of Theorem~\ref{thm:projectedcups742}.
At a node \(x\), keep its field \(E_x\), degree \(n_x=[E_x:\QQ]\),
and all its test data. The local annihilator in
Theorem~\ref{thm:unrestrictedcups762} bounds the order of an
individual cup. We sum the resulting lengths before passing
to the supremum over the interval. This allows a local
weight bound that depends on the node.

For a trace family at a spherical place \(w\mid\ell\) and a
target \(T\), use the three crystalline extensions of
Corollary~\ref{cor:threeextensions761} over \(K=E_{x,w}\).
Write \(h_X,h_{X'},h_Z\) for the respective upper bounds
on the jumps, including zero. Thus \(X\) is the graph
extension, \(X'\) the extension for the integral dual test,
and \(Z=(Y\otimes_{\ZZ_\ell}Y')^\vee(1)\) the kernel of
the third extension. The bounds include every jump after
restriction of coefficients. If the local matrix relations
admit a different representative of each trace class, we
may compute these heights on the replacement family.
Cup evaluation is unchanged by a coboundary.

\begin{lemma}[Prime sum at a single node]
\leavevmode\label{lem:sumcupheight763}
Fix an integer \(H\ge1\). At each datum \(D\) at \(x\)
discard any collection of spherical targets for which
\(\max(h_X,h_{X'},h_Z)\le H\). Let \(\Delta(D)\) be the
drop in the normalized lengths of the joint images. Then
\[
\begin{aligned}
 \sum_\ell (\log\ell)\sup_{D\ {\rm at}\ (x,\ell)}\Delta(D)
 \ \ll_d\ &n_x(t_0+1)H\,
 \bigl((n_x+1)(1+\log(2n_x))^3\\[-2pt]
 &\hspace{38mm}{}+(1+\log(2H))^2\bigr).
\end{aligned}
\]
The estimate also holds for a joint image after integral
scaling of some of its coordinates if the scaled classes
themselves admit the indicated representatives.
\end{lemma}
\begin{proof}
First bound the height constants in
Theorem~\ref{thm:allramifiedlifts761}. At any prime \(p=\ell\)
let \(e=e_K\le n_x\). Choose an integer \(J\) with
\(4e\le 2^J<8e\). Then
\[
 \frac{pe}{(p-1)^2p^J}\le\frac{2e}{2^J}\le\tfrac12,\qquad
 d_K\le 1+\log_2 e,\qquad v_p(h!)\le h
\]
for \(h\ge0\). Substituting in that theorem and rounding up
shows that, for \(0\le h\le H\), both \(C_K(h)\) and the
constant \(Q_K\) of Proposition~\ref{prop:commonlift762}
are \(O(H(1+\log(2n_x))^2)\). This holds at two and in
the presence of wild ramification. There are at most
\(4(n_x+d+1)\) primes up to that integer and their logarithms
are \(O_d(1+\log(2n_x))\). This gives the first term in
parentheses.

For \(p>4(n_x+d+1)\) the field is unramified at \(p\).
Indeed \(E_x=\QQ(\mu_{q_1},\sqrt{-d})\) as in
Proposition~\ref{prop:oddcyclo737}. For a prime \(p\mid q_1\),
the inclusion \(\QQ(\mu_p)\subseteq E_x\) implies
\(p-1\le n_x\). The quadratic factor can ramify only at \(d\).
At any of the remaining primes we have \(e=1,\ d_K=0\).
For \(h<p-1\), \(C_K(h)=0\) and \(v_p(h!)=0\).
For \(p-1\le h\le H\) the choice \(J=0\) in
Theorem~\ref{thm:allramifiedlifts761} gives
\[
 C_K(h)\le
    \left\lceil\frac{p h}{(p-1)^2-p}\right\rceil
       \ll H/p,\qquad v_p(h!)\le H/(p-1).
\]
The implied constant is absolute since \(p\ge 7\) here and
\(H\ge p-1\) in the nonzero case. Thus the sum of local
exponents at large primes is bounded, per possible free
coordinate, by
\[
             O\left(H\sum_{2\le m\le H+1}
                              \frac{\log m}{m}\right)
                 =O\bigl(H(1+\log(2H))^2\bigr).
\]
Integral comparison with \(\int_1^{H+1}\log(2t)\,dt/t\)
gives the last inequality. At both ranges of primes
Corollary~\ref{cor:sphericalcupbound762} bounds the loss by
the sum of the exponents times the local degrees. Their
sum over \(w\mid p\) is \(n_x\), and there are at most
\(t_0+1\) targets at any place. These majorants depend
only on \(x,H,d\) and \(p\), so they bound each supremum.
For the scaled image the proof uses the same injection of
the projection kernel into the discarded coordinates.
\end{proof}

\begin{theorem}[Omitting targets with bounded upper jumps on the interval]
\leavevmode\label{thm:subcriticalperiod763}
Throughout the original interval set
\[
       H_x=\left\lfloor
            \frac{\sqrt{b_x}}{x^{45}(1+\log b_x)^2}
                         \right\rfloor .
\]
At each spherical place omit all targets for which local
representatives of the original classes have
\(\max(h_X,h_{X'},h_Z)\le H_x\), and omit both original
level primes. Let \(r_*(D)\) be the sum of normalized
lengths of the joint global images on the retained
coordinates, with the same finite tests as before. Then
\[
      \sum_\ell(\log\ell)\sup_{D\ {\rm at}\ \ell}
        \frac{(p(D)-r_*(D))_+}{\sqrt{b_x}}
                            \ll_{q_0,d} V_0^{-10}.
\]
The same integral inequality therefore holds with the full
period cost \(p(D)\) wherever the omissions exhaust the
targets of every datum under consideration.
\end{theorem}
\begin{proof}
At the two original level primes use
Corollary~\ref{cor:oldperiod752}. At the spherical places
apply Lemma~\ref{lem:sumcupheight763} at \(H=H_x\).
This is at least one throughout the interval for large
\(V_0\), since \(b_x\ge x^{500}\). Use
\(n_x\ll_{q_0,d}x^{12}\), \(t_0\ll x^5\) and
\(\log(2H_x)\ll1+\log b_x\). After division by \(\sqrt{b_x}\)
the lemma gives a bound at each node of order
\[
 \frac{x^{29}(1+\log x)^3}{x^{45}(1+\log b_x)^2}
            +\frac{x^{17}}{x^{45}}
       \ \ll_{q_0,d}\ x^{-14}.
\]
For nonnegative terms the supremum over the nodes at a fixed
prime is at most their sum. Interchanging these nonnegative
sums gives a total additional loss \(O_{q_0,d}(V_0^{-13})\).
Add it to the estimate of Theorem~\ref{thm:projectedcups742}
and the previously bounded contribution of the level primes.
This uses the same choices of twists, node by node, as in
that theorem; taking a local representative does not
change the global source of a period evaluation.
\end{proof}

\begin{corollary}[Short-weight scalar tests at every spherical place]
\leavevmode\label{cor:shortscalar763}
For any scalar target \(T=\mathcal O(-d_0)\),
\(0\le d_0\le B=t'-1\), its free cup coordinates can be
omitted at every spherical place at which the residual
intervals of Theorem~\ref{thm:residualinterval755} have
only jumps in \([0,B]\). The interval estimate of
Theorem~\ref{thm:subcriticalperiod763} continues to hold.
This includes the trivial target and permits every rank.

For \(1\le d_0\le B\) the same omission is valid when
each trace representation is a rational direct sum of
blocks such that every block of dimension at least three
has only short jumps. The filtrations and decompositions
may depend on the place and target.
\end{corollary}
\begin{proof}
The residual interval theorem supplies representatives of
the same classes with the indicated matrix relations. Use
these representatives in the graph construction. Its kernel
\(Y\) then has only jumps \(j\) with \(0\le j\le B\).
The dual-test kernel \(Y'\) has jump \(-1-d_0\), including
after restriction of coefficients. The jumps of \(Z\) are
\(-1-j-(-1-d_0)=d_0-j\). It follows that
\(h_X\le B,\ h_{X'}=0,\ h_Z\le d_0\le B\).
In fact each such cup is killed by
\(p^{C_K(B)+C_K(d_0)}\): both first extensions have jumps
of one sign, and the upper bounds may be enlarged in the
definition of \(C_K\). Since \(B\ll x^5\) and
\((1+\log b_x)^2\ll b_x^{1/10}\), we have
\(H_x\gg b_x^{2/5}/x^{45}\gg x^5\) at the nodes in question.
The theorem proves the first assertion.

For a positive scalar target in the second assertion,
Proposition~\ref{prop:twoweights754} supplies a nonzero
integral multiplier \(k_{\alpha,w}\) whose scaled class
has a representative supported on the retained blocks.
The proposition removes the lines and every two-dimensional
block containing a far jump. All remaining blocks have
jumps in \([0,B]\) by hypothesis. For an omitted coordinate
with image \(I\), apply the short-weight bound to
\(k_{\alpha,w} I\). Its kernel in \(I\) has normalized
length at most \(v_\ell(k_{\alpha,w})\) per free coordinate.
For a joint image this follows by applying coordinate
multiplication first on the omitted coordinates and the
identity on all retained coordinates, then projecting.
The additional weighted cost at a node is at most
\(O_{q_0,d}(x^{63}b_x^{1/4})\) by the same proposition.
After division by \(\sqrt{b_x}\), summation over the
original nodes costs \(O_{q_0,d}(V_0^{-61})\).
Combining the two losses proves the assertion.
\end{proof}

In the unrestricted families a block of rank at least
three can contain a main jump \(J_\tau\) outside the
displayed range. The local theorem still bounds its cup
order. The prime sum proved here controls only the targets
meeting the hypotheses above.
\clearpage

\subsection{Large coefficient primes with a multiple near a period weight}
\label{sec:largeweightprimes766}
We return to the full marked torus data of
Theorem~\ref{thm:projectedcups742}, including trace
representations with a far jump. In local lifting arguments
one encounters the condition that a coefficient prime
divide a nonzero integer close to a large weight. The
distance allowed can itself be of order the weight divided
by the prime. For sufficiently large primes we can sum
all these exceptions on the original interval, with no
restriction on the local representations. The dependence
of the local cup bound on the prime is needed here.

\begin{lemma}[Counting prime divisors near finitely many weights]
\leavevmode\label{lem:nearweightprimes766}
Fix a node \(x\). Write \(b=b_x,\ n=n_x\) and take a finite
nonempty set \(\mathcal C\) of integer centers with
\(|c|\le M\) for every \(c\in\mathcal C\). Let
\(P>4(n+d+1)\), \(W\ge0\) and \(L\ge0\). Consider the
primes \(\ell\ge P\) for which some \(c\in\mathcal C\)
and some nonzero multiple \(v\) of \(\ell\) satisfy
\[
                       |v-c|\le W+Lb/\ell .
\]
With \(R=\lceil W+Lb/P\rceil\), their total contribution obeys
\[
 \sum_{\ell\ {\rm as\ above}}(\log\ell)
                \sup_{D\ {\rm at}\ (x,\ell)} r(D)
 \ll n(t_0+1)\,\frac bP\,
          |\mathcal C|(2R+1)\log(2+M+R).
\]
\end{lemma}
\begin{proof}
The two original level primes have zero contribution by
Corollary~\ref{cor:oldperiod752}. Every other coefficient
prime is spherical in the marked families under consideration.
For \(\ell\ge P\), all the completions \(E_{x,w}\) are
unramified by the field description in
Lemma~\ref{lem:sumcupheight763}. At every embedding, every
jump of the graph kernel is either a short jump in
\([0,t'-1]\) or a main jump \(J_\tau=z-1\pm k_\tau\),
where \(k_\tau\le4b\). The target is one of the same
characters and its dual has jumps \(-1-j_T\).
Consequently the three upper jumps in
Corollary~\ref{cor:sphericalcupbound762} are at most
\(12b\) for all sufficiently large \(x\), also after
restriction of coefficients. For the third kernel this
uses the sum of the two largest absolute jumps and one,
so it includes all pairs of coefficient embeddings.

As in the large-prime calculation of
Lemma~\ref{lem:sumcupheight763}, at an unramified completion
the local exponent is zero if every upper jump is less
than \(\ell-1\), and otherwise is \(O(b/\ell)\).
Indeed \(d_K=0\), \(v_\ell(h!)\le h/(\ell-1)\), and
the ceiling in \(C_K(h)\) only occurs for \(h\ge\ell-1\).
Applying the cup bound to all targets gives
\(r(D)\ll n(t_0+1)b/\ell\). This bound also holds for
zero lifts and is uniform in the local systems and all the
integral tests.

Each prime under consideration divides the positive integer
\[
          \prod_{c\in\mathcal C}
          \ \prod_{\substack{j\in\ZZ,\ |j-c|\le R\\j\ne0}}|j|,
\]
using \(1\) for an empty product. Its logarithm is at most
\(|\mathcal C|(2R+1)\log(2+M+R)\). The sum of the logarithms
of the distinct prime divisors is at most the logarithm
of the product. Multiplying by the uniform bound with
\(\ell\ge P\) proves the lemma.
\end{proof}

\begin{theorem}[The full period on the large primes near weight combinations]
\leavevmode\label{thm:exceptionalperiod766}
At each node let \(b_1=b_x,\ b_2=a_x/x_{1,x}\) be the two
integral base weights, let \(U_x=n_x(2t_0+1)\), and put
\[
 \begin{aligned}
   \mathcal C_x&=\{d_1b_1+d_2b_2:
                d_1,d_2\in\ZZ,\ |d_1|,|d_2|\le U_x\},\\
   P_x&=x^{50}(1+\log(2b_x))\, b_x^{3/4}.
 \end{aligned}
\]
Fix \(L\ge1\). Let \(\mathcal E_x\) be the coefficient primes
\(\ell\ge P_x\) with a nonzero multiple at distance at most
\(L(x^{125}+b_x/\ell)\) from some element of \(\mathcal C_x\).
On the original fields, for the original tests and interval,
\[
      \sum_\ell(\log\ell)
         \sup_{\substack{D\ {\rm at}\ \ell\\\ell\in\mathcal E_x}}
         \frac{p(D)}{\sqrt{b_x}}\ \ll_{q_0,d,L}\ V_0^{-10}.
\]
This controls all local cup coordinates on the indicated
data, including those of the main line and the weight-zero
target, for arbitrary ranks of the crystalline trace matrices.
\end{theorem}
\begin{proof}
Use the degree and rank bounds \(n_x\ll_{q_0,d}x^{12}\),
\(t_0\ll x^5\). Then \(U_x\ll x^{17}\),
\(|\mathcal C_x|\ll x^{34}\), and all centers have absolute
value at most \(2U_x b_x\). By the scale conditions
\(x^{10^6}\le x_{1,x}\), \(x_{1,x}^{10^6}\le b_x\) in
Theorem~\ref{thm:negativeinterval740}, \(P_x>4(n_x+d+1)\)
on a sufficiently distant interval. Also
\(\log(2+2U_xb_x+R)\ll_{q_0,d,L}1+\log(2b_x)\) for the
radius \(R=\lceil L(x^{125}+b_x/P_x)\rceil\).
Lemma~\ref{lem:nearweightprimes766} bounds the node sum of
\(r(D)/\sqrt{b_x}\) at the specified primes by
\[
 \ll_{q_0,d,L}
 \frac{x^{51}(1+\log(2b_x))}{\sqrt{b_x}}
     \left((1+x^{125})\frac{b_x}{P_x}
                              +\left(\frac{b_x}{P_x}\right)^2\right).
\]
After substitution the first term is
\(O_{q_0,d,L}(x^{126} b_x^{-1/4})\), and the second is
\(O_{q_0,d,L}(x^{-49}/(1+\log(2b_x)))\).
Their sum is \(O_{q_0,d,L}(x^{-48})\).
The use of a single product per node has included all
permitted signs, powers and coefficient embeddings before
taking the supremum. Sum over \(x\ge V_0\) to bound the
prime sum of the suprema of \(r(D)/\sqrt{b_x}\) on the
specified data by \(O_{q_0,d,L}(V_0^{-47})\).
Finally \(p(D)\le (p(D)-r(D))_++r(D)\); the prime sum of
the first term is bounded by Theorem~\ref{thm:projectedcups742}.
This proves the claimed estimate with the full cost \(p(D)\).
\end{proof}

For example, the original shifts from a base weight are
\((2j+1)r_x^8-h_\tau\), with \(j\le2\sqrt{x}\) and
\(0\le h_\tau\le4r_x\). The centering \(z-1\) has size
at most \(t_0\). Even after summing \(2U_x\) such errors,
the result is \(O_{q_0,d}(x^{59})\). Thus the theorem
includes the large primes with a nonzero multiple within
\(O(b_x/\ell+r_x^{24})\) of a signed sum of at most
\(2U_x\) actual jumps of the displayed types. Indeed
\(r_x^{24}\ll x^{120}\), and the fixed constants can be
absorbed in \(L\). Here the numbers of
occurrences of each base weight, counted with sign, must
give coefficients in \([-U_x,U_x]\). The remaining prime
sum includes large primes away from these neighborhoods
and primes below \(P_x\).
\clearpage

\subsection{The full marked torus period above the square-root cutoff}
\label{sec:fullperiodcutoff769}
Write \(n_x=[E_x:\QQ]\), \(t_0=5r_x\), and use \(p(D)\) and
\(r(D)\) as in Theorem~\ref{thm:projectedcups742}. The latter
is the normalized length of the \emph{joint} image on the
original free local cup coordinates. The difference between
the full period cost and the part already estimated there is
bounded by \(r(D)\). We first bound its prime sum at one node.

\begin{proposition}[A weighted sum of the joint local image]
\leavevmode\label{prop:weightedlocaltail769}
Fix \(q_0,d\). At every sufficiently large node \(x\) of the
shorter marked torus family, for every real \(L\ge x_{1,x}\),
\[
 \sum_{\ell\ge L}(\log\ell)
           \sup_{D\ {\rm at}\ (x,\ell)}r(D)
 \ \ll_{q_0,d}\
 n_x(t_0+1)(x^{42}+b_x/L)(1+\log(2b_x)).
\]
The supremum at each prime ranges over the original integral
tests, separating systems, coefficient embeddings and signed
powers at that node; an empty contribution is zero.
\end{proposition}
\begin{proof}
Put \(b=b_x\). The two original level primes contribute
zero by Corollary~\ref{cor:oldperiod752}. Every other
coefficient place is spherical in this family. At a place
\(w\mid\ell\), Corollary~\ref{cor:actualperiodmultiples768}
annihilates each free cup evaluation by \(\ell^{2S_{x,\ell}}\).
For a cyclic coordinate with coefficients \(\mathcal O/\pi^m\),
the subgroup killed by \(\ell^h\) has normalized
\(\mathcal O\)-length at most \(h\): its ordinary length is
\(\min(m,h e_{\mathcal O})\), where
\(\ell\mathcal O=\pi^{e_{\mathcal O}}\mathcal O\).
There are at most \([E_{x,w}:\QQ_\ell]\) coordinates per
target at this place, and at most \(t_0+1\) targets.
The joint global image is a submodule of the product of
these killed subgroups. Additivity of length in a product
and \(\sum_{w\mid\ell}[E_{x,w}:\QQ_\ell]=n_x\) give
\[
        r(D)\le 2n_x(t_0+1) S_{x,\ell}.
\]
This argument keeps the given global source, including any
relations between its local projections.

Take \(R=\lceil C_0(x^{42}+b/L)\rceil\), with \(C_0\) from
Corollary~\ref{cor:actualperiodmultiples768}. Let \(\mathcal N\)
consist of the nonzero integers within distance \(R\) of
at least one of \(0,\pm b,\pm a_x/x_{1,x}\). For every
\(\ell\ge L\) it contains the full set used to define
\(S_{x,\ell}\). Hence unique factorization gives
\[
 \sum_{\ell\ge L}(\log\ell) S_{x,\ell}
  \le \sum_{n\in\mathcal N}\sum_{\ell\mid n}
           v_\ell(|n|)\log\ell
   =\sum_{n\in\mathcal N}\log|n|.
\]
The number of integers is at most \(5(2R+1)\), and their
absolute values are at most \(b+R\). Since \(x^{42}\le b\),
the last sum is \(O_{q_0,d}((x^{42}+b/L)(1+\log(2b)))\).
This proves the proposition. Multiplicities of prime factors
are included in the factorization, as required for the
annihilators over the original finite rings.
\end{proof}

\begin{theorem}[The full period above a node-dependent cutoff]
\leavevmode\label{thm:fullperiodcutoff769}
For each node let
\[
          L_x=x^{50}(1+\log(2b_x))\sqrt{b_x}.
\]
On the same fields \(E_x\), on the same interval starting
at \(V_0\), and with the twists of
Theorem~\ref{thm:projectedcups742}, one has
\[
  \sum_\ell(\log\ell)
       \sup_{\substack{D\ {\rm at}\ \ell\\ \ell\ge L_x}}
                  \frac{p(D)}{\sqrt{b_x}}
       \ll_{q_0,d} V_0^{-10}.
\]
This holds for all sufficiently large \(V_0\). Every
coefficient prime above the displayed threshold is included,
regardless of its residue relative to the period weights
or of the ranks of the crystalline traces.
\end{theorem}
\begin{proof}
The inequality \(x_{1,x}^{10^6}\le b_x\) implies
\(L_x\ge x_{1,x}\). Dividing the proposition by \(\sqrt{b_x}\)
and using \(n_x\ll_{q_0,d}x^{12}\), \(t_0+1\ll x^5\),
we bound the node sum by
\[
 O_{q_0,d}\!\left(
    x^{59}\frac{1+\log(2b_x)}{\sqrt{b_x}}+x^{-33}\right).
\]
For large \(b_x\), \(1+\log(2b_x)\le b_x^{1/4}\).
Already \(b_x\ge x^{400}\) makes the first term
\(O(x^{-41})\). The sum over integer nodes \(x\ge V_0\)
is \(O_{q_0,d}(V_0^{-32})\). For each \(\ell\) the
supremum over admissible nodes of the nonnegative normalized
joint length is bounded by the sum of the corresponding
node suprema. Summing over \(\ell\) therefore gives the
same upper bound on the restricted prime sum of
\(r(D)/\sqrt{b_x}\). Finally
\[
             p(D)\le \max\{0,p(D)-r(D)\}+r(D).
\]
Theorem~\ref{thm:projectedcups742} bounds the prime sum
of the normalized first term, including at these primes,
by \(O_{q_0,d}(V_0^{-10})\). The two estimates use the
same choices of data and give the stated inequality.
\end{proof}

\begin{corollary}[The local contribution from the first auxiliary scale]
\leavevmode\label{cor:firstauxiliarytail769}
At a single node the free cup contribution on the longer
range \(\ell\ge x_{1,x}\) satisfies
\[
 \sum_{\ell\ge x_{1,x}}(\log\ell)
         \sup_{D\ {\rm at}\ (x,\ell)}r(D)
   \ll_{q_0,d} n_x b_x x^{-20}(1+\log(2b_x)).
\]
\end{corollary}
\begin{proof}
Apply the proposition with \(L=x_{1,x}\). Since
\(t_0+1\ll x^5\), \(x_{1,x}\ge x^{10^6}\) and
\(b_x\ge x^{400}\), both
\((t_0+1)x^{42}\) and \((t_0+1)b_x/x_{1,x}\)
are \(O(b_x x^{-20})\). This proves the inequality.
\end{proof}
The theorem bounds the full period at the square-root
normalization on its stated range. The corollary uses a
larger normalization on its wider range. Controlling the
remaining coefficient primes at the required scale and
comparing the integral periods of the higher-rank lifts
are separate steps in the passage to Dirichlet GRH.
\clearpage

\subsection{The remaining main cup coordinates at intermediate primes}
\label{sec:remainingmain771}
At a node \(x\) let \(b=b_x\), \(b_2=xb/x_{1,x}\), and call
these the two base weights of the marked data. At a coefficient
place we consider all labelled embeddings there for the
given datum, including the labels obtained after extending
coefficients. The sign at a label is the sign of its main
jump \(J_i\). The short shifts in the jump do not change
the sign of the signed base weight for sufficiently large \(x\).

\begin{lemma}[Main targets at places with constant sign or base weight]
\leavevmode\label{lem:onesignmain771}
For \(p\ge x^{50}\), let \(w\mid p\) be spherical. If the
sign is constant over \(w\), or if the base weight is constant
over \(w\), then \(p^{2S_{x,p}}\) kills each cup of the main
target, with the five intervals of
Corollary~\ref{cor:allscalartests770}. In particular the
sign condition holds at every such place for a rational prime
split in \(\QQ(\sqrt{-d})\).
\end{lemma}
\begin{proof}
All these fields are unramified at \(p\) by
Lemma~\ref{lem:sumcupheight763}, and \(B=t'-1<p-1\).
First suppose the sign is constant. Use
\(a=\lceil12b/(p-3)\rceil\) to lift \(E^a\) in all
three extensions, as in Corollary~\ref{cor:actualperiodmultiples768}.
At a positive place the graph has only nonnegative jumps,
and the main dual has only nonpositive jumps. The first
therefore has the required gap, and in the second the
unit already lifts by Theorem~\ref{thm:integrallifts759}.
At a negative place the graph has upper jump at most \(B\)
at every embedding, so its unit lifts by the same theorem
since \(B<p-1\). The main dual has only nonnegative jumps
and satisfies the gap condition. These two cases meet the
respective labelled alternatives of
Theorem~\ref{thm:twolatticecup768}. Its costs lie in the
five intervals by the shift bounds. At a rational prime
split in the quadratic field, all embeddings of a fixed
completion induce the same embedding of that field. The
CM type described in Theorem~\ref{thm:negativeinterval740}
then gives the same sign.

If the base weight is a fixed \(\beta\in\{b,b_2\}\) at the
place, every main jump has absolute value between
\(\beta-O_{q_0,d}(x^{41})\) and
\(\beta+O_{q_0,d}(x^{41})\). We have
\(\beta\ge b^{1-10^{-6}}\). In all three extensions every
upper jump is bounded by \(12\beta\), also for the third
kernel since its jumps are differences at the same label.
Use \(a=\lceil12\beta/(p-3)\rceil\) in the lifting
argument. Now \(|J_i|>a+B+2\) at all labels for large
\(x\), without requiring the sign to be constant.
The graph and main dual both have the gap, and the
two labelled alternatives hold as in the proof of
Corollary~\ref{cor:actualperiodmultiples768}.
The radius needed for each valuation sum is
\(O_{q_0,d}(x^{42}+\beta/p)\) about the appropriate centers
and is bounded by the radius in the assertion.
\end{proof}

\begin{lemma}[Places having only the smaller base weight]
\leavevmode\label{lem:lightplaces771}
At all spherical places of a datum where every base weight
is \(b_2\), omit all targets. Let \(\Delta_2(D)\) be the
resulting drop in joint normalized length. Then
\[
 \sum_\ell(\log\ell)\sup_{D\ {\rm at}\ (x,\ell)}\Delta_2(D)
 \ll_{q_0,d} n_x b_x x^{-45}(1+\log(2b_x)).
\]
This includes ramified places and places above \(2\).
\end{lemma}
\begin{proof}
All three upper jumps are at most \(H=12b_2\) at each
omitted place, also in restriction of scalars.
For primes up to \(4(n_x+d+1)\) the local estimates in
Lemma~\ref{lem:sumcupheight763} bound the sum by
\[
 O_d\big(n_x(t_0+1)H(n_x+1)(1+\log(2n_x))^3\big).
\]
For larger primes its proof bounds each exponent by
\(O(H/\ell)\) when \(\ell\le H+1\), and by zero thereafter.
Put \(N=\lfloor H+1\rfloor\). For every prime \(\ell\le N\)
we have \(v_\ell(N!)\ge\lfloor N/\ell\rfloor\ge N/(2\ell)\).
Consequently
\[
       \sum_{\ell\le N}\frac{\log\ell}{\ell}
          \le \frac{2\log(N!)}{N}\le 2\log N.
\]
Thus the larger primes together cost
\(O(n_x(t_0+1)H(1+\log(2H)))\). In both sums we have
used the original joint length, as in the cited lemma,
and the bounds hold after the supremum at the node.
Since \(b_2=xb_x/x_{1,x}\), \(x_{1,x}\ge x^{10^6}\),
\(n_x\ll_{q_0,d}x^{12}\) and \(t_0+1\ll x^5\), each of
these two bounds satisfies the asserted inequality.
\end{proof}

\begin{theorem}[A period bound after retaining the mixed main coordinates]
\leavevmode\label{thm:remainingmain771}
Use the interval and twists of Theorem~\ref{thm:projectedcups742}.
For each datum project its original joint image onto the
coordinates left after the following omissions.
\begin{enumerate}
\item Omit both original level primes. At spherical places
with \(\ell\ge x_{1,x}\), omit all targets.
\item At spherical places with \(\ell\ge x^{50}\), omit
every scalar target. Also omit the main target if either
the sign or the base weight is constant across the labels.
\item At a spherical place with base weight \(b_2\) at
every label, omit all targets at every prime.
\end{enumerate}
Write \(r_{\mathrm{rem}}(D)\) for the normalized length of
this projection. Set \(f_x=n_x/2\) and \(k_x=r_x+1-m_x\)
for the ranks of the shorter family. Then
\[
 \sum_\ell(\log\ell)\sup_{D\ {\rm at}\ \ell}
   \frac{\max\{0,p(D)-r_{\mathrm{rem}}(D)\}}
        {f_x(m_x/k_x)b_x(1+\log(2b_x))}
     \ll_{q_0,d} V_0^{-39}.
\]
\end{theorem}
\begin{proof}
The projection kernel embeds in the product of the
discarded local coordinates. The two original level primes
have zero free image. At primes \(\ell\ge x_{1,x}\),
Corollary~\ref{cor:actualperiodmultiples768} bounds every
free evaluation by the exponent \(2S_{x,\ell}\). At the
other primes \(\ell\ge x^{50}\) the coordinates specified
in (2) have exponents at most \(3S_{x,\ell}\) by
Corollary~\ref{cor:allscalartests770} and
Lemma~\ref{lem:onesignmain771}. This argument uses the
same \(C_0\) throughout, enlarged once for the displayed
families. The proof of Proposition~\ref{prop:weightedlocaltail769}
applies to these discarded coordinates with \(L=x^{50}\):
its factorization only requires containment in the fixed
five intervals of radius \(\lceil C_0(x^{42}+b_x/L)\rceil\).
Thus their weighted cost at the node is
\[
 O_{q_0,d}\big(n_x(t_0+1)(x^{42}+b_x/x^{50})
                              (1+\log(2b_x))\big).
\]
Here \(b_x\ge x^{400}\). With \(t_0+1\ll x^5\) this is
\(O_{q_0,d}(n_x b_x x^{-45}(1+\log(2b_x)))\).
The remaining omissions in (3) satisfy this same bound
by Lemma~\ref{lem:lightplaces771}. More explicitly, one
can project first away from the coordinates in (3).
Its kernel has the bound of that lemma, and the kernel
of the next projection injects into the product of the
other discarded coordinates. This proves the node bound
for \(r(D)-r_{\mathrm{rem}}(D)\).

Since \(m_x\ge1\) and \(0<k_x\le r_x+1\ll x^5\), division
by the denominator in the theorem gives \(O_{q_0,d}(x^{-40})\)
for this node sum. Passing to the supremum over the interval
costs at most the sum over all \(x\ge V_0\), as in
Theorem~\ref{thm:fullperiodcutoff769}, hence
\(O_{q_0,d}(V_0^{-39})\). Finally bound the positive part
of \(p(D)-r_{\mathrm{rem}}(D)\) by the positive part of
\(p(D)-r(D)\) plus the drop in joint length. For the first
term use Theorem~\ref{thm:projectedcups742}: the ratio of
\(\sqrt{b_x}\) to the stated denominator is
\(O_{q_0,d}(x^5/\sqrt{b_x})=O_{q_0,d}(V_0^{-195})\).
This completes the interval estimate.
\end{proof}
In particular, for \(x^{50}\le\ell<x_{1,x}\) the retained
coordinates can occur only for main targets at places with
both signs and both base weights. At smaller primes this
argument still retains some scalar targets as well.
\clearpage

\subsection{A shorter range for the mixed main cups}
\label{sec:mixedrange773}
The order at the second Frobenius prime in
Theorem~\ref{thm:mixedcup772} costs at most a constant times
the smaller base weight at each coefficient prime. Summing
this cost only up to a prescribed point lets us shorten
the range left in Theorem~\ref{thm:remainingmain771}.
Near its other end, the ratio \(b_x/b_2=x_{1,x}/x\) is the
threshold that suffices for the first Taylor expansion.

\begin{proposition}[The sum of the main coordinates up to \(Z\)]
\leavevmode\label{prop:mixedrange773}
At a sufficiently large node let \(x^{50}\le Z\le x_{1,x}\).
At spherical places above \(x^{50}\le\ell\le Z\) discard
all main targets, and let \(\Delta_{\rm main}(D)\) denote
the drop in normalized length of the original joint image
or of any projection of that image onto coordinates.
With \(b=b_x\) and \(b_2=xb/x_{1,x}\), the resulting sum obeys
\[
 \begin{aligned}
 \sum_{x^{50}\le\ell\le Z}
       (\log\ell)\sup_{D\ {\rm at}\ (x,\ell)}\Delta_{\rm main}(D)
 \ll_{q_0,d} n_x\big(& (x^{42}+b/x^{50})(1+\log(2b))\\
                     &+Z(b_2+x^{42}+1)+Z\log(2b)\big).
 \end{aligned}
\]
At all primes \(\ell\ge64x_{1,x}/x\), the annihilator
\(\ell^{2S_{x,\ell}}\) of
Corollary~\ref{cor:actualperiodmultiples768} already holds
at every spherical place, for every target.
\end{proposition}
\begin{proof}
There is one main target per place. Its number of free
coordinates is at most the local degree, whose sum over
places is \(n_x\). The projection kernel embeds into
these coordinates; by Theorem~\ref{thm:mixedcup772} each
has normalized length bounded by the exponent in that
theorem. The part \(2S_{x,\ell}\), summed with \(\log\ell\),
is handled by the integer factorization in
Proposition~\ref{prop:weightedlocaltail769} with lower
cutoff \(x^{50}\). This gives the first term in the bound.

For the term constant in \(\ell\), we recall the elementary
inequality \(\sum_{\ell\le Z}\log\ell\ll Z\). Each prime
\(n<\ell\le2n\) divides \(\binom{2n}{n}\), so the sum of
its logarithm over that interval is at most \(2n\log2\).
Summing on dyadic intervals up to the first power of two
at least \(Z\) proves the inequality. The term
\(\log(2b)/\log\ell\) costs at most \(Z\log(2b)\).
These estimates also hold for suprema in the data because
their bounds depend only on \(x\) and \(\ell\).

For the final assertion, take
\(a=\lceil12b/(\ell-3)\rceil\). If \(\ell\ge64x_{1,x}/x\),
then \(\ell\ge x^{50}\) for large \(x\), and
\(a\le16b/\ell+1\le b_2/4+1\).
At every label \(|J_i|\ge b_2-O_{q_0,d}(x^{41})\) and
\(B=t'-1\ll x^5\). Since \(b_2\ge b^{1-10^{-6}}\), the
gap \(|J_i|>a+B+2\) holds everywhere. This is the only
weight separation needed in the proof of
Corollary~\ref{cor:actualperiodmultiples768}, after
unramifiedness and \(B<\ell-1\). Applying the same
two-lattice theorem gives the claimed annihilator.
\end{proof}

\begin{theorem}[The shortened interval of retained main coordinates]
\leavevmode\label{thm:middlerange773}
Fix \(8\le\gamma\le46\). At each node, begin with the
retained coordinates of Theorem~\ref{thm:remainingmain771}.
Among them at primes \(\ell\ge x^{50}\), retain only
coordinates at primes in
\[
                 \frac{x_{1,x}}{x^\gamma}<\ell<
                         \frac{64x_{1,x}}{x}.
\]
These are main targets at places having both signs and both
base weights. Let \(r_\gamma(D)\) be the normalized length
of the joint projection onto the resulting coordinates,
including those retained below \(x^{50}\). For the same
interval and twists as before,
\[
 \sum_{\ell}(\log\ell)\sup_{D\ {\rm at}\ \ell}
 \frac{\max\{0,p(D)-r_\gamma(D)\}}
      {f_x(m_x/k_x)b_x(1+\log(2b_x))}
 \ll_{q_0,d} V_0^{-39}+V_0^{\,7-\gamma}.
\]
\end{theorem}
\begin{proof}
Write \(U=x_{1,x}/x^\gamma\), so \(x^{50}\le U\) for all
large \(x\) in the stated range of \(\gamma\).
At a node the additional drop for primes through \(U\)
is controlled by the proposition at \(Z=U\). The
term \(U b_2\) equals \(b_x x^{1-\gamma}\). The other
terms, apart from this one, total
\[
        O_{q_0,d}\big(n_x b_x x^{-50}(1+\log(2b_x))\big).
\]
Indeed \(x^{42}\ll b_x/x^{50}\) and
\(U(x^{42}+1)\ll b_x/x^{50}\), since
\(U\le x_{1,x}\le b_x^{10^{-6}}\) and \(b_x\ge x^{400}\).
Also \(U\le b_x/x^{50}\). At primes at least
\(64x_{1,x}/x\) the additional main coordinates cost
at most \(2S_{x,\ell}\) each and there are at most \(n_x\)
per prime. The same factorization bound, even using the
smaller lower cutoff \(x^{50}\), gives the same last
error. These estimates concern projection kernels and
apply after the earlier omissions.

Divide by the denominator of the theorem and use
\(n_x=2f_x\) and \(k_x/m_x\ll x^5\).
The extra drop at the node is
\(O_{q_0,d}(x^{-45}+x^{6-\gamma})\) after division.
The sum of these bounds for integer \(x\ge V_0\) is
\(O_{q_0,d}(V_0^{-44}+V_0^{7-\gamma})\), with a uniform
constant since \(\gamma\ge8\). The supremum at each prime
is bounded by this summation of its nonnegative node
contributions. Adding Theorem~\ref{thm:remainingmain771}
gives the claimed estimate.
\end{proof}

\begin{corollary}[Size of the residual contribution on the shortened interval]
\leavevmode\label{cor:coarsemixed773}
At one node the weighted sum of the retained main lengths
on the displayed interval is
\[
       O_{q_0,d}\big(n_x b_x\{1+x^{-50}(1+\log(2b_x))\}\big).
\]
\end{corollary}
\begin{proof}
Use the proposition at \(Z=64x_{1,x}/x\), which lies
between \(x^{50}\) and \(x_{1,x}\) for large \(x\).
Now \(Z b_2=64 b_x\) and the other terms have the bounds
used in the preceding proof. For a datum at such a prime
the retained image has only main coordinates; omitting
all of them drops precisely its length.
\end{proof}
The term \(n_x b_x\) in the corollary still requires a
sharper estimate to obtain the stated interval normalization
for the full period. There are also retained coordinates
below the polynomial cutoff and the higher-rank integral
theta-period comparison.
\clearpage

\subsection{The full torus period at large coefficient primes}
\label{sec:fulltoruslarge777}
The local estimate of Theorem~\ref{thm:mainorbit776}
allows us to sum over the entire range in which
Theorem~\ref{thm:middlerange773} retained mixed main
coordinates. We use the original field
\(E_x=\QQ(\mu_{q_{1,x}},\sqrt{-d})\) at each node, its
degree \(n_x\), and the shorter torus data and projected
integral tests of Theorem~\ref{thm:projectedcups742}.

\begin{proposition}[The sum of the mixed main lengths]
\leavevmode\label{prop:sumallmixed777}
At a sufficiently large node put \(L=x^{50}\) and
\(Z=x_{1,x}\). For any projection of the joint local
image, write \(\Delta(D)\) for its drop in normalized
length upon further omitting any collection of the main
coordinates at spherical places. Then
\[
 \sum_{L\le\ell\le Z}(\log\ell)
     \sup_{D\ {\rm at}\ (x,\ell)}\Delta(D)
   \ll_{q_0,d} n_x b_x x^{-50}(1+\log(2b_x)).
\]
\end{proposition}
\begin{proof}
Write \(b=b_x,\ G=xb/Z\), and use \(W=O_{q_0,d}(x^{42})\)
as in Theorem~\ref{thm:mainorbit776}. At a place \(w\)
there is one main target with at most
\([E_{x,w}:\QQ_\ell]\) free coordinates. An \(\ell^h\)-killed
submodule in one copy of \(\mathcal O/\pi^m\) has length
at most \(h\,e_{\mathcal O}\), where \(v_\pi(\ell)=e_{\mathcal O}\).
The projection kernel injects into the product of the
omitted coordinates. Summing the local degrees gives
\(n_x\), so the theorem bounds \(\Delta(D)\), before or
after the indicated supremum, by
\[
 C_{q_0,d} n_x
 \left(W+1+\frac{b}{\ell^2}+\frac{G}{\ell}
                         +\frac{\log(2b)}{\log\ell}\right).
\]
The elementary bound \(\sum_{\ell\le y}\log\ell\ll y\)
was proved in Proposition~\ref{prop:mixedrange773}. On
\([2^j L,2^{j+1}L)\) it gives
\(\sum(\log\ell)/\ell^2\ll (2^jL)^{-1}\), so the sum
for all \(\ell\ge L\) is \(O(1/L)\). On each dyadic
interval below \(2Z\) it also gives
\(\sum(\log\ell)/\ell\ll1\). Consequently, after taking
out the factor \(n_x\), the total is bounded by a constant
times
\[
          (W+1)Z+b/L+G(1+\log(2Z))+Z\log(2b).
\]
Here \(x^{10^6}\le Z\le b^{10^{-6}}\) and \(b\ge x^{400}\).
In particular \(x^{92}Z\le b\) for sufficiently large
\(x\), by \(b^{1-10^{-6}}\ge x^{400(1-10^{-6})}\).
It follows that \((W+1)Z\ll b/L\) and \(Z\le b/L\).
Also \(G\le b x^{1-10^6}\le b/L\) and \(Z\le b\).
Substitution gives the assertion.
\end{proof}

\begin{theorem}[The marked torus period with only small-prime coordinates retained]
\leavevmode\label{thm:smallcoords777}
Use the fields, interval starting at \(V_0\), and
simultaneous twists of Theorem~\ref{thm:projectedcups742}.
At primes \(\ell<x^{50}\) keep the projection of
Theorem~\ref{thm:remainingmain771}; at all other primes
take the zero projection. Write \(r_{<}(D)\) for its
normalized length. Set
\[
     \mathcal B_x=f_x(m_x/k_x)b_x(1+\log(2b_x)),
                \qquad f_x=n_x/2 .
\]
For all sufficiently large \(V_0\) one has
\[
 \sum_\ell(\log\ell)\sup_{D\ {\rm at}\ \ell}
          \frac{(p(D)-r_{<}(D))_+}{\mathcal B_x}
             \ll_{q_0,d} V_0^{-39}.
\]
In particular, with the supremum restricted to data with
\(\ell\ge x^{50}\), the same bound holds for
\(p(D)/\mathcal B_x\) itself at the original integral tests.
\end{theorem}
\begin{proof}
Write \(r_{\rm rem}\) for the earlier projection. It
already omits the original level primes, and at
\(\ell\ge Z=x_{1,x}\) it contains no coordinates.
For \(x^{50}\le\ell<Z\), the only coordinates still
present come from main targets at spherical places.
Thus for fixed \(x\) the last proposition bounds the sum
of \((\log\ell)\sup(r_{\rm rem}-r_{<})\) at that node.
At \(\ell<x^{50}\) the difference is zero.
Dividing by \(\mathcal B_x\), using \(n_x=2f_x\) and
\(k_x/m_x\ll x^5\), gives \(O_{q_0,d}(x^{-45})\).
For each prime, the supremum over the interval of these
nonnegative quantities is bounded by the sum of their
node suprema. Summing in either order costs at most
\(O_{q_0,d}(\sum_{x\ge V_0}x^{-45})=O_{q_0,d}(V_0^{-44})\).
Finally
\[
 (p(D)-r_{<}(D))_+
 \le (p(D)-r_{\rm rem}(D))_++r_{\rm rem}(D)-r_{<}(D).
\]
The first term has the bound of
Theorem~\ref{thm:remainingmain771}, which proves the
display. At a datum with \(\ell\ge x^{50}\), \(r_{<}(D)=0\),
giving the final assertion.
\end{proof}
The estimate uses the torus normalization \(\mathcal B_x\).
The stronger normalization \(\sqrt{b_x}\) for the full
period is established on the range of
Theorem~\ref{thm:fullperiodcutoff769}. The small-prime
term \(r_{<}(D)\) and the integral comparison for the
higher-rank lifts are separate outstanding steps toward
the zero-location theorem.
\clearpage

\subsection{The torus period above the product \(k_x t_0\)}
\label{sec:rankcutoff778}
For the shorter family put \(Q_x=k_x t_0\), with
\(t_0=5r_x\) and \(k_x=r_x+1-m_x\). We retain the
original fields and twists. The interval supremum can
be estimated directly at a coefficient prime: every
node with \(Q_x<\ell\) satisfies \(x\ll\ell^{1/10}\).
This matters for a contribution proportional to
\(b_x/\ell^2\), which occurs for each scalar target.

\begin{lemma}[The large-prime local bounds below \(x^{50}\)]
\leavevmode\label{lem:rankcutoff778}
At sufficiently large nodes and at an unramified spherical
place over \(p\ge7\) with \(p>t_0+2\), the bound of
Theorem~\ref{thm:mainorbit776} holds for every main cup.
For each scalar cup one may omit the \(G/p\) term.
In particular these estimates hold at every spherical
place over a prime \(p> Q_x\).
\end{lemma}
\begin{proof}
Take the two weights \(b=b_x,\ G=xb/x_{1,x}\), and the
integer \(W=O_{q_0,d}(x^{42})\) of
Section~\ref{sec:mainorbit776}. For a scalar target
Corollary~\ref{cor:allscalartests770} applies at every
unramified \(p\ge5\). In each of its five intervals there
are \(O(W+1+b/p)\) consecutive integers of absolute value
at most \(C b\), with fixed \(C\). After deleting zero,
summing the counts of multiples of successive powers
of \(p\) bounds the total valuation by
\(O((W+1+b/p)/(p-1)+\log(2b)/\log p)\).
This proves the scalar assertion.

For a main target of constant sign, the proof of
Lemma~\ref{lem:onesignmain771} uses two numerical facts
after unramifiedness. The upper jump bound on the side
with only short positive jumps is \(B=t'-1<p-1\), so
its quotient unit lifts. The remaining side either has
nonnegative jumps or is nonpositive throughout, in
which latter case its unit lifts as well. More explicitly,
at positive sign the graph is nonnegative and the dual
nonpositive. At negative sign the graph has upper jump
\(B\) and the dual is nonnegative. These give precisely
the two alternatives of Theorem~\ref{thm:twolatticecup768}
with \(a=\lceil(b+W)/(p-3)\rceil\). All three \(E^a\)-lifts
follow from their upper jumps, and the two valuation sums
obey the interval count just given.

If the base weight \(\beta\) is constant, use
\(a_\beta=\lceil(\beta+W)/(p-3)\rceil\) on all three
extensions, whose upper bounds are now \(\beta+W\).
Uniformly for \(p\ge7\), \(a_\beta\le(\beta+W)/4+1\);
\(\beta\ge G\ge b^{1-10^{-6}}\) gives
\(\beta-W>a_\beta+B+2\) for large \(x\).
There is no negative jump in \([-a_\beta,0)\). The two
sign alternatives again give Theorem~\ref{thm:twolatticecup768}
and the same valuation bound.

For a mixed orbit, keep the two choices of labels and
the two-step comparison in Theorem~\ref{thm:mainorbit776}.
We check its numerical hypotheses. The bound
\(a\le b/p+8H_p\) only used \(p\ge7\), and
\(a\le(b+W)/4+1<b-W\) since \(W=o(G)\).
The light gap in the first choice follows as before
from \(s\le G-W-2\). In the second choice the degree
of the monic polynomial satisfies the strict inequality
there by \(b/p-G>40H_p\), independently of the size
of \(W/p\). The cost at the second prime uses \(B<p\).
These are all size requirements in constructing the
points and comparing the cochains; the sign alternation
holds for the same reason at every unramified mixed
place. The same exponent follows.

Finally \(q_{1,x}=3q_0 l_0 2^s\), with \(l_0\) a power of
three, and \(E_x\subseteq\QQ(\mu_{q_{1,x}d})\) by
Corollary~\ref{cor:oddwindow677}. Thus every ramified
prime divides \(6q_0d\). Since \(r_x\asymp x^5\) and
\(\sqrt{x}\le m_x\le2\sqrt{x}\), one has
\(k_x\asymp x^5,\ Q_x\asymp x^{10}\). Primes above
\(Q_x\) are consequently unramified and exceed \(t_0+2\)
at all sufficiently large nodes.
\end{proof}

\begin{theorem}[The full integral torus period above \(k_x t_0\)]
\leavevmode\label{thm:rankcutoff778}
In Theorem~\ref{thm:smallcoords777} let \(r_Q(D)=r_{<}(D)\)
when \(\ell\le Q_x\), and \(r_Q(D)=0\) otherwise. With its
normalization \(\mathcal B_x=f_x(m_x/k_x)b_x(1+\log(2b_x))\),
\[
 \sum_\ell(\log\ell)\sup_{D\ {\rm at}\ \ell}
      \frac{(p(D)-r_Q(D))_+}{\mathcal B_x}
       \ll_{q_0,d} V_0^{-1/2}.
\]
Thus for data with \(\ell>Q_x\) the bound holds for the
full period, including all scalar targets and all main
targets at the original twist and precision.
\end{theorem}
\begin{proof}
Only primes \(Q_x<\ell<x^{50}\) entail a further
projection loss. The two old level primes already have
zero retained image. In the spherical case there are
at most \(n_x=[E_x:\QQ]\) coordinates per target and
at most \(t_0+1\) targets. The local lemma and the
length argument of Proposition~\ref{prop:sumallmixed777}
therefore bound the additional normalized loss by
\[
 C_{q_0,d}\frac{x^{19/2}}{b_x(1+\log(2b_x))}
 \left(W+1+\frac{b_x}{\ell^2}+\frac{G}{\ell}
                       +\frac{\log(2b_x)}{\log\ell}\right).
\]
Here \(n_x/f_x=2\) and
\((t_0+1)k_x/m_x\ll x^{19/2}\); this uses the lower
bound \(m_x\ge\sqrt{x}\).

Consider first \(b_x/\ell^2\). There are constants
\(c,C>0\) such that \(c x^{10}\le Q_x\le Cx^{10}\)
at all these nodes. Thus \(\ell\ge cV_0^{10}\) and
\(x^{19/2}\ll\ell^{19/20}\). Even after the supremum
over all nodes at this prime the resulting summand,
with its factor \(\log\ell\), is at most
\(C'(\log\ell)\ell^{-21/20}\).
The inequality \(\sum_{\ell\le y}\log\ell\ll y\) bounds
its sum on the dyadic interval \([2^j L,2^{j+1}L)\)
by \(O((2^j L)^{-1/20})\). With \(L=cV_0^{10}\) the
geometric series costs \(O(V_0^{-1/2})\).

For the other terms take a separate upper bound at
each node and sum the nonnegative bounds over the
nodes. Up to fixed constants the node costs for \(W+1\)
and for \(\log(2b_x)/\log\ell\) are at most
\[
             \frac{x^{19/2+50}(W+2)}{b_x}.
\]
Here we used the same prime logarithm bound and at most
\(x^{50}\) primes. For \(G/\ell\) the cost is at most
\[
      C x^{19/2}\frac{x}{x_{1,x}}\big(1+\log(2x^{50})\big),
\]
using \(\sum_{\ell\le y}(\log\ell)/\ell\ll1+\log(2y)\).
As \(b_x\ge x^{400}\), \(W\ll x^{42}\) and
\(x_{1,x}\ge x^{10^6}\), both bounds are \(O(x^{-100})\).
Their sum is \(O(V_0^{-99})\). Add the estimate of
Theorem~\ref{thm:smallcoords777} using the positive-part
inequality in its proof. This gives the assertion.
\end{proof}
\clearpage

\subsection{The torus period away from the fixed conductor}
\label{sec:fixedconductor781}
For the shorter family put \(Q_x=k_x t_0,\ 2^s=m_x^{22}\)
and \(q_1=3q_0l_0 2^s\), where \(l_0\) is the power of
three in Corollary~\ref{cor:oddwindow677}. In particular
\(27s\le l_0\le81s\). The characteristic-zero fields are
\(E_x=\QQ(\mu_{q_1},\sqrt{-d})\), with CM type induced
from \(K_d=\QQ(\sqrt{-d})\). We first describe the actual
spacing of the places assigned the larger weight and its
behavior in a local Frobenius orbit.

\begin{lemma}[The spacing of the larger weights on the period field]
\leavevmode\label{lem:spacing781}
Identify the infinite places with \((\ZZ/q_1\ZZ)^\times\)
by choosing the embedding over the specified embedding
of \(K_d\). At sufficiently large nodes there is a set
\(\mathcal T\) of places of class \(-1\bmod q_0\), of size
at least \(2^s\exp(-C\sqrt{\log(2x)\log\log(3x)})\),
with the following property. For every prime \(p\le Q_x\)
and \(0<i<j<l_0\), there are no three members \(a,b',c\)
(allowing repetitions) with
\[
   b'/a\equiv\pm p^i,\qquad c/a\equiv\pm p^j
                   \pmod {q_1/p^{v_p(q_1)}}.
\]
The signs are independent. In particular the set has at
least \(3l_0 2^s/\log b_x\) members on the original
scales with \(\log b_x\ge x^\epsilon\) for fixed
\(\epsilon>0\). The spacing is preserved by translation
by a fixed unit. On an unramified completion over \(p\le Q_x\)
with \(p\nmid q_1 d\), labels carrying the weight \(b_x\)
from any such set cannot occur three times with positive
successive gaps whose sum is less than \(l_0\).
\end{lemma}
\begin{proof}
Since \(K_d\) has discriminant \(-d\) and \(d\nmid q_1\),
it intersects \(\QQ(\mu_{q_1})\) in \(\QQ\): otherwise a
field unramified at \(d\) would contain \(K_d\).
The oriented embeddings therefore give the indicated
group of unit classes. Take \(v=v_2(q_0)\ge2\) and
\(g=1+3l_0q_0\), a unit modulo \(q_1\).
Its order modulo \(2^{s+v}\) is \(2^s\), since
\(v_2(g^n-1)=v+v_2(n)\) for \(n\ge1\). For odd \(n\)
this follows by the geometric sum, and each subsequent
squaring adds one since \(g\equiv1\bmod4\).

Set \(M=2^s\) and \(L=2l_0\). Apply
Lemma~\ref{lem:sphere48} with modulus \(M\) and bound
\(L=O(\log M)\). Its exponents \(e\) lie in
\([0,M/(8L))\). For three such exponents and positive
\(\alpha,\alpha'\) with \(\alpha+\alpha'\le2l_0\), a
congruence
\[
     (\alpha+\alpha')e_2=\alpha e_1+\alpha'e_3\pmod M
\]
lifts to an equality of integers by this size bound.
The lemma then forces the three exponents to coincide.
Take \(\mathcal T\) to consist of the places \(-g^e\).
The order calculation makes them distinct and all have
class \(-1\bmod q_0\). The cardinality assertion follows
from the lemma and \(x^{11}\le M\le 2^{22}x^{11}\).
The exponential loss times \(3l_0\) is less than
\(x^\epsilon\) for all sufficiently large \(x\).

Suppose the two stated ratios exist for an odd \(p\).
The modulus includes \(2^{s+v}\). Raise the first ratio
to \(2j\) and the second to \(2i\). The signs disappear
and both values are \(p^{2ij}\). The order of \(g\)
therefore gives the displayed exponent congruence with
\(\alpha=2(j-i)\), \(\alpha'=2i\). All three exponents
coincide. In particular \(p^{2i}\equiv1\bmod 2^{s+v}\).
For odd \(p\), factoring \(p^{2i}-1\) at its odd and
two-power exponents gives
\[
 v_2(p^{2i}-1)=v_2(p-1)+v_2(p+1)+v_2(i).
\]
Indeed \(p^2\equiv1\bmod8\), and each squaring raises
the valuation by one. One of \(p-1,p+1\) has valuation
one. Hence \(2^{s+v}\le 2(p+1)i\le2(Q_x+1)l_0\).
This contradicts \(2^s\ge x^{11}\), since \(Q_x\ll x^{10}\)
and \(l_0\ll\log(2x)\). At \(p=2\) reduce instead
modulo \(3l_0\): every ratio is \(1\), so squaring the
first would give \(4^i\equiv1\bmod 3l_0\).
But \(v_3(4^i-1)=1+v_3(i)\) by the binomial formula,
forcing \(l_0\mid i\), again impossible. Multiplying all
three places by one unit preserves both ratios.

At an unramified prime outside \(q_1d\), arithmetic
Frobenius on a primitive \(q_1\)-th root is the
\(p\)-th power map. Conjugating an embedding to return it
to the chosen CM type negates the exponent of that root.
Hence after \(n\) steps the ratio of the associated
place classes is \(\pm p^n\) modulo \(q_1\), with the
reverse order of labels if inverse Frobenius indexes
the cycle. Reversing a triple preserves its total gap.
The two forbidden congruences prove the final assertion.
\end{proof}

\begin{theorem}[A period bound with only the fixed conductor primes retained]
\leavevmode\label{thm:fixedconductor781}
Use the sparse choice of the larger-weight places of
the shorter family, with the explicit spacing in
Lemma~\ref{lem:spacing781}. Make this choice in the
period data before choosing the simultaneous twists of
Theorem~\ref{thm:projectedcups742}. Let
\(\mathcal R=\{p:p\mid 6q_0d\}\), and put
\(r_{\mathcal R}(D)=r_Q(D)\) for \(\ell\in\mathcal R\),
with \(r_{\mathcal R}(D)=0\) otherwise. Then on the
same interval and fields as in Theorem~\ref{thm:rankcutoff778},
\[
 \sum_\ell(\log\ell)\sup_{D\ {\rm at}\ \ell}
   \frac{(p(D)-r_{\mathcal R}(D))_+}
         {f_x(m_x/k_x)b_x(1+\log(2b_x))}
                \ll_{q_0,d} V_0^{-1/2}.
\]
In particular all coefficient primes outside the fixed
set \(\mathcal R\) satisfy the full torus-period comparison
in this prime sum, including primes below \(k_x t_0\).
\end{theorem}
\begin{proof}
Fix \(p\notin\mathcal R,\ p\le Q_x\). The field is
unramified there because \(E_x\subseteq\QQ(\mu_{q_1d})\).
We have \(p\ge5\). At any spherical place choose
\[
               J=\left\lfloor\frac{120\log x}{\log p}
                                                    \right\rfloor .
\]
Here \(J\ge1\) for large \(x\), since \(Q_x\ll x^{10}\).
As \(s=22\log_2 m_x\ge11\log_2 x\), we have
\(2J\le240\log_2 x<l_0\). The preceding lemma supplies
the precise triple condition in Theorem~\ref{thm:sparsecups780}.
Writing \(G=xb_x/x_{1,x}\) and using its \(W,\eta\),
we have
\[
  p^J\le x^{120}<p^{J+1},\qquad
  \frac{\eta}{b_x}\ll_{q_0,d}
                  x^{1-10^6}+\frac{x^{42}+1}{b_x}.
\]
Since \(b_x\ge x^{400}\) and \(J\ll\log(2x)\), these imply
\(16(J+1)p^J\eta\le b_x\). They also bound
\((J+1)^2\eta\) by \(O(b_x x^{-120})\).
The logarithmic term in the local bound has the same
majorant, since \(\log(2b_x)\ll \sqrt{b_x}\) and
\(\sqrt{b_x}\ge x^{200}\). It follows that every
original main or scalar cup at this place is killed
by a power of \(p\) of exponent \(O_{q_0,d}(b_x x^{-120})\).
These estimates concern all tests in the chosen datum.
In particular membership in the smaller prime range did
not require another field, rank or twist.

At a node the projection removed when passing from
\(r_Q\) to \(r_{\mathcal R}\) has at most \(n_x(t_0+1)\)
coordinates per prime and zero length at the old level
primes. By \(\sum_{p\le Q_x}\log p\ll Q_x\ll x^{10}\),
the weighted sum of its lengths is at most
\(C_{q_0,d}\,n_x(t_0+1)b_x x^{-110}\).
After division by the displayed denominator, use
\(n_x=2f_x\) and \((t_0+1)k_x/m_x\ll x^{19/2}\).
The result is \(O_{q_0,d}(x^{-100})\). At \(p>Q_x\)
the preceding projection was already zero. Summing the
nonnegative node bounds over all integers \(x\ge V_0\)
costs \(O(V_0^{-99})\), also after the primewise
supremum. Add this to Theorem~\ref{thm:rankcutoff778}
by the positive-part inequality, obtaining the claim.
\end{proof}
The remaining torus contribution is now confined to
primes dividing a fixed integer. Their ramification
indices can grow with the node. Controlling them, and
the integral projections in higher rank, is still needed
to derive the asserted zero location for Dirichlet
\(L\)-functions.
\clearpage

\subsection{All coefficient primes for the marked torus period}
\label{sec:alltorus783}
We finish the sum for the shorter marked torus family
at normalization \(\mathcal B_x=f_x(m_x/k_x)b_x(1+\log(2b_x))\).
The primes left by Theorem~\ref{thm:fixedconductor781}
divide \(6q_0d\). Two properties of their embeddings are
needed simultaneously. The larger weight is sparse even
after grouping the embeddings by their restriction to the
maximal unramified subfield, and all its occurrences in
one completion have a common sign.

\begin{lemma}[Spacing and sign at the conductor primes]
\leavevmode\label{lem:conductorsign783}
Fix a place \(w\mid p\), \(p\mid6q_0d\), in the family
of Theorem~\ref{thm:fixedconductor781}, and write
\(K=E_{x,w}\), \(K_0\) for its maximal unramified subfield.
Mark an embedding of \(K_0\) if some extension to \(K\)
has base weight \(b_x\). On the periodic residue-field
orbit there are no three marked occurrences with positive
successive gaps whose sum is \(<l_0\). At every full
embedding over \(w\) carrying the larger weight, the
main jump has the same sign.
\end{lemma}
\begin{proof}
All these primes satisfy \(p\le Q_x\) for large \(x\).
Put \(q'=q_1/p^{v_p(q_1)}\). The \(q'\)-th roots of unity
in \(K\) belong to \(K_0\): their reductions have the
same order, and each lifts uniquely from the residue field
as a root of \(X^{q'}-1\). Hence the image of a primitive
\(q'\)-th root is determined by a residue-field embedding
alone, and Frobenius changes it by raising to \(p\).
For three purported marked occurrences, choose a full
embedding over each one that carries the larger weight.
Changing these embeddings to the chosen CM orientation
changes the exponents only by signs. As in
Lemma~\ref{lem:spacing781} their ratios would give the
two forbidden congruences modulo \(q'\), reversing their
order for inverse Frobenius. This holds also for repeated
visits to the same residue-field label.

For \(p\mid q_1\) one has \(p\mid q_0\), since \(6\mid q_0\)
and \(l_0\) is a power of three. The quadratic field
\(K_d\) splits at every such prime. Indeed for odd \(p\)
one has \(-d\equiv1\bmod p\) from
\(d\equiv q_0/2-1\bmod q_0\), and the square root lifts
by Hensel's lemma. At two, \(16\mid q_0\) gives
\(-d\equiv1\bmod8\), so \(-d\) is a square in \(\QQ_2\).
All embeddings of this fixed completion then induce
the same embedding on the quadratic field.

The other possibility is \(p=d\). Now all \(q_1\)-th
roots of unity are unramified at \(p\), and the ratio
of the exponents of any two local embeddings on one
primitive root lies in \(\langle d\rangle\bmod q_1\).
This follows from the cyclic Galois action on the residue
field and uniqueness of the prime-to-\(p\) lifts. If
two larger-weight embeddings had opposite CM orientations,
their unoriented exponents would have ratio \(-1\bmod q_0\):
their oriented exponents both belong to the same class
\(-1\bmod q_0\). But \(d^2\equiv1\bmod q_0\) and neither
\(1\) nor \(d\) equals \(-1\bmod q_0\), by the specified
congruence and \(4\mid q_0\). This rules out opposite
orientations. The same argument applies if a uniform
translation changes the common class of the larger
weights. The main jump sign is determined by the
orientation, up to a common reversal for a conjugate
character, as in Theorem~\ref{thm:negativeinterval740}.
\end{proof}

\begin{proposition}[Integral cups at every conductor prime]
\leavevmode\label{prop:conductorcups783}
With these data, for \(p\mid6q_0d\) at each spherical
place \(w\mid p\), every original main or scalar cup
is killed by \(p^h\) for an integer
\[
                      h\ll_{q_0,d} b_x x^{-108}.
\]
The bound is uniform in \(w\), the trace rank and all
integral tests and precisions at \(x\).
\end{proposition}
\begin{proof}
Take the coefficient extension containing \(K\) and use
its unramified label cycle as in
Section~\ref{sec:ramifiedorders782}. Use \(b=b_x,\ G=xb/x_{1,x}\),
\(W\ll x^{42}\), \(\eta=2(G+W+1)\) as before.
The graph and test extensions are crystalline at the
spherical place. The third extension is defined by
lifting the original equivariant map modulo \(\pi^a\)
to an \(\mathcal O\)-linear map \(\psi\) and using the
cocycle \((g\psi-\psi)/\pi^a\); its contraction recovers
that same finite pairing as in
Corollary~\ref{cor:periodlabels764}. At a full embedding,
the jumps of its kernel are differences at that embedding.
The tensor and dual calculation is unchanged by ramification.
On an unmarked residue-field label all three upper jumps
are therefore at most \(\eta\), and everywhere they are
at most \(b+W\), after fixing \(W\) sufficiently large.

Choose \(J=\lfloor120\log x/\log p\rfloor\),
\(L=2J+1\). The preceding lemma gives the no-triple
condition since \(2J<l_0\). Choose the cuts followed
by \(J\) unmarked labels. The elementary gap argument
of Theorem~\ref{thm:sparsecups780} applies verbatim to
this periodic orbit: successive cuts are at distance
\(1\le l\le L\). Write \(e=[K:K_0]\le n_x\ll_{q_0,d}x^{12}\).
Since \(p^{J+1}>x^{120}\), it is at least \(8e\).
Theorem~\ref{thm:ramifiedcut782} supplies a common
\(p^\kappa\)-lift at each cut on all three extensions with
\[
          \kappa\ll 1+J(d_K+1)\eta+e(b+W)/p^{J+1}.
\]
Scaling to a common exponent is legitimate around the
entire orbit. Use the third lift for the pairing
\(\widetilde Q=p^\kappa Q\). Its bound on level zero
kernels is \(\tau^{-1}A_i(1)\); at level \(l\) from the
preceding cut it is \(\tau_l^{-1}A_i(1)\), where
\(\tau_l/\tau\) is a unit times \(\prod_{j<l}P_j\).
This follows by applying Frobenius to the integral
contraction and \(\mu/\tau\in E A^\times\). The same
lattices and contractions are used at the original
precision in the finite Witt ring.

At each cut put \(D=\eta+2\) and apply
Proposition~\ref{prop:ramifiedorders782} to the two
level zero points, imposing depth \(D\) at \(0\le j<L\).
The increase in the quotient exponent is at most
\[
 C_1=C_{q_0,d}\big((L+1)(L+d_K+4)\eta
                       +(L+1)\log(2b)/\log p\big).
\]
Indeed on all full embeddings over a residue label
the low jumps lie in \([-\eta,\eta]\); the far jumps
of the given side, if present, lie within \(W\) of
the same signed \(b\), by the sign lemma. Thus the
union of the shifted integer intervals is contained
in at most two intervals totaling \(O(\eta)\) indices
at each \(j\). Each has absolute values at most \(8b\).
Counting multiples of powers as in
Theorem~\ref{thm:sparsecups780} proves the bound.
On the subgroup \(H_L=G_{K(\varpi^{1/p^L})}\) the
displacements now lie in
\(\delta\prod_{j<L} P_j^D\) times their kernel lattices.

There is a choice of side \(U\) with lower jumps
\(\ge-\eta\) at every full embedding and opposite side
\(V\) with upper jumps \(\le\eta\) everywhere. At
unmarked embeddings both sides satisfy both bounds.
At marked embeddings use the common sign of the larger
jump, assigning opposite signs to the graph and dual;
an absent far jump imposes no condition. Consequently
for consecutive cuts \(c,i=F^l c\), with
\(S=\prod_{j<l}P_j^\eta\), the ramified height
inclusions give
\[
      S M_{l,i}(U)\subseteq M_{0,i}(U),\qquad
      S M_{0,i}(V)\subseteq M_{l,i}(V).
\]
In \(M_{l,i}(U)\) take the \(l\)-fold Frobenius of a
pure \(p^\kappa\)-lift at \(c\), and impose the same
orders \(D\) at \(0\le j<l\). These lattices are regular
at those \(P_j\), so the cost for the resulting \(y_i\)
is at most \(\kappa+C_1\). Take a common exponent
\(e_*\le\kappa+C_1\), enlarged by one if needed for
rounding, for all cut and auxiliary points, by scaling
them individually. As before, at noncuts use Frobenius
of the preceding pair.

At a cut first change its \(U\)-point to \(y_i\).
The kernel difference belongs to \(S^{-1}M_{0,i}\).
The other displacement and level zero pairing leave
a factor \(u\). Next change its \(V\)-point to the
\(l\)-fold Frobenius of the point at \(c\). This
difference is in \(S^{-1}M_{l,i}\); the displacement
of \(y_i\) supplies \(\delta\prod_{j<l}P_j^D\),
clearing both \(S\) and the level \(l\) Tate denominator
and leaving a factor \(u\). Finally change \(y_i\) to
the corresponding \(l\)-fold Frobenius on \(U\); the
difference is in \(M_{l,i}\) and the fixed displacement
has factor \(u^{p^l}\tau_l\).
This argument also handles translated differences since
the group preserves the lattices and each principal
ideal involved. The initial cup and the three cochain
corrections are therefore small, in the sense of
Theorem~\ref{thm:cuppath757}. Frobenius permutes the
residue-field idempotents in one cycle and their sum
is one. The terminal pair at each label is exactly the
single Frobenius of the preceding initial pair. Summing
and taking constants in finite Witt vectors gives
annihilation by \(p^{2e_*+\kappa}\) on \(H_L\), by the
same cochain descent as in Theorem~\ref{thm:sparsecups780}.
This descent uses the Witt constants over any finite
extension of \(\QQ_p\), including \(p=2\), by
Lemma~\ref{lem:wittconstants757}. The group index is
\(p^L\), since \(X^{p^L}-\varpi\) is Eisenstein over
\(K\). Corestriction and the integral coefficient
retraction thus give \(h=2e_*+\kappa+L\) over the
original field and ring.

Here \(L+d_K\ll_{q_0,d}1+\log(2x)\). The term
\(e(b+W)/p^{J+1}\) is \(O_{q_0,d}(b x^{-108})\).
All terms involving \(\eta\) or \(\log(2b)\) satisfy
this bound as well, using \(G/b\le x^{1-10^6}\),
\(W\ll x^{42}\) and \(b\ge x^{400}\). This proves the
claimed uniform estimate.
\end{proof}

\begin{theorem}[The full projected torus period at the source normalization]
\leavevmode\label{thm:alltorus783}
Fix \(q_0,d\) in the shorter marked torus family, and use
the sparse choices of Lemma~\ref{lem:spacing781}. For
the simultaneous twists of Theorem~\ref{thm:projectedcups742}
on the same fields \(E_x\), the same interval beginning
at \(V_0\), and the original integral tests, one has
\[
  \sum_{\ell}(\log\ell)\sup_{D\ {\rm at}\ \ell}
       \frac{p(D)}{f_x(m_x/k_x)b_x(1+\log(2b_x))}
               \ll_{q_0,d} V_0^{-1/2}.
\]
The sum includes every coefficient prime and the supremum
includes all the original trace ranks, separating families
and precisions in that comparison.
\end{theorem}
\begin{proof}
It remains in Theorem~\ref{thm:fixedconductor781} to
bound \(r_{\mathcal R}(D)\) for \(\mathcal R=\{p:p\mid6q_0d\}\).
At a node the original level primes have zero retained
length. At the spherical places above a given
\(p\in\mathcal R\), the retained joint image injects
into a product of at most \(n_x(t_0+1)\) cyclic
coordinates. The proposition bounds the normalized
length of each by \(O_{q_0,d}(b_x x^{-108})\).
The sum of \(\log p\) in this fixed set depends only
on \(q_0,d\). Divide by the displayed denominator,
and use \(n_x=2f_x\) and
\((t_0+1)k_x/m_x\ll x^{19/2}\).
The weighted sum at a node is \(O_{q_0,d}(x^{-98})\).
Its sum over the integer nodes \(x\ge V_0\) is
\(O_{q_0,d}(V_0^{-97})\), which bounds the corresponding
sum of interval suprema. Now use
\(p(D)\le(p(D)-r_{\mathcal R}(D))_++r_{\mathcal R}(D)\)
and Theorem~\ref{thm:fixedconductor781}.
\end{proof}
This proves the estimate for the original marked torus
projections. Section~\ref{sec:missing} also requires
integral control of the higher-rank theta chain in the
comparison between local divisibility and the archimedean
upper bound for the prime-power estimate. That implication
is still outstanding.
\clearpage

\section{Congruence modules and integral projections}\label{nav:torus-modules}

\subsection{The torus budget and the determinant of a higher-rank projection}
\label{sec:projectionbudget784}
In the denominator comparison for a chain with \(m_x\)
lifts there are two tests at a coefficient prime. One
concerns the projected inner product for a designated
torus line. The other asks for an operator on the image
of a higher-rank lift which clears its projected integral
forms and has controlled determinant. We convert the
bound on the former to the units in which both tests
are measured, and identify the exact loss for the latter.

\begin{proposition}[One family of logarithmic bounds for every main pair]
\leavevmode\label{prop:torusbudget784}
At each node write \(a_x=xb_x,\ H_x=\log b_x\), and put
\[
       B_x^{\rm det}=\frac{f_x a_x H_x^3}{m_x k_x}.
\]
For the fields, data and simultaneous twists of
Theorem~\ref{thm:alltorus783} there are nonnegative
numbers \(\gamma^{(1)}_{x,\ell}\) such that, for every
main pair, every permitted twist power and every pair
of its integral test forms,
\[
     |\iota\langle P\Xi,P\Xi'\rangle|_\ell
                     \le \exp(\gamma^{(1)}_{x,\ell}).
\]
One set of numbers works for all lift indices,
coefficient embeddings, and separating families. For
every finite set of primes \(S\),
\[
  \sum_{\ell\in S}\gamma^{(1)}_{x,\ell}
   \ll_{q_0,d} \frac{V_0^{-1/2}}{H_x^2}\,B_x^{\rm det}.
\]
\end{proposition}
\begin{proof}
Let \(\omega_\ell=(\log\ell)\sup_{D\ {\rm at}\ \ell}
p(D)/\mathcal B_x\), where the denominator inside the
supremum is the one for the node of \(D\). Use zero for
an empty set. Theorem~\ref{thm:alltorus783} gives
\(\sum_\ell\omega_\ell\ll_{q_0,d}V_0^{-1/2}\). Define
\(\gamma^{(1)}_{x,\ell}=\mathcal B_x\omega_\ell\).
The definition of \(p(D)\) in
Section~\ref{sec:projectedcups742} gives the asserted
inequality also for a zero projection. For the budget,
\[
  \frac{\mathcal B_x}{B_x^{\rm det}}
      =\frac{m_x^2}{x}\frac{1+\log(2b_x)}{H_x^3}
      \ll H_x^{-2},
\]
because \(m_x\le2\sqrt{x}\) and \(H_x\ge1\).
Multiplication by the bound for \(\sum_\ell\omega_\ell\)
proves the claim.
\end{proof}

\begin{theorem}[The exact determinant loss on integral tested forms]
\leavevmode\label{thm:projectionindex784}
At one coefficient prime \(\ell\) and embedding, transport
the definite level space, its true lift range \(Z\), and
its orthogonal projection \(P\) to the local coefficients.
Let \(T\) be the space orthogonal to all false lifts
in \(Z^\perp\). In a finite local field of definition
\(C/\QQ_\ell\) let \(L\) be the lattice of forms integral
at this embedding and let \(\mathcal O\) be the integers
of \(C\). Write
\[
 L_T=T\cap L,\qquad L_Z=Z\cap L,\qquad
 L_0=T\cap\ker P\cap L,\qquad
 \lambda=\frac{\operatorname{length}_{\mathcal O}
                   (L_T/(L_Z+L_0))}{v_\pi(\ell)} .
\]
The minimum of \(-\log|\det\Lambda|_\ell\), over invertible
operators on \(Z\) such that \(\Lambda P(L_T)\subseteq L_Z\),
is \(\lambda\log\ell\). It is attained, and extending the
local coefficient field preserves the minimum. For \(Z=0\)
both sides are zero. If \(D_Z=\dim Z>0\) and \(r_x\)
is the target rank of the marked data, the least
nonnegative \(\gamma\) for which such an operator has
\(|\det\Lambda|_\ell\ge \exp(-D_Zr_x\gamma)\) is
\(\lambda\log\ell/(D_Zr_x)\).
\end{theorem}
\begin{proof}
We describe the finite local model used here. Definite
level forms are given by their values at finitely many
double-coset representatives. The values lie in finite
Weyl coefficient spaces and satisfy linear stabilizer
relations. At a fixed embedding the local gauge is the
coordinate maximum after the inverse weight action of
the representative at the places over \(\ell\). This finite
list computes integrality: rational translation cancels
the weight action by equivariance of the form, and right
translation by the level preserves the integral weight
lattice. For the latter assertion the local matrices
and their inverses belong to the integral general linear
group for the standard level and its subgroups; the
integral Weyl representation and integral determinant
twists take these to integral invertible matrices.
A finite level character has unit values as well.
Equivalently one may use representatives at the level
kernel and impose the finite level relations.
Thus the lattice is the inverse image of the integral
coordinate cubes under finitely many linear evaluations.
Together the evaluations are injective. Over a finite
local field their integral inverse image is a full
lattice: a common scalar clears their entries in a fixed
basis, and an invertible coordinate minor bounds the
inverse image inside a second multiple of that lattice.

A choice of basis for \(Z\), the matrix of \(P\), and a
basis for the sum of all false-lift spaces require only
finitely many coefficients, also when the possible
Hecke tests range over infinitely many places. Write
the orthogonality conditions for that basis before
transporting their coefficients. These coefficients,
the stabilizer equations and gauge matrices all belong
to a common finite extension of \(\QQ_\ell\): each
transported coefficient belongs to \(\overline{\QQ}_\ell\),
and the local components of the finitely many adelic
representatives lie in the completions of the given
number field. This also explains how to take the lattice
at the specified embedding while retaining the actual
orthogonality tests.

Since the false lifts are in \(Z^\perp\), we have
\(Z\subseteq T\). Thus \(P(L_T)\) contains \(L_Z\) and
is a lattice in \(Z\). Projection induces an isomorphism
\[
         L_T/(L_Z+L_0)\ \simeq\ P(L_T)/L_Z .
\]
Indeed if \(Pv\in L_Z\) then \(v-Pv\in L_0\), and the
reverse inclusion and surjectivity follow directly.
Choose a basis \(e_1,\ldots,e_{D_Z}\) of \(L_Z\) in
which the larger lattice has basis \(\pi^{-a_i}e_i\)
with integers \(a_i\ge0\), by the elementary divisor
theorem. Its quotient length is \(\sum_i a_i\).
Every eligible operator sends these last basis vectors
to integral columns. Taking their determinant forces
\(v_\pi(\det\Lambda)\ge\sum_i a_i\). The diagonal
operator \(e_i\mapsto\pi^{a_i}e_i\) attains equality.
Since \(v_\ell(\ell)=1\), conversion to logarithms gives
the formula. Under a finite coefficient extension the
same bases have elementary exponents multiplied by
the relative ramification index, just as \(v_\pi(\ell)\)
is. Intersections and images used above commute with
this flat scalar extension: after clearing a common
denominator they are kernels and images of integral
matrices. The same determinant lower bound therefore
applies to any operator over \(\overline{\QQ}_\ell\)
by adjoining its finitely many entries. The empty
determinant equals one in dimension zero.
\end{proof}

\begin{corollary}[The remaining higher-rank inequality in the denominator budget]
\leavevmode\label{cor:remainingbudget784}
Fix a node and its simultaneous twists as above. For
a given lift index \(j\) let \(h_{x,j,\ell}\) be the
supremum of the last quotient in
Theorem~\ref{thm:projectionindex784} over the designated
coefficient embeddings and their corresponding twists
at \(\ell\); use zero for a zero-dimensional range.
If for every finite set of primes \(S\)
\[
                   \sum_{\ell\in S}h_{x,j,\ell}
                              \le \tfrac12 B_x^{\rm det},
\]
then the integral operator test and all main-pair tests
at this index admit a common logarithmic bound of total
sum at most \(B_x^{\rm det}\), for sufficiently large
\(V_0\). Conversely any common bound of this total
size implies \(\sum_{\ell\in S}h_{x,j,\ell}\le B_x^{\rm det}\).
\end{corollary}
\begin{proof}
At a single prime take the maximum of \(h_{x,j,\ell}\)
and \(\gamma^{(1)}_{x,\ell}\). The theorem supplies the
operators, and the proposition supplies the main-pair
inequalities for the same twists. The sum of the maxima
is at most the sum of the two sums. The proposition's
bound is at most \(B_x^{\rm det}/2\) for sufficiently
large \(V_0\). For the converse, the exact local
determinant minimum bounds \(h_{x,j,\ell}\) from above
by any permissible common logarithmic bound at that
prime; sum over \(S\).
\end{proof}
The sum of the higher-rank indices in the hypothesis
requires an arithmetic estimate on the actual tested
lattice. The torus-period estimate proves its own term
in the common budget on this same family.
\clearpage

\subsection{Split Hecke values on the complement of the true lifts}
\label{sec:splitcost785}
Fix a node \(x\), one of its prescribed twists, a lift index
\(j\), and an embedding at \(\ell\). Use \(Z,P,T,L\) and
\(\lambda\) from Theorem~\ref{thm:projectionindex784}.
The mark at the two split level primes \(p_0,p_1\) is a
central idempotent \(e\) for the finite group
\[
  \Gamma_x=\prod_{i=0}^1\prod_{v\mid p_i,\ v\ {\rm in}\ E_x^+}
                     {\rm GL}_{r_x}(\mathbf F_{p_i}).
\]
Write \(M_x=|\Gamma_x|\). At an unramified split place \(w\)
of \(E_x\) away from the level, the twist, and \(\ell\), let
\(q=Nw\), let \(s<n\) be the two ranks in the lift, and put
\(t=n-s\). Write \(\mu_W\) for the unitary splitting character
attached to the source and put
\[
   \alpha_w=q^{s/2}\mu_{W,w}(\varpi_w),\qquad
   X_w=(q^{\,t-1}\alpha_w)^{-1},\qquad
   A_w=\sum_{i=0}^{n}(-X_w)^i q^{i(i-1)/2}T_{w,i}.
\]
Here \(T_{w,i}\) is the unnormalized double coset for the
diagonal matrix with \(i\) entries \(\varpi_w\) and the
other entries one, in the convention of the arithmetic
Hecke polynomial. We use the definite inner product to
form orthogonal complements before transporting them to
the local coefficient field.

\begin{proposition}[The split Hecke operators on tested forms]
\leavevmode\label{prop:testedkernel785}
At each of these places \(Z\subset eV\cap\ker A_w\), where
\(V\) is the whole definite level space. Put
\(U=eT\cap Z^\perp\). Then \(eT=Z\oplus U\);
each \(A_w\) acts invertibly on \(U\), and both \(A_w\)
and \(M_xe\) preserve the integral lattice \(L\).
For \(y\in T\) the orthogonal projection of \(y\) onto
\(eV\cap\ker A_w\) is \(Py\).
\end{proposition}
\begin{proof}
The good source forms belong to the sum of the marked
types defining \(e\). The lift intertwines their action.
Indeed, changing a marked frame on the source changes
the matching frame in the theta kernel. In its matrix
coordinate the two changes act on opposite sides of the
identity block. Simultaneously changing the two frames
preserves the block congruence, and the determinant
character in the form action accounts for the scalar
of the splitting. Changing variables by this frame
change in the source integral gives equivariance for
each group element and hence for \(e\). Thus \(eZ=Z\).

The good split Hecke algebra is a commuting algebra of
normal operators for the definite pairing. In fact the
adjoint of each double coset is the inverse double coset,
and the spherical algebra is commutative. The actions at
the marked and good places also commute. On a spherical
theta constituent the added unitary Satake entries are
\(\mu_{W,w}(\varpi_w)q^{(t-1)/2-i}\), \(0\le i<t\), by
Proposition~\ref{prop:split175}. Passage to the
arithmetic polynomial multiplies each entry by
\(q^{(n-1)/2}\). The added factors are therefore
\(1-\alpha_w q^i X\), \(0\le i<t\), in reverse order.
The last of these factors vanishes at \(X_w\). Decomposing
the source into its spherical constituents gives \(A_wZ=0\).

The central idempotent \(e\) is self-adjoint. This follows
from the unitarity of the finite group action and the
decomposition into its isotypic spaces. For each place
used in the false-lift tests the corresponding false
space is the intersection of \(eV\), the kernel of the
evaluated Hecke polynomial there, and \(Z^\perp\).
It lies in \(eV\), so \(e\) preserves \(T\). Each of
these three spaces is stable under \(A_w\) and its
adjoint. For \(Z^\perp\) this follows from normality
and \(A_wZ=0\). For the other two it follows from
commutation in the spherical and marked algebras.
It follows that \(U\) is stable. A vector of \(U\)
killed by \(A_w\) belongs to the false space at \(w\)
and is orthogonal to that space. Positivity makes it
zero. Finite dimensionality gives invertibility on \(U\).
Finally, \(eV\cap\ker A_w\) is the orthogonal direct
sum of \(Z\) and the false space at \(w\), which proves
the projection formula.

Each \(T_{w,i}\) is a sum of translations with trivial
component at \(\ell\), so it preserves the gauge lattice.
The rank-normalized character has a continuous
\(\ell\)-adic avatar as in Corollary~\ref{cor:normalizedtransfer722}.
Its values at good Frobenius elements are units by
compactness. Thus \(\alpha_w\) and \(q\) are
\(\ell\)-units, and \(A_w L\subseteq L\).
The coefficient of a group element \(g\) in \(M_x e\)
is a sum of terms \(\chi(1)\chi(g^{-1})\) for
irreducible characters of \(\Gamma_x\). These are
algebraic integers. Every \(g\) acts on a marked frame
by an invertible integral change and a root of unity.
This also preserves the gauge when \(\ell=p_i\), since
the change and its inverse act on the integral weight
lattice. The character formula now gives \(M_x eL\subseteq L\).
\end{proof}

\begin{theorem}[A bound with separate multiplicities for the competing systems]
\leavevmode\label{thm:heckeflag785}
Let \(\mathcal W\) be any nonempty finite set of the split
places above. Split \(U\) into the joint eigenspaces
\(U_\nu\) of \((A_w)_{w\in\mathcal W}\), grouping equal
tuples of eigenvalues. Put \(d_\nu=\dim U_\nu\),
\(D=\dim Z\), and let \(a_{\nu,w}\) be the eigenvalue
on \(U_\nu\). With \(v_\ell(\ell)=1\) one has
\[
 \lambda\ \le\ D v_\ell(M_x)
       +\sum_\nu \min(D,d_\nu)
                         \min_{w\in\mathcal W}v_\ell(a_{\nu,w}).
\]
All the eigenvalues in this sum are nonzero and integral.
In particular the congruence index vanishes if
\(\ell\nmid M_x\) and each competing tuple has a unit
entry, even if different places supply the units.
\end{theorem}
\begin{proof}
Enlarge the local field to define the finitely many
eigenspaces, and denote its valuation ring by \(R\),
with uniformizer \(\pi\). Use \(v_\pi\) and ordinary
\(R\)-length until the last step. The commuting normal
operators admit a joint eigenspace decomposition over
the complex numbers. It transports to the enlarged
local field. Integrality of their matrices on \(L\)
makes all eigenvalues integral. Their nonvanishing on
\(U\) follows from Proposition~\ref{prop:testedkernel785}.

Put \(B=L\cap eT,\ B_Z=L\cap Z,\ B_U=L\cap U\).
Then \(B_Z\subseteq PB\subseteq P(L\cap T)\), and
\[
                 M_x P(L\cap T)\subseteq PB.
\]
Indeed \(Pe=P\), and \(M_x e\) takes \(L\cap T\)
into \(B\). The length of \(P(L\cap T)/PB\) is at
most \(D v_\pi(M_x)\) by the elementary divisor theorem.
The other quotient \(PB/B_Z\) is isomorphic to
\[
         C=\widetilde B_U/B_U,\qquad
                      \widetilde B_U=(1-P)B.
\]
Both quotients identify with \(B/(B_Z+B_U)\) by the
kernel calculation of Theorem~\ref{thm:projectionindex784}.
In particular \(C\) is generated by at most \(D\) elements.
For every \(w\in\mathcal W\) one has
\(A_w\widetilde B_U\subseteq B_U\), because \(A_wP=0\)
and \(A_w B\subseteq B\cap U\).

Order the different tuples in any way and filter \(U\)
by their successive direct sums \(U^{(i)}\), starting
with zero. Intersect \(\widetilde B_U\) with these
subspaces and take their images in \(C\). In the quotient
\(U^{(i)}/U^{(i-1)}=U_{\nu_i}\) write \(S_i,S_i^0\)
for the images of the intersections of \(U^{(i)}\)
with \(\widetilde B_U,B_U\), respectively. The \(i\)-th
successive quotient of the filtration of \(C\) is
\(S_i/S_i^0\). Indeed a lift from
\(\widetilde B_U\cap U^{(i)}\) has image in \(S_i^0\)
exactly when subtracting a vector of \(B_U\cap U^{(i)}\)
moves it to \(U^{(i-1)}\).

The two images are full lattices, so \(S_i/S_i^0\)
has at most \(d_{\nu_i}\) generators. It also has at
most \(D\). To see this, pull back any submodule of
\(C\) to \(R^D\) under a surjection onto \(C\).
The pullback is free of rank at most \(D\), and
passing to a quotient preserves this bound on generators.
The inclusion \(A_w\widetilde B_U\subseteq B_U\)
on the \(i\)-th quotient reads
\(a_{\nu_i,w}S_i\subseteq S_i^0\). Its length is
therefore at most \(\min(D,d_{\nu_i})\) times the least
of the \(v_\pi(a_{\nu_i,w})\). Summing the lengths
and dividing by \(v_\pi(\ell)\) proves the estimate.
The coefficient-extension assertion of
Theorem~\ref{thm:projectionindex784} puts it back in
the original normalization.
\end{proof}

\begin{corollary}[The contribution of the marking on the period interval]
\leavevmode\label{cor:markbudget785}
For \(D>0\) the first term of
Theorem~\ref{thm:heckeflag785} contributes at most
\(f_x r_x\log(p_0p_1)\) to the sum of the logarithmic
bounds per \(D r_x\), uniformly in all \(j\), embeddings
and twists at this node. Its ratio to \(B_x^{\rm det}\)
is \(O(x^{10}/(b_x(\log b_x)^2))\).
\end{corollary}
\begin{proof}
There are \(f_x\) selected places above each \(p_i\).
Counting all matrices gives
\(\log M_x\le f_x r_x^2\log(p_0p_1)\).
The sum of \(v_\ell(M_x)\log\ell\) over primes is
exactly \(\log M_x\). Divide by \(r_x\) and use
\(\log(p_0p_1)\le22\log b_x\) on the prescribed
interval. In the budget of Proposition~\ref{prop:torusbudget784}
we have \(a_x=xb_x,\ m_x\le2\sqrt{x}\), and
\(k_x,r_x\ll x^5\). Substitution gives the stated ratio.
\end{proof}
The second term of the theorem concerns only systems
which survive the false-lift tests and for which every
selected split Hecke value is divisible at \(\ell\).
Controlling the sum of their weighted valuations on the
period interval is a further arithmetic question.
\clearpage

\subsection{The full cyclotomic string in the congruence module}
\label{sec:fullstring786}
Use the field, twist, ranks and marked idempotent of
Section~\ref{sec:splitcost785}. On \(Z\) and on \(U\)
now use the full collection of spherical Hecke coefficients
at good split places away from \(\ell\). Write \(\xi\)
for a distinct system on \(Z\) and \(\nu\) for one on \(U\).
Their multiplicities in these spaces are \(D_\xi,d_\nu\),
respectively. Thus \(D=\sum_\xi D_\xi\). The individual
theta lifts of spherical eigenvectors are spherical
eigenvectors by Proposition~\ref{prop:split175}. This
makes \(Z\) stable under these coefficients and their
adjoints. The kernels in the false-lift tests and the
mark commute with them, so \(U\) is stable as well.

For either kind of system let \(u_\sigma\) be the
representation of \(G_{E_x}\) from
Theorem~\ref{thm:splitattachment722}, tensored with
the character of Corollary~\ref{cor:normalizedtransfer722}
having value \(\alpha_w^{-1}\) at geometric Frobenius.
Use stable lattices over a common finite local coefficient
field. Put
\[
 P_\sigma(g;Y)=\det(Y-u_\sigma(g)),\qquad
 Q_g(Y)=\prod_{a=0}^{t-1}(Y-\kappa(g)^{-a}),
\]
where \(\kappa\) is the cyclotomic character. For a
competing system define
\[
 \delta_\nu=\min_{\substack{g\in G_{E_x}\\0\le k<t}}
 v_\ell\!\left([Y^k]\big(P_\nu(g;Y)\bmod Q_g(Y)\big)\right),
                 \qquad v_\ell(0)=+\infty .
\]
The remainder is from division by a monic polynomial over
the valuation ring. Multiplicities in \(Q_g\) are retained
even when two or more of its roots coincide after reduction.

\begin{proposition}[Integral Hecke operators for the remainder]
\leavevmode\label{prop:remhecke786}
The remainder of \(P_\xi(g;Y)\) is zero for every true
system and every \(g\). For each \(\nu\) the depth
\(\delta_\nu\) is finite and is attained at a good split
Frobenius. Each coefficient of the remainder at such a
place is the eigenvalue of an integral Hecke operator
which kills \(Z\).
\end{proposition}
\begin{proof}
At a good split place let \(g_w\) be geometric Frobenius.
The elementary symmetric coefficients of \(u_\sigma(g_w)\),
with alternating signs, are the eigenvalues of
\[
            (-1)^i q^{i(i-1)/2}\alpha_w^{-i}T_{w,i}
                      \qquad(0\le i\le n).
\]
This follows from split local compatibility and the
normalization defining \(u_\sigma\). These operators
preserve the lattice by Proposition~\ref{prop:testedkernel785}.
The added roots on a true lift are \(1,q,\ldots,q^{t-1}\),
and \(\kappa(g_w)=q^{-1}\). Thus \(Q_{g_w}\) divides
\(P_\xi(g_w;Y)\). Monic long division expresses each
remainder coefficient as an integral linear combination
of the input polynomial coefficients with scalars in the
ring generated by the coefficients of \(Q_{g_w}\).
Applying the same linear combination to the displayed
operators gives the required integral operator.

We record why using split places loses no congruence
information. The representations and \(\kappa\), reduced
at any fixed precision, factor through a finite Galois
extension of \(E_x\). Enlarge it to a finite Galois
extension of \(E_x^+\). An element of \(G_{E_x}\)
induces an element fixing \(E_x\) in the latter finite
group. By Chebotarev choose a prime of \(E_x^+\) and a
place above it with this geometric Frobenius, avoiding
any fixed finite set. Its restriction to \(E_x\) is
trivial, so the prime splits in \(E_x/E_x^+\) and the
relative residue degree is one. The Frobenius at the
resulting place of \(E_x\) has the desired reductions.
All remainder coefficients depend polynomially on the
characteristic coefficients and on units
\(\kappa(g)^{-a}\). Their vanishing for a true system
at good split places therefore implies vanishing on
\(G_{E_x}\).

For a competing system the evaluated operator \(A_w\)
has nonzero eigenvalue at every such place by
Proposition~\ref{prop:testedkernel785}. This eigenvalue
equals \(q^{-n(t-1)} P_\nu(g_w;q^{t-1})\), and
\(Q_{g_w}(q^{t-1})=0\). At least one remainder coefficient
is consequently nonzero. The finite valuations lie in a
discrete set bounded below by zero, so a minimum exists.
Apply the preceding Frobenius argument one uniformizer
power beyond a coefficient attaining it. This realizes
the same minimum at a good split place.
\end{proof}

Identify systems whose residual representations have
isomorphic semisimplifications over the algebraic closure
of the residue field. For such a class \(b\) put
\(D_b=\sum_{\xi\in b}D_\xi\). For a true and a competing
system put
\[
 e_{\xi\nu}=\min_{\substack{g\in G_{E_x}\\0\le i<n}}
     v_\ell\!\left([Y^i](P_\xi(g;Y)-P_\nu(g;Y))\right).
\]
This number is finite. Indeed equality of all coefficients
at every good place would give the same eigenvalue of
\(A_w\), which vanishes on \(Z\) and is nonzero on \(U\).

\begin{theorem}[Two bounds inside each residual representation]
\leavevmode\label{thm:residualflag786}
For a competing system \(\nu\) write \(b=b(\nu)\) for its
residual class. Define
\[
\begin{aligned}
   H_\nu&=\min(D_b,d_\nu)\,\delta_\nu,\\
   J_\nu&=\sum_{\xi\in b}\min(D_\xi,d_\nu)\,e_{\xi\nu}.
\end{aligned}
\]
In the sum for \(J_\nu\) the index \(\xi\) runs through
the true systems in \(b\). The index of the tested
projection satisfies
\[
        \lambda\ \le\ Dv_\ell(M_x)
                      +\sum_{\nu}\min(H_\nu,J_\nu).
\]
In particular every system with \(D_{b(\nu)}=0\)
contributes zero to the second term. This estimate
uses all coefficient primes, including primes at
which some of \(1,q,\ldots,q^{t-1}\) coincide in the
residue field.
\end{theorem}
\begin{proof}
Use ordinary lengths and the uniformizer valuation
over a common splitting field for the proof. At the
end divide all lengths and valuations by \(v_\pi(\ell)\).
By the finite precision argument in
Proposition~\ref{prop:remhecke786}, equality of all
residual characteristic coefficients at good split
places is equality on the whole Galois group.
On a common finite residual quotient this is equivalent
to equality of the semisimplifications by the
Brauer--Nesbitt theorem. In its finite group formulation
one can read this implication from the Brauer characters.
The characteristic polynomial on each \(\ell\)-regular
element gives its eigenvalues with multiplicities, and
the sum of their prime-to-\(\ell\) root-of-unity lifts
gives its Brauer character. Irreducible Brauer characters
are linearly independent. This argument uses polynomial
coefficients at every residue characteristic.

The projection onto each residual class in \(eT\)
preserves \(B=L\cap eT\). To see this select a
distinguishing good split coefficient for each pair
of distinct residual classes. For any one selected
operator, factor the product over its distinct
characteristic-zero eigenvalues of \(Y-\lambda'\)
according to the residues of \(\lambda'\). Factors for
different residues have unit resultant and hence are
comaximal over the valuation ring. The Chinese remainder
idempotents are integral polynomials in the operator.
Their products for the finitely many selected operators
give the claimed projectors. They commute with \(P\):
they preserve both \(Z\) and \(U\) in \(eT=Z\oplus U\).
They also preserve the corresponding intersections with
\(L\). Thus the module \(C=PB/(L\cap Z)\) of the
preceding theorem is the direct sum of modules \(C_b\).
Writing \(B_b\) for the projected \(B\), each \(C_b\)
has both descriptions
\[
  PB_b/(B_b\cap Z)\ \simeq\ B_b/\big((B_b\cap Z)+(B_b\cap U)\big)
             \ \simeq\ (1-P)B_b/(B_b\cap U).
\]
Each isomorphism respects the action of all good
spherical coefficients. In particular \(C_b\) needs at
most \(D_b\) generators.

In the rightmost description filter \(C_b\) by ordering
the competing systems and intersecting with their
successive sums as in Theorem~\ref{thm:heckeflag785}.
Write \(G_j/G_{j-1}\) for the piece belonging to \(\nu\).
It needs at most \(\min(D_b,d_\nu)\) generators.
Each remainder operator kills \(Z\), so it carries
\((1-P)B_b\) into \(B_b\cap U\). On the successive
lattices at \(\nu\) it acts by the indicated remainder
coefficient. Choosing a place and coefficient attaining
\(\delta_\nu\) bounds the normalized length of this
piece by \(H_\nu\).

To obtain the other bound, also filter the leftmost
description of \(C_b\) by successive sums of the true
eigenspaces. Denote this filtration by \(F_i\).
The subquotients
\[
 \frac{F_i\cap G_j}
 {(F_{i-1}\cap G_j)+(F_i\cap G_{j-1})}
\]
filter \(G_j/G_{j-1}\) as \(i\) varies. Indeed the
image of \(F_i\cap G_j\) in that quotient has inverse
image of the previous step equal to the displayed
denominator. By the same argument with the indices
interchanged, each displayed module is a subquotient
of \(F_i/F_{i-1}\) as well. It therefore needs at most
\(\min(D_\xi,d_\nu)\) generators when \(i\) corresponds
to \(\xi\). Here we use the elementary fact about
subquotients of a module over a discrete valuation ring
proved in Theorem~\ref{thm:heckeflag785}.

On \(F_i/F_{i-1}\) each good coefficient acts by its
value at \(\xi\), as one sees by projecting its defining
lattices to the \(i\)-th eigenspace. On \(G_j/G_{j-1}\)
it acts by its value at \(\nu\). Their difference
annihilates the common subquotient. The least valuation
of these differences at good split places is
\(e_{\xi\nu}\), by the same finite precision argument.
The elementary divisor theorem bounds its normalized
length by \(\min(D_\xi,d_\nu)e_{\xi\nu}\).
Summing first over \(i\) bounds the length of the
fixed piece by \(J_\nu\). Use the smaller of the two
bounds on each piece and sum. Passing from \(PB\)
to \(P(L\cap T)\) costs at most \(D v_\ell(M_x)\)
as in Theorem~\ref{thm:heckeflag785}.
\end{proof}

\begin{corollary}[A test at a repeated cyclotomic root]
\leavevmode\label{cor:repeatedroot786}
If \(\bar\kappa(g)=1\) for some \(g\in G_{E_x}\)
and the algebraic multiplicity of \(1\) on
\(\bar u_\nu(g)\) is less than \(t\), the corresponding
piece of \(C\) is zero. More generally it is zero
whenever \(P_\nu(g;Y)\bmod\pi\) is not divisible by
\(Q_g(Y)\bmod\pi\) for some \(g\).
\end{corollary}
\begin{proof}
Reduction commutes with monic long division. A nonzero
remainder in the residue field gives \(\delta_\nu=0\).
If \(\bar\kappa(g)=1\), the divisor is \((Y-1)^t\),
giving the asserted multiplicity test.
\end{proof}
The theorem leaves only residual classes that occur in
both the true lifts and the tested complement. A sum
bound of the required size for their positive depths
is still needed on the common interval.
\clearpage

\subsection{The long character block in the definite target}
\label{sec:longblock787}
Return to the higher-rank descent pairs on the shorter
period interval. Their ranks are \(s=r_x\) and \(n=s+t\),
where \(t=5s\). At a split place use unitary Satake
parameters in this subsection unless the arithmetic
normalization is specified. A block \((\rho,b)\) of
the discrete-spectrum base change consists of a unitary
cuspidal representation \(\rho\) of \({\rm GL}_a\)
over \(E_x\) together with its \(b\) norm twists. It
has rank \(ab\). At an unramified place its entries
are the parameters of \(\rho\) multiplied in turn
by \(q^{(b-1)/2-i}\), \(0\le i<b\).

\begin{theorem}[A character block of exact length on every true lift]
\leavevmode\label{thm:longcharacter787}
The base change of every true system \(\xi\) in the
marked definite target contains the block
\((\mu_W,t)\) exactly once. All the other blocks
together have rank \(s\). At every coefficient prime
there is a continuous semisimple representation
\(R_\xi\) of \(G_{E_x}\) of rank \(s\) for which
\[
             u_\xi \simeq R_\xi\ \oplus\
                     \bigoplus_{a=0}^{t-1}\kappa^{-a}.
\]
Consequently every competing residual semisimplification
that contributes to Theorem~\ref{thm:residualflag786}
contains the indicated residual characters as composition
factors with their full multiplicities.
\end{theorem}
\begin{proof}
Apply the cohomological base change and the description
of the discrete spectrum in the proof of
Theorem~\ref{thm:splitattachment722} to the definite
target. It expresses the good local parameters as
the union over the global blocks just defined. The
unramified parameters of a unitary cuspidal
\({\rm GL}_a\) representation have modulus strictly
less than \(q^{1/2}\). For \(a=1\) their modulus is
one. For \(a>1\) one may use the bound of Luo,
Rudnick and Sarnak with exponent
\(1/2-1/(a^2+1)\), as in
Proposition~\ref{prop:basechange258}.~\cite{X-VII-lrs787}

At all good split places Proposition~\ref{prop:split175}
supplies the entry \(\mu_{W,w}(\varpi_w)q^{(t-1)/2}\).
The other added entries have smaller modulus. The
old entries are the Satake entries of the source
up to a unitary character and the possible
contragredient convention. Their moduli lie between
\(q^{-(s-1)/2}\) and \(q^{(s-1)/2}\): the source is
unitary, and the split spherical estimate in
Lemma~\ref{lem:sphericalcore252} applies to its local
component. Equivalently each block in the spherical
unitary classification of Lemma~\ref{lem:sphericalchain716}
has length at most \(s\). A block shifted by a
complementary exponent \(0<|a'|<1/2\) has a partner
of the same length, so its length is at most \(s/2\).
Both descriptions give the asserted upper bound and,
by duality, the lower bound. Thus the largest modulus
in the target is exactly \(q^{(t-1)/2}\).

If \(b\le t-1\), every entry of \((\rho,b)\) has
modulus strictly less than
\(q^{(b-1)/2}q^{1/2}\le q^{(t-1)/2}\).
A block giving the largest entry must have \(b\ge t\).
Since \(n=s+t<2t\), there can be only one such
block and its cuspidal degree is one. Write
\(\omega\) for its unitary Hecke character. Its top
entry has modulus \(q^{(b-1)/2}\), so \(b=t\).
Every other block has rank at most \(s\). The top
entry of this one block must therefore be the known
entry of the true lift. We get
\(\omega_w(\varpi_w)=\mu_{W,w}(\varpi_w)\) at
all good split places.

For completeness, put \(\eta=\omega/\mu_W\). At
split places outside a finite set the Euler factors
of \(L_E(z,\eta)\) and \(\zeta_E(z)\) agree. At an
inert place over \(v\) of \(E_x^+\) the norm in \(E_x\)
is \((Nv)^2\). Hence the Euler product for their
quotient over the inert places converges absolutely
to a nonzero value near \(z=1\), by comparison with
\(\sum_v(Nv)^{-2\operatorname{Re}z}\). The finitely
many remaining factors are finite and nonzero at
\(z=1\). Thus \(L_E(z,\eta)\) has a pole at one,
and the pole theorem for unitary Hecke characters
gives \(\eta=1\).

The Galois representation in
Theorem~\ref{thm:splitattachment722} is constructed
as a direct sum over the blocks of this base change.
The character block contributes a sum of one-dimensional
representations by class field theory. Its eigenvalues
at geometric Frobenius, after the common multiplication
by \(q^{(n-1)/2}\alpha_w^{-1}\) defining \(u_\xi\),
are \(1,q,\ldots,q^{t-1}\). They give exactly the
direct sum of the indicated cyclotomic powers, by
density of split Frobenius and characteristic-zero
semisimplicity. The other blocks give \(R_\xi\).
Reduce a stable direct sum of lattices in these
summands. The semisimplified reduction is independent
of the lattice by the characteristic-polynomial
form of Brauer--Nesbitt. Theorem~\ref{thm:residualflag786}
now gives the final assertion, also when several
cyclotomic powers reduce to the same character.
\end{proof}

\begin{lemma}[The unique long string at an infinite place]
\leavevmode\label{lem:archstring787}
At every infinite place the holomorphic infinitesimal
parameter of the definite target has exactly one
string of \(t\) consecutive entries of spacing one.
That string cannot be extended at either end.
\end{lemma}
\begin{proof}
Fix a complex label above the given place and put
\(h=s^8,\ k=s+1-m_x,\ a=k+j-1\), with indexing of
the source entries starting at zero. At this label
write \(w_i=w_0^{\rm base}+(2i+1)h\), where
\(w_0^{\rm base}\) is the prescribed nonnegative
base weight. The \(s\) entries \(\sigma_p\) of the
row parameter, in decreasing order, are
\[
 \sigma_p=
 \begin{cases}
 k/2-1-p,&0\le p<k-1,\\
 h-w_{p+2-k},& k-1\le p<s .
 \end{cases}
\]
Up to a common additive shift, its target list is
obtained by inserting, after the first \(a\) entries,
\[
     u_i=-w_j+(t-1)/2-i,\qquad 0\le i<t .
\]
Here is a verification from the weight columns. The
\(\sigma_p+w_j\) for \(p<k-1\) are at least \(h\),
since \(w_j\ge3h\) and \(k\le s\). On the second
range these values are
\((1+2(j-(p+2-k)))h\). They are at least \(h\)
for \(p<a\) and at most \(-h\) for \(p\ge a\).
In particular the indicated insertion is already
in decreasing order. At a column height \(p\) with
\(1\le p<n\), the defining positive-column count
uses the drop from \(\sigma_{p-1}\) to the larger
of \(\sigma_p\) and \(u_0\), minus one, truncated
at zero. The negative-column count is read at the
reflected height \(n-p\); in the lower part it uses
the drop from the smaller of the preceding old
entry and \(u_{t-1}\) to the next old entry, minus
one, again truncated at zero. In particular, above
and below the insertion the counts are the old
successive differences minus one. At height \(a\)
the count is \(\sigma_{a-1}-u_0-1\). At height
\(a+t<n\) it is \(u_{t-1}-\sigma_a-1\). At the
interior heights \(a<p<a+t\) both counts are zero.
These are exactly the consecutive differences of
the inserted list minus one. For a highest weight
in rank \(n\), consecutive differences of the
Harish-Chandra parameter equal the column counts
plus one, from adding the half sum of positive
roots. This verifies the list up to translation,
including when the insertion ends the list. The
opposite complex label reverses and negates the
list, preserving the assertion about strings.

The distance between the inserted block and either
adjacent old entry, when present, is at least
\(h-(t-1)/2>1\) for the ranks in the interval.
There are only \(s<t\) old entries in total, so
they cannot supply a second string of length \(t\).
\end{proof}

\begin{theorem}[A split-place bound on the entire tested complement]
\leavevmode\label{thm:compgap787}
Assume \(Z\ne0\). Every global block \((\rho,b)\)
of every system on \(U=eT\cap Z^\perp\) has \(b<t\).
At each good split place \(w\) as in
Section~\ref{sec:splitcost785} and each such system
the eigenvalue in Proposition~\ref{prop:testedkernel785}
satisfies, in the complex realization of the data,
\[
        (1-q^{-1/2})^n\ \le\ |a_{\nu,w}|\
                                \le\ (1+q^{-1/2})^n .
\]
In particular \(1/2\le|a_{\nu,w}|\le\sqrt e\) if
\(q\ge4n^2\). The bounds are independent of the lift
index and of the twist.
\end{theorem}
\begin{proof}
A block of length \(b\ge t\) has degree one since
\(n<2t\). Call its unitary Hecke character \(\omega\).
At a complex place its holomorphic infinitesimal
entries are consecutive integers apart and centered
at the holomorphic exponent of \(\omega\): the local
norm is \(z\bar z\), and the \(b\) twists have
exponents \((b-1)/2-i\). All the forms in the level
space have the given archimedean weight. By
Lemma~\ref{lem:archstring787} this forces \(b=t\).
Choose a true system in \(Z\). Its long block in
Theorem~\ref{thm:longcharacter787} supplies the unique
string at this same weight. Thus the holomorphic
exponents of \(\omega\) and \(\mu_W\) agree at every
infinite place. For a unitary character of
\(\CC^\times\) written \(z^p\bar z^{p'}\), one has
\(p'=-\overline p\), by evaluating on positive real
numbers and on the unit circle. Consequently their
quotient is trivial at infinity.

This quotient is a finite order Hecke character:
it factors through the ray class group of a finite
modulus because its archimedean component is
trivial. Take a good split place which splits
completely in the normal closure over \(E_x^+\) of
its class field. At this place \(\omega_w=\mu_{W,w}\)
and the whole predicted string occurs. In particular
\(A_w\) kills the system. This contradicts its
invertibility on \(U\) and proves \(b<t\).

Let \(z_i\) be a unitary Satake entry on such a
system. A character block of length \(b<t\) gives
\(|z_i|\le q^{(t-2)/2}\). For a block of cuspidal
degree at least two one has \(b\le n/2=3s\), and
the strict cuspidal bound in the first theorem gives
\(|z_i|<q^{b/2}\le q^{(t-2)/2}\) since \(3s\le5s-2\).
Now the eigenvalue of \(A_w\) in unitary parameters is
\[
          \prod_{i=1}^n
             \left(1-\frac{z_i}
                  {\mu_{W,w}(\varpi_w)q^{(t-1)/2}}\right).
\]
Each fraction has absolute value at most \(q^{-1/2}\).
The triangle and reverse triangle inequalities give
the two bounds. If \(n/\sqrt q\le1/2\), Bernoulli's
inequality gives the last lower bound and
\((1+q^{-1/2})^n\le\exp(n/\sqrt q)\) gives the upper.
\end{proof}
These estimates identify the characteristic-zero blocks
on the period fields. The positive finite-place lengths
between distinct systems in a shared residual class
still enter Theorem~\ref{thm:residualflag786}.
\clearpage

\subsection{Decomposition of the tested index by residue type}
\label{sec:typedindices799}
Fix the original node, twist, lift index \(j\), coefficient
prime \(\ell\), and embedding in Section~\ref{sec:projectionbudget784}.
Let \(\Gamma=\Gamma_x\) be the finite group acting on the
marked frames, and \(M_x=|\Gamma|\). Over a finite splitting
coefficient field write \(e=\sum_{\omega\in\Omega} e_\omega\)
for the marking idempotent, decomposed into its distinct
irreducible character idempotents. A character of this
product group fixes a tuple of residue types at the original
level places. All lattices below are intersections with the
original integral gauge \(L\) at the prescribed level.

\begin{theorem}[The correction for splitting the marked types]
\leavevmode\label{thm:typedindex799}
Let \(T\) be the tested subspace, \(Z\) the true lift range
and \(P\) its orthogonal projection, transported to the
local coefficient field as in Theorem~\ref{thm:projectionindex784}.
Put \(T_\omega=e_\omega T,\ Z_\omega=e_\omega Z\),
\(D_\omega=\dim Z_\omega,\ D=\dim Z\).
Let \(\lambda\) be the normalized length in that theorem
and let \(\lambda_\omega\) be the corresponding index for
\(P:T_\omega\to Z_\omega\). Both use the same embedding
and gauge. Then
\[
        \left|\lambda-\sum_{\omega\in\Omega}\lambda_\omega\right|
                       \le D\,v_\ell(M_x),
                  \qquad v_\ell(\ell)=1 .
\]
In particular equality \(\lambda=\sum_\omega\lambda_\omega\)
holds at every \(\ell\nmid M_x\). For \(D=0\) all these
lengths are zero.
\end{theorem}
\begin{proof}
The frame action is unitary on the definite space. It
preserves \(Z\) by equivariance of the lift. The good
Hecke operators in the false-lift tests commute with the
frame action, and the marking idempotent is central.
Each false space \(eV\cap\ker A_w\cap Z^\perp\), and hence
its orthogonal complement, is stable under \(\Gamma\).
This proves stability of \(T\). In particular \(P\)
commutes with each \(e_\omega\) and \(Pe=P\).
The frame changes and their inverses preserve the integral
gauge even at \(\ell=p_0,p_1\), by the proof of
Proposition~\ref{prop:testedkernel785}. Since
\[
 M_x e_\omega=\sum_{g\in\Gamma}
                   \chi_\omega(1)\chi_\omega(g^{-1})g
\]
has algebraic-integer coefficients, \(M_x e_\omega L\subset L\).
The coefficient extension used to split \(\Gamma\) preserves
normalized lengths of the original lattices by flatness and
the elementary-divisor calculation of
Theorem~\ref{thm:projectionindex784}.

In \(Z\) write \(A=P(L\cap T)\), \(B=L\cap Z\). In
\(Z_\omega\) put \(A_\omega=P(L\cap T_\omega)\),
\(B_\omega=L\cap Z_\omega\). Define the direct sums
\(A_*=\bigoplus_\omega A_\omega\) and
\(B_*=\bigoplus_\omega B_\omega\). They have full rank
\(D\). We have \(B\subset A,\ B_*\subset A_*\), and
\[
                M_x A\subset A_*\subset A,\qquad
                M_x B\subset B_*\subset B.
\]
For the first inclusion take \(a=Pt\) with \(t\in L\cap T\).
Then \(M_x a=\sum_\omega P(M_x e_\omega t)\), with
\(M_xe_\omega t\in L\cap T_\omega\). For the second
chain apply the same identity directly to \(b\in L\cap Z\).
The remaining inclusions follow by summing integral points.
These arguments allow \(T\) to contain the other marked
components on which \(P\) vanishes.

Write \(\ell_p\) here with \(p=\ell\) for normalized local
length. Compute the length of \(A/B_*\) along the two
chains through \(A_*\) and through \(B\). This gives the
exact identity
\[
      \lambda=\sum_{\omega}\lambda_\omega+
                   \ell_p(A/A_*)-\ell_p(B/B_*) .
\]
Each of the last two quotients is generated by at most
\(D\) elements and is killed by \(M_x\). Both lengths
therefore belong to \([0,Dv_\ell(M_x)]\), proving the
claim. In particular the correction involves the dimension
of the true range only once.
\end{proof}

\begin{corollary}[The type comparison on the original period interval]
\leavevmode\label{cor:typedbudget799}
For \(D>0\) set \(g=\lambda\log\ell/(D s)\), \(s=r_x\),
and set \(g_\omega=\lambda_\omega\log\ell/(D_\omega s)\)
when \(D_\omega>0\). For every choice of data at that
prime,
\[
       \left|g-\sum_{D_\omega>0}
                     \frac{D_\omega}{D}g_\omega\right|
                       \le\frac{v_\ell(M_x)\log\ell}{s}.
\]
The sum of these error bounds over all coefficient primes
at a node, including both level primes, divided by
\(B_x^{\rm det}\), is
\(O_{q_0,d}(x^{10}/(b_x(\log b_x)^2))\).
Consequently the sum over primes of the interval supremum
of these normalized errors is \(O_{q_0,d}(V_0^{-389})\).
\end{corollary}
\begin{proof}
The dimension weights sum to one by the direct decomposition
of \(Z\). Multiply the theorem by \(\log\ell/(Ds)\).
Its right-hand side is independent of the trace level, the
coefficients and their embedding. Since
\(\sum_\ell v_\ell(M_x)\log\ell=\log M_x\), the node cost
is bounded by \((\log M_x)/s\le f_xs\log(p_0p_1)\);
this is the group order computation of
Corollary~\ref{cor:markbudget785}. There
\(\log(p_0p_1)\le22\log b_x\),
\(B_x^{\rm det}=f_x x b_x(\log b_x)^3/(m_x k_x)\),
\(m_x\le2\sqrt{x}\) and \(k_x,s\ll x^5\). Substitution
gives the stated node bound. As \(b_x\ge x^{400}\),
it is \(O_{q_0,d}(x^{-390})\). Bound the nonnegative
supremum at each prime by the sum over eligible nodes,
interchange sums and use \(\sum_{x\ge V_0}x^{-390}
\ll V_0^{-389}\).
\end{proof}
This identity permits a separate local trace family for
each fixed tuple of residue types. Its dimension weights
must still be used at the same node and embedding as the
tested projection. The spherical cohomology estimate in
Theorem~\ref{thm:allspherical798} is uniform over those
choices. An estimate for the indices within the individual
types is still needed to bound the full theta projection.
\clearpage

\subsection{The integral projection and the Selmer group of a true system}
\label{sec:markeddescent802}
Fix a node \(x\), a lift index \(j\), a coefficient embedding at
\(\ell\), and its auxiliary twist. Write \(s=r_x,\ t=5s\),
\(n=s+t\). Use \(Z,T,L,P,\lambda\) from
Theorem~\ref{thm:projectionindex784}. The finite group of frame
changes has order \(M_x\). For an individual marked tuple
\(\omega\) put \(e=e_\omega\). At an auxiliary trace level
write \(M\) for the integral lattice of forms, and use the lattice
\(M_\omega=M\cap eM[1/\ell]\) in this paragraph. We work over
a finite coefficient field splitting the finite data in use.

\begin{lemma}[Averaging on the lattice of a single type]
\leavevmode\label{lem:compressednorm802}
The descent of Theorem~\ref{thm:topdescent723} and
Corollary~\ref{cor:baseseparation723} applies on \(M_\omega\)
using \(C_*=M_xc_0\) in place of \(c_0\), where
\(c_0=\gcd_i|\mathrm{GL}_n(\mathbb F_{p_i})|\).
The trace family in this descent can consist just of the systems
occurring in the \(\omega\)-component at the required levels.
\end{lemma}
\begin{proof}
The auxiliary residue group acts at places disjoint from both
original levels and commutes with the old frame group. Each
spherical Hecke sum away from these places commutes with both.
The local translations, the weighted auxiliary norms and the
finite sums defining the parahoric operators at the auxiliary
places commute with \(e\). Their restrictions which map between
integral level lattices therefore map between the corresponding
intersections with \(eM[1/\ell]\). Congruences restrict at full
precision because \(M_\omega\) is saturated in \(M\).
Base invariants in \(M_\omega\) are exactly the intersection of
the original base lattice with the same component. Likewise all
scalar Hecke families used in the construction on this component
act integrally there. For the characteristic coefficients at a
general Galois element this follows by approximation by good
split Frobenii: the restricted integral operator algebra is a
finite submodule of the endomorphisms of \(M_\omega\), hence
closed. This is the integral realization used in
Corollary~\ref{cor:normalizedtransfer722}.

The auxiliary orbit endomorphisms need separate attention. For
every subgroup \(H\) used in Proposition~\ref{prop:righttransfer721},
its orbit construction on \(M\) gives \(\Psi_H\) with
\(\sum_{h\in H}h^{-1}\Psi_Hh=c_0\,1_M\). The same \(c_0\)
works for the subgroup since its point stabilizers are subgroups
of the stabilizers in Lemma~\ref{lem:basestabilizer723}. Use
\[
          \Psi_{H,\omega}=M_xe\Psi_H|_{M_\omega}.
\]
This is an integral endomorphism by the character idempotent
formula in Theorem~\ref{thm:typedindex799}. Its conjugation sum
is \(C_*1_{M_\omega}\). In particular
\(N\Psi_{\Gamma,\omega}z=C_*z\) on base invariants, for the
auxiliary group \(\Gamma\) and its norm \(N\).
These conjugation sums are exactly the orbit identities used for
the principal parahoric and the full group in
Propositions~\ref{prop:righttransfer721} and~\ref{prop:normfactor723}.
The integral Hecke operator on the indicated parahoric fixed
lattice also commutes with \(e\) there, so the same divided
difference and right-translation identities apply to the
intersections. In the induction for \(Y_{S,i}\) use the new
orbit operators and \(\delta=C_*^2D\). Every step then takes
place on \(M_\omega\). The divisibility condition becomes
\((n_0C_*^2)^{|J|+1}\mid m'\). The finite base systems, and
therefore \(n_0\), are fixed before choosing the precision.
Increasing it achieves this condition. In the last division
of Theorem~\ref{thm:topdescent723} these powers cancel. Its
integral conclusion incurs just the factor \(C_*\).
\end{proof}

For a true system \(\xi\) in \(Z_\omega=eZ\) write \(d_{\omega,\xi}\)
for its multiplicity as a dimension of the joint eigenspace
in \(Z_\omega\). Its actual normalized representation is
\(R_\xi\oplus\bigoplus_{a=0}^{t-1}\kappa^{-a}\) by
Theorem~\ref{thm:longcharacter787}. For a stable lattice \(B\)
in \(R_\xi\) put \(A_B=B[1/\ell]/B\). Denote by \(S_\xi\)
the normalized length of the subgroup of \(H^1(E_x,A_B)\)
with the following local conditions. At \(v\mid\ell\) use the
image of \(H^1_f(E_{x,v},B[1/\ell])\). At the original levels
and the twist places away from \(\ell\) allow arbitrary classes.
At every other finite place require an unramified class. An
infinite length is given value \(+\infty\). Here the two original
levels and the twist are those of this very projection.

\begin{proposition}[Clearing the projection on one marked component]
\leavevmode\label{prop:markedclear802}
Let \(D_{x,\ell}^{\rm loc}\) be the supremum of the sum of the
finite-precision local defects of Section~\ref{sec:highsmall789}
over \(v\mid\ell\) and the \(t+1\) target summands, for individual
marked families. It may include a separate choice of designated
true system, precision and auxiliary trace level for each value
of the supremum. Write \(Z_x\) for the positive integer of
Section~\ref{sec:projectedcups742}, with \(t_0=t\), and let
\[
              E_{\rm tw}=t\sum_{w\in Q\cup Q^c}
                                v_\ell(1-\zeta_w)
\]
at the new twist places, with the notation of
Corollary~\ref{cor:alltwistcost738}. If \(D=\dim Z>0\), then
\[
 \frac{\lambda}{D}\ \le\
       \sum_{\omega,\xi}\frac{d_{\omega,\xi}}{D}\,S_\xi
                    +D_{x,\ell}^{\rm loc}+v_\ell(Z_x)+E_{\rm tw}.
\]
Only pairs with positive multiplicity enter the sum. All local
lengths and valuations have \(v_\ell(\ell)=1\).
\end{proposition}
\begin{proof}
Fix \(\omega,\xi\) and suppose the Selmer lengths appearing
on the right are finite. At an original prime \(q\ne\ell\)
the tuple has a primitive member at some place. The first
lemma of Section~\ref{sec:markedlevel801} determines the smooth
type there just as at a coefficient place. Split-place
compatibility away from \(\ell\) in Theorem~\ref{thm:splitattachment722}
identifies it with actual inertia, since monodromy is zero.
The \(t\) trivial characters on inertia are precisely the
ones contributed by the Tate lines. Thus \(R_\xi\) has the
primitive degree-\(s\) orbit. The auxiliary twist is unramified
at this place. Theorem~\ref{thm:markedmatrix715} supplies a
tame element \(\eta\) acting by a scalar \(y\) with \(y-1\)
a unit, and the full integral matrix span on \(B\). It also
shows that any stable lattice is homothetic to \(B\): the
integral diagonal matrix units split it into fractional ideals
and the other units identify these ideals. Its Selmer length
is consequently independent of this choice. These scalar
hypotheses hold for every true system in the component.
Corollary~\ref{cor:markeddet715} supplies the auxiliary primes
for the designated system with unit determinant at the chosen
Frobenii, at any prescribed precision. The common polynomials
and the two trace identities at each pair are then those of
Theorem~\ref{thm:commonhecke716} and
Corollary~\ref{cor:normalizedtransfer722}.

Take the row construction of Theorem~\ref{thm:ribettrace720}
on this family. We record the finite-level estimate on its
main determinant. In Proposition~\ref{prop:blockrows719} the
annihilator of the row lattice, at precision \(\ell^\nu\), maps
injectively to normalized raw cocycles with values in
\(A_B[\ell^\nu]\). They vanish at \(\eta\) and at the two
auxiliary tame generators at each trace pair. At an auxiliary
trace place the target is unramified. Its restriction to
inertia is an additive homomorphism on the pro-\(\ell\) tame
quotient and vanishes on its generator. At any further good
place the matrix relations themselves force the inertia
values to vanish, since all the matrices there act trivially.
Localize these cocycles at the coefficient places and divide
their images by the reductions of the integral finite local
conditions. The kernel maps into the Selmer group defining
\(S_\xi\) by division by \(\ell^\nu\). A coboundary in \(A_B\)
vanishing at \(\eta\) is zero because \(y-1\) is a unit. The
map on this kernel is therefore injective. Its quotient has
normalized length at most the sum of the local defects for
\(B\). For a scalar target, the quotient defining \(l_a\)
in Theorem~\ref{thm:charrows719} injects into the corresponding
sum of local quotients by restriction. Indeed these integral
evaluation equations characterize the reductions of local finite
cocycles by Lemma~\ref{lem:localevaluation728}, and raw
coboundaries over the finite coefficient ring lift integrally.
All these estimates use the same trace family, including its
precision. Their total is bounded by \(S_\xi+D_{x,\ell}^{\rm loc}\).

The scalar row costs at the original levels away from \(\ell\)
are at most \(2v_\ell(G_0)\) by Corollary~\ref{cor:levelcost725}.
The new twist places contribute \(E_{\rm tw}\). The class and
auxiliary terms of Theorem~\ref{thm:ribettrace720} are bounded by
\[
    v_\ell(h_E\mathcal J_{N,t})+
                (4t+3)(N+t+1)^6 v_\ell(G_0),
           \qquad N=q_{1,x}d,
\]
by Corollary~\ref{cor:actualmarks714}. These are the remaining
costs at the designated eigenvalue of the top trace coefficient.

Use Corollary~\ref{cor:baseseparation723} on the intersection
lattice of the lemma. For \(l\in L\cap eT\) take \(z=Pl\).
It is supported on the true systems. On any base system with a
trivial Galois composition factor, the arithmetic Hecke polynomial
at every good split place has its root \(\alpha_w\), because the
normalized Galois polynomial has root one there. Since
\(|\alpha_w|^2=q_w^s\), Lemma~\ref{lem:sphericalchain716} gives
the root \(\alpha_wq_w^{t-1}\) as well. Thus \(A_w\) of
Proposition~\ref{prop:testedkernel785} kills that eigenspace.
The projection formula there gives equality of the components of
\(l\) and \(z\) on this system. More explicitly, the joint
eigenspace projector commutes with \(e\) and \(P\), and its
component of \(l-z\) belongs to the tested space, to \(Z^\perp\)
and to the kernel of \(A_w\). Their intersection in \(eT\) is
zero. This verifies the base agreement even in an eigenspace
of dimension greater than one. The lemma and base separation
now give a diagonal operator \(s_\xi\) on \(Z_\omega\) such
that \(C_*s_\xi Pl\in L\cap Z_\omega\) for all such \(l\).
At the designated eigenspace its eigenvalue has valuation at
most the top coefficient bound just obtained. With integral
\(l\) already in the component the factor \(h\) in that
corollary is one.

Repeat the construction for each distinct true system on
\(Z_\omega\), allowing its own trace precision. The finitely many
operators are defined over one finite extension. One integral
linear combination simultaneously attains the minimum of their
eigenvalue valuations at each true system: divide the coefficients
of each resulting linear form by one of minimal valuation, then
choose the combining scalars with residues outside the finitely
many zero hyperplanes. A finite unramified extension supplies
such residues because their union is proper over the algebraic
closure of the residue field. The new operator still clears the
same projected lattice after multiplication by \(C_*\). Its
determinant valuation is the sum of its eigenvalue valuations
with multiplicities. The exact minimum in
Theorem~\ref{thm:projectionindex784} bounds \(\lambda_\omega\)
by this value, with the factor \(C_*\) counted
\(\dim Z_\omega\) times. Coefficient extension preserves the
index. Sum over the tuples and apply Theorem~\ref{thm:typedindex799}.
The additional normalized cost per dimension is at most
\(v_\ell(C_*M_x)\). Since \(s,n\le st\), block embeddings
of the finite general linear groups give \(M_x\mid G_0\) and
\(c_0\mid G_0\). Thus \(C_*M_x\mid G_0^3\). The two old level
rows and these integrality factors together cost at most five
powers of \(G_0\), included in \(Z_x\). This proves the result.
If any relevant global length is infinite the inequality holds
in the extended nonnegative reals.
\end{proof}

\begin{theorem}[Comparison on the period interval]
\leavevmode\label{thm:markedperiod802}
On the marked levels of Corollary~\ref{cor:markedlocal801},
write \(\overline S=\sum_{\omega,\xi}d_{\omega,\xi}S_\xi/D\)
for a positive-dimensional true range with finite indicated
lengths. The sum with weights \(\log\ell\) of the primewise
suprema of
\[
             \frac{(\lambda/D-\overline S)_+}{sB_x^{\rm det}}
\]
on the shorter interval is
\(O_{q_0,d}((\log(2V_0))^{-2})\). The supremum uses the same
node, lift index, embedding and twist for the two terms of
each difference. Use zero for an empty supremum or a zero range.
\end{theorem}
\begin{proof}
The coefficient places in these levels are the spherical ones
and the two original levels. Indeed the compact level imposes
the frame conditions at the original places and parahoric
conditions at the twist places. Away from these it uses the
standard maximal compact for the self-dual lattice. The parity
factor in the normalizing splitting character is unramified
on finite units. The twist pairs and the later trace pairs
are chosen away from \(\ell\). At an inert place
\(E_x/E_x^+\) is unramified, so the standard self-dual lattice
there gives a hyperspecial compact, as in
Lemma~\ref{lem:inertcrys789}. Corollary~\ref{cor:markedlocal801}
therefore gives precisely the required estimate for
\(D_{x,\ell}^{\rm loc}\).

At each node \(\sum_\ell v_\ell(Z_x)\log\ell=\log Z_x\).
The calculation in Theorem~\ref{thm:projectedcups742} gives
\(\log Z_x\ll_{q_0,d}x^{110}(1+\log b_x)\). Here
\(B_x^{\rm det}=f_xxb_x(\log b_x)^3/(m_xk_x)\),
\(m_x\le2\sqrt x,\ k_x\le s+1,\ f_x\ge1\) and \(b_x\ge x^{400}\).
Consequently \(\log Z_x/(sB_x^{\rm det})\ll_{q_0,d}x^{-289}\).
The sum of the normalized twist-place costs with weights
\(\log\ell\) at each node can be taken at most one by
Corollary~\ref{cor:alltwistcost738}, during the same choice as
for the critical ratios. For these last two nonnegative errors
replace each primewise supremum by the sum over eligible integer
nodes. The resulting tail is \(O_{q_0,d}(V_0^{-288})\).
Combining the preceding proposition with the local bound proves
the comparison.
\end{proof}
To obtain a bound for the projection itself one must still control
\(\overline S\) for the actual true systems with the twists used
for the periods. A bound for Selmer groups of independently
chosen twisted representations would not verify that condition.
\clearpage

\subsection{The source representation and the common period denominator}
\label{sec:sourceperiod803}
At a node of Corollary~\ref{cor:actualmarks714} let \(s=r_x\),
\(t=5s\) and \(n=s+t\). The index \(j\) runs from \(1\) to \(m_x\).
The character of the norm-one torus used in the \(j\)-th auxiliary
lift before changing the finite level will be denoted by \(\mu_j\).
For any character \(\nu\) of this torus put
\(\widetilde\nu(a)=\nu(a/\bar a)\) on the ideles of \(E=E_x\).
Write \(\vartheta_\ell\) for the Galois character of
\(\widetilde\vartheta\) when \(\vartheta\) has finite order. We use
geometric Frobenius, so \(\kappa(\mathrm{Frob}_w)=q_w^{-1}\) for
\(w\nmid\ell\), where \(q_w\) is the residue cardinality. In
particular the unramified norm character \(|\cdot|_E\) has avatar
\(\kappa\).

\begin{proposition}[The main summand of an actual twisted lift]
\leavevmode\label{prop:sourcetwist803}
Fix a coefficient prime and embedding. There is a finite set
\(T_*\) of auxiliary primes and a finite list of rank-\(s\)
representations \(V_{\pi,j}\), both chosen before selecting the
twists for the Selmer estimates, with the following property.
If the conductor of \(\vartheta\) avoids \(T_*\) and the lift
using \(\mu_j\vartheta^j\) has a true system
\(\xi\), its summand in Theorem~\ref{thm:longcharacter787} satisfies
\[
                  R_\xi\simeq V_{\pi,j}\otimes\vartheta_\ell^{-j}
\]
for some member of the list. Each \(V_{\pi,j}\) is de Rham at
the coefficient places and unramified away from these places and
the two original levels. Its polarization is
\(V_{\pi,j}^c\simeq V_{\pi,j}^{\vee}\kappa^{1-t}\).
\end{proposition}
\begin{proof}
For index \(j\) take the actual source on \(W_j\). Its fixed
archimedean discrete series and finite principal level are
independent of the twist; the marking is an idempotent in its
finite level group. There are finitely many spherical systems
\(\pi\) in this source. Indeed the space of multiplicities at
this level for the prescribed real type is finite dimensional by
Harish-Chandra finiteness. As in the proof of
Proposition~\ref{prop:heckecore252}, evaluation on a cyclic
compact-finite vector realizes it in the space of automorphic
forms at fixed level, infinitesimal character and compact type.
The center is the norm-one torus and has compact adelic quotient.
The commuting spherical Hecke algebras are closed under
adjoints on the unitary discrete spectrum, so their restrictions
to the finite multiplicity space are simultaneously diagonalizable.
The frame idempotent commutes with them. Linearity of the lift
therefore expresses its range as a sum of images of the marked
source eigenspaces. The spherical local theta correspondence
specifies the good split parameters on each nonzero image by
Proposition~\ref{prop:split175}. This argument uses \(W_j\) for
each \(j\), including its indefinite signatures.

At a good split place let \(B_w\) denote the Satake
multiset of a fixed source system \(\pi\) in unitary normalization.
Write \(\chi_S,\chi_T\) for the two splitting characters
and evaluate them at the chosen uniformizer. Here
\(\chi_T=\widetilde{\mu_j\vartheta^j}(w)\). Indeed, if
\(\upsilon\) is the fixed splitting character, the correction
on the norm-one torus is \((\upsilon|_{U(1)})^{n/2}\). Its
pullback is \(\upsilon^n\), since \(\upsilon^c=\upsilon^{-1}\).
The definition of the rank-\(n\) space character multiplies by
\(\upsilon^{-n}\), giving the asserted equality.
The good-place parameter formula of Proposition~\ref{prop:split175}
now gives the multiset
\[
       (\chi_S/\chi_T)B_w\ \sqcup\
       \chi_S\{q_w^{-(t-1)/2+i}:0\le i<t\}.
\]
The arithmetic normalization in
Corollary~\ref{cor:normalizedtransfer722} multiplies this list by
\(q_w^{(n-1)/2}/(q_w^{s/2}\chi_S)\). The second part becomes
\(\{q_w^i:0\le i<t\}\). Cancelling the Tate factors using
Theorem~\ref{thm:longcharacter787}, the eigenvalues of
\(R_\xi\otimes\vartheta_\ell^j\) are precisely
\[
                 q_w^{(t-1)/2}B_w/\widetilde\mu_j(w)
\]
in the chosen embedding. Thus for fixed \(\pi,j\) any two
nonzero lifts, possibly for different twists and different
marked vectors, give the same semisimple representation after
this character twist.
These split places suffice for equality of continuous semisimple
representations over \(E\). In any finite Galois quotient over
\(E\), take a Galois closure over \(E^+\) and lift a specified
element to the subgroup fixing \(E\). Chebotarev over \(E^+\)
gives a prime split in \(E\) and a choice of place with Frobenius
in the prescribed class. Finitely many excluded places do not
affect this argument. Equality of traces on this dense set gives
the asserted isomorphism by characteristic-zero Brauer--Nesbitt.

We explain how to arrange unramifiedness outside the old set
without requiring a nonvanishing statement at the trivial twist.
Consider the finite list of pairs \((\pi,j)\). Use as candidate
twists all the permitted \(m_x\)-fold finite-order torus characters
supported at new split pairs, prime to the old levels and trivial
at the places above \(\ell\). For each pair which gives a nonzero
lift at some candidate, choose one such witness \(\beta_{\pi,j}\)
and a true system in its image. Let \(T_*\) be the finite union
of the rational primes in the conductors of these witnesses.
Retain a pair if it has a witness \(\delta_{\pi,j}\) whose
conductor avoids \(T_*\), choosing one such witness. Every true
system for a final twist avoiding \(T_*\) comes from a retained
pair, by the definition of retention.

For a retained pair take \(V_{\pi,j}\) to be the common
semisimple representation just identified using its two witnesses.
The true summand in each definite target is de Rham and polarized
with exponent \(1-t\) by Proposition~\ref{prop:highwidth788}
and its proof. Finite-order twisting preserves the de Rham
condition and anticyclotomic twisting preserves the polarization.
Fix a place away from \(\ell,p_0,p_1\). At least one of the two
witnesses is unramified throughout its rational residue
characteristic, since their conductors are disjoint. For that
witness the target is spherical there and its normalizer
unramified. Its local base change is unramified at a split place
by Theorem~\ref{thm:splitattachment722}, and at an inert place by
the hyperspecial base change used in Lemma~\ref{lem:inertcrys789}.
The local parameter has zero monodromy, so compatibility away
from \(\ell\) gives an unramified Galois representation. Its
summand and its unramified twist are unramified. This proves the
assertion at every such place for \(V_{\pi,j}\).
\end{proof}

\begin{lemma}[One old inertia test for the entire list]
\leavevmode\label{lem:sourcebudget803}
Let \(G_0=\prod_{i=0}^1|\mathrm{GL}_{st}(\mathbb F_{p_i})|^{f}\),
where \(f=[E^+:\mathbb Q]\). For the source list of the proposition,
the constants \(B_i\) of Theorem~\ref{thm:twistedselmer730} are
at most \((4s+1)v_\ell(G_0)\). Its distinguished generator \(g_*\)
may be taken at any one prescribed old marked place away from
\(\ell\), simultaneously for this list and the torus tests of
Theorem~\ref{thm:rankoneselmer737}.
\end{lemma}
\begin{proof}
Use a true system for a witness of the preceding proposition.
Choose a marked tuple in which it has a nonzero component.
For each old place over \(p_i\ne\ell\),
Theorem~\ref{thm:markedtype801} and its proof identify
the smooth target parameter there. Its old-frame calculation uses
only the prescribed tuple, the old determinant character and the
fact that the witness is unramified at the old primes, all of
which hold for the candidates used above. The parameter has zero monodromy, the
\(s\) shifted roots of the degree-\(s\) label \(c\), and \(t\)
copies of the trivial character. The shift is by a character
of order dividing \(Q\) of Section~\ref{sec:markedmatrix715}.
An orbit of size greater than one keeps its size, and no allowed
degree-one root has order dividing \(Q\). Thus all \(s\) shifted
roots are nontrivial. The full \(t\)-dimensional trivial part
on inertia is already accounted for by the Tate lines.
Split-place compatibility away from \(\ell\) and the vanishing
of monodromy identify the other roots with the actual inertia
eigenvalues of the true summand. Twisting back by the witness
has no effect on inertia at the old places. At a primitive
member of the tuple the shifted root remains primitive by
Lemma~\ref{lem:tameorbit715}.
The conjugate place has the inverse inertia eigenvalues, by
polarization. Each tuple has a primitive member above each
of \(p_0,p_1\). For every representation choose such a member
away from \(\ell\). Theorem~\ref{thm:markedmatrix715} gives an
inertial scalar \(y\) with \(y^2-1\) a unit and the full matrix
span on a stable lattice. This marked scalar may depend on
the representation. The possibly different generator \(g_*\)
in the lemma is common: every eigenvalue at every old place
is nontrivial, including at \((g_*^c)^{-1}\). The rank-one tests
also have this property there by Theorem~\ref{thm:rankoneselmer737}.

Put \(g_i=|\mathrm{GL}_{st}(\mathbb F_{p_i})|\).
A root \(\lambda\ne1\) from an orbit of degree \(e\le s\)
has order dividing \(p_i^e-1\), also after the scalar shift,
since \(Q\mid p_i-1\). The polynomial \(X^{p_i^e-1}-1\)
splits into roots of unity. Dividing by \(X-1\) and evaluating
at one shows \(v_\ell(1-\lambda)\le v_\ell(p_i^e-1)\).
The integer \(p_i^e-1\) divides \(g_i\) by the order formula
for the finite general linear group. The determinant valuation at
a full tame generator is therefore at most \(s v_\ell(g_i)\).
At a noncoefficient old place the local dual invariants at every
precision are in the kernel of this generator minus one on the
Tate dual. Its determinant has the same valuation. The
elementary divisors bound the normalized length of the kernel
by that valuation. Summing over at most \(2f\) places of \(E\)
over each \(p_i\) gives \(2s v_\ell(G_0)\). There are no further
old terms if \(S\) consists of these places and the coefficient
places. The two distinguished determinants together cost at most
\(2s v_\ell(G_0)\). Finally Corollary~\ref{cor:congcost728}
gives \(v_\ell(I_t)\le v_\ell(G_0)\) using a primitive member
away from \(\ell\). Adding the terms gives the bound.
\end{proof}

\begin{theorem}[Selmer lengths for the summands in the projected periods]
\leavevmode\label{thm:actualselmer803}
Fix a node and positive tolerances \(\epsilon_\ell\) with
\(\sum_\ell\epsilon_\ell\log\ell\le1\). The twists for the
critical ratios and the torus periods may be chosen so that
every nonzero true higher-rank system on these same twists satisfies
\[
                 S_\xi\le (4s+1)v_\ell(G_0)+\epsilon_\ell .
\]
Here \(S_\xi\) is the divisible Selmer length used in
Proposition~\ref{prop:markedclear802}.
\end{theorem}
\begin{proof}
At each embedding take the finite list of lattices of
Proposition~\ref{prop:sourcetwist803} for all \(j\), together
with the rank-one list. Enlarge the coefficient field to contain
their finite data. The polarization with \(t\ge1\), the de Rham
condition, the full matrix span and the unit scalar verify the
hypotheses of Theorems~\ref{thm:balancedcount728}
and~\ref{thm:tiedlength729}. The preceding lemma checks the two
determinants for the chosen \(g_*\) of
Theorem~\ref{thm:criticalheights734}. Apply that theorem to this
joint list and the critical pairs at the node, excluding \(T_*\)
from all new conductor primes. This exclusion is legitimate in
the Chebotarev choices of Theorems~\ref{thm:tiedlength729}
and~\ref{thm:commoncritical731}: deleting a fixed finite set
of rational primes leaves every required Frobenius class available,
including at the first pair and at the smoothing prime. The group
of allowed global cocycles still imposes the unramified condition
at the excluded auxiliary places. Thus this choice adds no term
to \(B_i\). The construction selects
one character at the embedding for all the signed powers; it has
the simultaneous divisible-coefficient and new-inertia estimates
of Theorem~\ref{thm:twistedselmer730}. Its ray characters descend
to the norm-one torus as in the construction of
Theorem~\ref{thm:commoncritical731}. The critical ratio and
new-inertia tolerances may be decreased to meet
Corollaries~\ref{cor:quotientchainheight734}
and~\ref{cor:alltwistcost738} for the finitely many powers.
Taking the power \(-j\) identifies the twisted lattice up to
homothety by Proposition~\ref{prop:sourcetwist803}. Indeed every
stable lattice is homothetic by the full integral matrix span,
which a twist by scalar units preserves.

To see that the bound applies to the precise Selmer group, let
\(A=B[1/\ell]/B\) for a stable lattice \(B\) in this twisted
representation. A class has a unique cocycle representative
vanishing at the chosen scalar element: subtract the coboundary
of \((y-1)^{-1}z(\eta)\). The twist is unramified there.
Away from \(S\) and the new pairs, the action is unramified,
so a locally unramified class has zero raw restriction to inertia.
At \(v\mid\ell\) the twist is trivial. A class in the image of
\(H_f^1(E_v,B[1/\ell])\) has a rational finite cocycle whose
reduction in \(A\) agrees with the actual local cocycle. To
obtain equality on cocycles, lift the element of \(A\) giving
their difference as a coboundary to \(B[1/\ell]\) and adjust
the rational representative. Continuous cocycles on a compact
Galois group have bounded image in the rational representation
and finite image in the discrete module \(A\). Choose a power
\(\ell^k\) clearing all the local rational values and killing the
global image. Under \(A[\ell^k]\simeq B/\ell^kB\), each local
restriction is then a reduction of an integral finite cocycle.
It satisfies the evaluation equations of
Lemma~\ref{lem:localevaluation728}. There is no restriction at
the remaining old or new pairs. Thus normalization embeds the
entire Selmer group into the bounded group of cocycles of
Theorem~\ref{thm:twistedselmer730}. This proves the inequality,
including after extending the coefficient field.
\end{proof}

\begin{corollary}[A common determinant and torus denominator on the period interval]
\leavevmode\label{cor:actualdenom803}
On the shorter interval with fixed \(q_0,d\), make the
simultaneous choices of the theorem at each node. For the
actual true projection let \(D=\dim Z\), with zero quotient
by convention when \(D=0\). Then
\[
    \sum_\ell(\log\ell)
       \sup_{x,j,\iota}
          \frac{\lambda}{Ds B_x^{\rm det}}
           \ll_{q_0,d}(\log(2V_0))^{-2}.
\]
For all sufficiently large \(V_0\) the integral operator and
all main-pair period tests in Corollary~\ref{cor:remainingbudget784}
admit common bounds of total size at most \(B_x^{\rm det}\)
at each \(x,j\) of the interval.
\end{corollary}
\begin{proof}
The multiplicities \(d_{\omega,\xi}\) sum to \(D\).
The same bound on \(S_\xi\) thus holds for their weighted mean
\(\overline S\) in Theorem~\ref{thm:markedperiod802}.
At a node its contribution, after the sum over primes with
weights \(\log\ell\), is bounded by
\[
         \frac{(4s+1)\log G_0+1}{sB_x^{\rm det}} .
\]
Here \(B_x^{\rm det}=fxb_x(\log b_x)^3/(m_x k_x)\), and
\(\log G_0\le f(st)^2\log(p_0p_1)\) by
Corollary~\ref{cor:actualmarks714}. Since \(\log(p_0p_1)
\ll\log b_x,\ s\le3x^5,\ t=5s,\ m_x\le2\sqrt{x}\) and
\(k_x\le s+1\), the quotient is
\(O_{q_0,d}(x^{25}/(b_x(\log b_x)^2))\).
For a sum of nonnegative errors the supremum over nodes is at
most their sum. Summing the last bound over integer nodes
\(x\ge V_0\), with \(b_x\ge x^{400}\), gives
\(O_{q_0,d}(V_0^{-374})\). Now apply
Theorem~\ref{thm:markedperiod802} to the positive excess over
the same mean for each twist. This proves the first assertion.
The sum of \(h_{x,j,\ell}\) of
Corollary~\ref{cor:remainingbudget784} is at most \(B_x^{\rm det}\)
times this interval sum, so eventually it is at most
\(B_x^{\rm det}/2\). That corollary and
Proposition~\ref{prop:torusbudget784}, on the very same twists,
give the stated common denominator.
\end{proof}
This denominator controls projection onto the true lift range.
To reach the prime-power estimate one still needs the comparison
of the determinant of the composite theta lift with the local
divisibility gain for the same family of maps.
\clearpage

\subsection{Minors of the projected theta-period pairing}
\label{sec:testedminors804}
Let \(Z\) be the true range of the higher-rank theta lift at a
fixed node, index and twist. Write \(P\) for its orthogonal
projection and \(T\) for the tested subspace of
Theorem~\ref{thm:projectionindex784}. Use the Hermitian pairing
on the definite target which is linear in its second variable
and gives the defining compact level volume one. For two lists
of tested vectors, of equal length \(h\), put
\[
   H(\mathbf y,\mathbf z)=
       \det\bigl(\langle P y_i,P z_k\rangle\bigr)_{1\le i,k\le h}.
\]
We determine how the index of the integral projection enters
this minor, and how its complex conjugates are bounded by the
norms of the original test vectors.

\begin{proposition}[The tested projection at every complex embedding]
\leavevmode\label{prop:archcontraction804}
For an automorphism \(\sigma\) of \(\mathbb C\) let \(\|v\|_\sigma\)
be the class-sum norm of the definite form after applying \(\sigma\)
to its coefficient coordinates. On \(T\) the transported map \(P\)
is the restriction of orthogonal projection for this norm onto
the transported subspace \(Z\). In particular the map and all its
exterior powers are contractions on the respective tested spaces.
\end{proposition}
\begin{proof}
We specify the Hermitian structure at \(\sigma\). Fix level class
representatives and an integral tensor realization of the Weyl
coefficient space. Apply \(\sigma\) to the coordinates of the
values and sum their Hermitian pairings over classes, with weight
the reciprocal of the stabilizer order. The Weyl Gram matrix
is the Gram matrix of integral tensors in a tensor basis with
identity Gram matrix. It has rational entries and is positive
definite. This construction gives an ordinary positive Hermitian
space on the transported coordinates. The class pairing is
independent of the rational representatives. Indeed a rational
unitary matrix acts unitarily on each tensor factor belonging
to a complex embedding \(\tau\) of \(E_x\). After transport the
factor belongs to \(\sigma\circ\tau\); this embedding, too, takes
the CM involution to ordinary complex conjugation on its image.
The induced action on the Weyl space, including its determinant
power, is unitary. Finite compact character factors are roots
of unity at every conjugate.

Choose one split spherical place \(w\) off the level and use
\(e,A_w\) of Proposition~\ref{prop:testedkernel785}. In the
transported structure the frame action is unitary: translation
permutes the classes with their stabilizer orders, with the
unitary actions on values just described. A central idempotent
in a finite complex group algebra acts as identity or zero on
each irreducible summand of a unitary representation. Its image
and kernel are therefore orthogonal; the same applies to the
central idempotent after applying \(\sigma\) to its coefficients.
At \(w\), the adjoint of a double-coset operator is the inverse
double-coset operator. To see this, lift both vectors to the
intersection of the two compact levels, and change variables
by right translation in the class integral. Haar measure is
right invariant and the two indices are equal by unimodularity.
With level volume one this gives precisely the sums over right
cosets defining the two operators. The calculation also applies
to transported values. The scalar coefficients of the adjoint
polynomial are the ordinary conjugates of the transported
coefficients. Both polynomials belong to the commutative spherical
algebra, so the transported \(A_w\) is normal. Its image is
orthogonal to its kernel.

Before transport, \(eZ=Z\), \(A_wZ=0\), and \(A_w\) is invertible
on \(eT\cap Z^\perp\) by Proposition~\ref{prop:testedkernel785}.
For \(y\in T\) decompose its remainder as
\[
       y-Py=(1-e)y+(ey-Py)
                       \ \in\ \ker(e)+\operatorname{im}(A_w).
\]
Indeed the second summand belongs to \(eT\cap Z^\perp\). Every
vector of the subspace on the right is orthogonal to \(Z\)
after transport by any \(\sigma\), by the preceding two
orthogonality statements and the inclusions
\(Z\subset\operatorname{im}(e)\cap\ker(A_w)\).
Thus the transported \(Py\) is exactly the orthogonal
projection of the transported \(y\). In an orthonormal basis
starting with a basis of the image, an orthogonal projection
retains the indicated coordinates. In the corresponding
exterior basis its action likewise retains a subset of the
coordinates. This proves both contraction assertions.
\end{proof}

\begin{proposition}[One local projection loss for a period minor]
\leavevmode\label{prop:extprojection804}
Use a finite coefficient field at \(\iota\) with uniformizer
\(\pi\), and let \(e_\ell=v_\pi(\ell)\). Put \(D=\dim Z\) and
write \(a_1\ge\cdots\ge a_D\ge0\) for the elementary divisor
exponents of \(L_Z=Z\cap L\) in \(P(L_T)\), where \(L_T=T\cap L\).
For \(1\le h\le D\) put
\(\lambda_h=e_\ell^{-1}\sum_{i=1}^h a_i\), so
\(\lambda_h\le\lambda_D=\lambda\). Suppose each \(z_k\)
is integral at \(\iota\) and each \(y_i\) is integral at
\(\iota\circ c\), where \(c\) is ordinary complex conjugation.
In the class-sum normalization, for
\(c_0=\gcd_{i=0,1}|\mathrm{GL}_{s+t}(\mathbb F_{p_i})|\),
\[
      \log\max\{1,|\iota H(\mathbf y,\mathbf z)|_\ell\}
         \le \bigl(\lambda_h+h\,v_\ell(c_0)\bigr)\log\ell .
\]
The elementary divisors are understood after adjoining the
finitely many coefficient values used in the calculation.
\end{proposition}
\begin{proof}
Each row defines a linear functional on \(Z\) at \(\iota\),
by \(u\mapsto\iota(c_0\langle y_i,u\rangle)\), with the scalar
coordinate of \(u\) transported by \(\iota\). This functional
is integral on \(L_Z\). In fact the coefficient pairing in the
class sum uses the conjugate of the first coordinate array and
has an integral Gram matrix. The changes of basis used by the
local gauges on the two inputs cancel in this pairing by
\(\bar g^{\,t}g=1\) at the two conjugate components. This is
the pairing calculation in Theorem~\ref{thm:rankonepair736}
and uses only the integral tensor realization on the definite
target. Thus every summand before division by its stabilizer
order is integral. Each stabilizer order divides \(c_0\) by
Lemma~\ref{lem:basestabilizer723}. Multiplication by \(c_0\)
clears all the summands together.

The columns \(Pz_k\) belong to \(P(L_T)\). Orthogonality gives
\(\langle y_i,Pz_k\rangle=\langle Py_i,Pz_k\rangle\).
Choose a basis \(u_1,\ldots,u_D\) of \(P(L_T)\) for which
\(\pi^{a_i}u_i\) form a basis of \(L_Z\). The inclusions of
these lattices and the length equality \(\lambda_D=\lambda\)
are the projection exact sequence of
Theorem~\ref{thm:projectionindex784}. An exterior basis vector
\(u_{i_1}\wedge\cdots\wedge u_{i_h}\) becomes integral in
\(\bigwedge^h L_Z\) on multiplication by
\(\pi^{a_{i_1}+\cdots+a_{i_h}}\). The largest such exponent
is \(\sum_{i=1}^h a_i\). Consequently the wedge of the \(h\)
integral functionals evaluated on
\(Pz_1\wedge\cdots\wedge Pz_h\) has valuation at least
\(-\lambda_h\) in the normalization \(v_\ell(\ell)=1\).
This evaluation is \(c_0^h\iota H(\mathbf y,\mathbf z)\).
It proves the estimate. Extending the coefficient field
multiplies both the exponents and \(e_\ell\) by its relative
ramification index, preserving \(\lambda_h\).
\end{proof}

\begin{theorem}[Archimedean and finite bounds on the same projected periods]
\leavevmode\label{thm:actualminors804}
Take the shorter fields and the twists of
Corollary~\ref{cor:actualdenom803}. For the complex bounds put
\(\sigma^\#=c\circ\sigma\circ c\). At every node and index, for
any of its twists and all its tested vectors,
\[
 \begin{split}
 |\sigma H(\mathbf y,\mathbf z)|
  &\le \|y_1\wedge\cdots\wedge y_h\|_{\sigma^\#}
         \|z_1\wedge\cdots\wedge z_h\|_\sigma\\
  &\le \prod_{i=1}^h\|y_i\|_{\sigma^\#}
         \prod_{k=1}^h\|z_k\|_{\sigma}.
 \end{split}
\]
On the same interval one has
\[
  \sum_\ell\ \sup_{\substack{x,j,\iota,h\\\mathbf y,\mathbf z}}
     \frac{\log\max\{1,|\iota H(\mathbf y,\mathbf z)|_\ell\}}
          {D s B_x^{\rm det}}
       \ll_{q_0,d}(\log(2V_0))^{-2}.
\]
The finite supremum uses the two integrality conditions in the
preceding proposition and the actual twist at \((x,\ell,\iota)\).
Contributions of a zero range are taken as zero.
\end{theorem}
\begin{proof}
Compute the period matrix in one set of class representatives.
All the class weights and coefficient Gram entries are rational.
Applying \(\sigma\) to an entry therefore pairs the coordinates
of \(Py_i\) transported by \(\sigma^\#\) with those of \(Pz_k\)
transported by \(\sigma\), in the same ambient weighted
coordinate space: \(\sigma(\bar u)=\overline{\sigma^\#(u)}\).
The determinant of these mixed pairings is the Hermitian pairing
on their \(h\)-th exterior powers. Apply Cauchy--Schwarz,
followed by Proposition~\ref{prop:archcontraction804} at the
two indicated automorphisms. Hadamard's inequality, or ordinary
Gram--Schmidt on each list, proves the second inequality.

For \(h>D\) the period matrix has rank at most \(D\), so its
determinant is zero. For \(1\le h\le D\), the previous
proposition bounds the quotient in the second assertion by
\[
       \frac{\lambda\log\ell}{DsB_x^{\rm det}}+
       \frac{v_\ell(c_0)\log\ell}{sB_x^{\rm det}} .
\]
The sum of interval suprema of the first term is bounded by
Corollary~\ref{cor:actualdenom803}. For the second, \(s+t\le st\)
at the nodes under consideration, so block embedding in
\(\mathrm{GL}_{st}\) shows that \(c_0\) divides \(G_0\) of
Lemma~\ref{lem:sourcebudget803}. Replace the supremum over
integer nodes by their sum and use
\(\sum_\ell v_\ell(c_0)\log\ell=\log c_0\) at each node.
The calculation at the end of Corollary~\ref{cor:actualdenom803}
bounds the resulting sum by \(O_{q_0,d}(V_0^{-374})\).
The empty minor is one and costs zero. This proves the assertion
for every size simultaneously.
\end{proof}
The complex inequality retains the exterior norms of the actual
test arrays. A lower bound for a chain volume requires a nonzero
minor for the specified period maps and a quantitative estimate
for its size. The inequalities above allow the minor to vanish.
\clearpage

\clearpage
\section{References}\label{nav:bibliography-X}
\begingroup\def\section*#1{}\begin{thebibliography}{99}\small
\bibitem{X-VII-VII-chttransfer721} L. Clozel, M. Harris and R. Taylor, \emph{Automorphy for some \(\ell\)-adic lifts of automorphic mod \(\ell\) Galois representations}, Publ. Math. IH\'ES \textbf{108} (2008), 1--181, \S3.1. \href{https://virtualmath1.stanford.edu/~rltaylor/twugfin.pdf}{Authors' text}.
\bibitem{X-VII-zelevinsky738} A. V. Zelevinsky, \emph{Induced representations of reductive \(p\)-adic groups. II. On irreducible representations of \(\GL_n\)}, Ann. Sci. \'Ecole Norm. Sup. (4) \textbf{13} (1980), 165--210. \href{https://www.numdam.org/articles/10.24033/asens.1379/}{Original paper}.
\bibitem{X-VII-conradunits744} B. Conrad, \emph{Lifting global representations with local properties}, Appendix~B on locally algebraic characters. \href{https://math.stanford.edu/~conrad/papers/locchar.pdf}{Author's text}.
\bibitem{X-VII-localduality745} J. S. Milne, \emph{Arithmetic Duality Theorems}, second edition, BookSurge, 2006, Chapter~I, \S2. \href{https://www.jmilne.org/math/Books/ADTnot.pdf}{Author's text}.
\bibitem{X-VII-lrs787} W. Luo, Z. Rudnick and P. Sarnak, \emph{On the generalized Ramanujan conjecture for $\mathrm{GL}(n)$}, in \emph{Automorphic Forms, Automorphic Representations, and Arithmetic}, Proc. Sympos. Pure Math. \textbf{66}, part 2 (1999), 301--310. \href{https://doi.org/10.1090/pspum/066.2/1703764}{Published article}.
\end{thebibliography}\endgroup
\end{document}
